% =====================================================================
%  SJTU Mathematical Analysis
%  Chapter 2 : Limits of Sequences and the Fundamental Theorems
%              of Real Numbers
%
%  Every problem of the chapter -- the worked examples of the body of
%  the text and the four exercise sets -- transcribed and translated
%  into English, followed by a separate part containing all solutions.
%
%  Numbering convention
%    "Example 2.1.4"  = worked example 2.1.4 of the body of the text
%    "Exercise 2.1.4" = problem 4 of exercise set 2.1
%
%  Build with:  pdflatex <this file>     (run twice for the contents list)
% =====================================================================
\documentclass[11pt,a4paper]{article}

\usepackage[margin=2.4cm]{geometry}
\usepackage{amsmath,amssymb,amsthm,mathtools}
\usepackage[shortlabels]{enumitem}
\usepackage{xcolor}
\usepackage{microtype}
\usepackage[hidelinks]{hyperref}

\hypersetup{
  pdftitle={SJTU Mathematical Analysis, Chapter 2: Limits of Sequences},
}

\allowdisplaybreaks
\setlist[enumerate]{itemsep=0.35ex,topsep=0.5ex,parsep=0.2ex}
\setlist[itemize]{itemsep=0.35ex,topsep=0.5ex,parsep=0.2ex}

% ---- markup for the three kinds of block ----------------------------
\definecolor{problemcolor}{RGB}{20,60,130}
\definecolor{solutioncolor}{RGB}{20,20,20}
\definecolor{sourcecolor}{RGB}{105,105,105}

\newenvironment{problem}[1]{%
  \par\addvspace{1.7ex}%
  \noindent{\color{problemcolor}\textbf{#1}}\hspace{0.55em}\ignorespaces}%
  {\par\addvspace{0.4ex}}

\newenvironment{solution}[1]{%
  \par\addvspace{1.7ex}%
  \noindent{\color{solutioncolor}\textbf{#1}}\hspace{0.55em}\ignorespaces}%
  {\par\addvspace{0.9ex}}

% small grey line naming the pages of the textbook each group comes from
\newcommand{\srcline}[1]{%
  \par\addvspace{1.6ex}%
  \noindent{\color{sourcecolor}\small\itshape #1}%
  \par\addvspace{0.2ex}}

\newcommand{\R}{\mathbb{R}}
\newcommand{\N}{\mathbb{N}}
\newcommand{\Z}{\mathbb{Z}}
\newcommand{\eps}{\varepsilon}
\DeclareMathOperator{\sgn}{sgn}

\title{\bfseries Limits of Sequences\\[0.5em]
       \large\mdseries Problems and Solutions from SJTU Mathematical Analysis Textbook}
\date{}

\begin{document}
\maketitle

\tableofcontents
\pagebreak
\part*{Part I\quad Problems}

\addcontentsline{toc}{part}{Part I --- Problems}


\section*{A.\quad Worked examples from the text}

\addcontentsline{toc}{section}{A. Worked examples from the text}


\srcline{Worked examples 2.1.2--2.1.5}

\begin{problem}{Example 2.1.2}\label{e:2.1.2}
Using the $\varepsilon$-$N$ language of limits of sequences, verify that
\[
\lim_{n\to\infty} q^{n}=0 \qquad (0<\lvert q\rvert<1).
\]
\end{problem}

\begin{problem}{Example 2.1.3}\label{e:2.1.3}
Using the $\varepsilon$-$N$ language of limits of sequences, prove that
\[
\lim_{n\to\infty}\frac{n^{2}-n+2}{3n^{2}+2n-4}=\frac{1}{3}.
\]
\end{problem}

\begin{problem}{Example 2.1.4}\label{e:2.1.4}
Let $\{x_{n}\}$ be a sequence such that $\displaystyle\lim_{n\to\infty}x_{n}=a$. Prove that
\[
\lim_{n\to\infty}\frac{x_{1}+x_{2}+\cdots+x_{n}}{n}=a .
\]
\end{problem}

\begin{problem}{Example 2.1.5}\label{e:2.1.5}
Prove that
\[
\lim_{n\to\infty}\frac{1}{n^{2}}\neq 1 .
\]
\end{problem}


\srcline{Worked examples 2.1.12--2.1.15}

\begin{problem}{Example 2.1.12}\label{e:2.1.12}
Suppose that $\displaystyle\lim_{n\to\infty}x_{n}=a$. Using the definition of the limit of a sequence, prove that
\[
\lim_{n\to\infty}x_{n}^{\frac{1}{3}}=a^{\frac{1}{3}} .
\]
\end{problem}

\begin{problem}{Example 2.1.13}\label{e:2.1.13}
Prove that
\[
\lim_{n\to\infty}\frac{10^{n}}{n!}=0 .
\]
\end{problem}

\begin{problem}{Example 2.1.14}\label{e:2.1.14}
Let
\[
x_{n}=\frac{(2n-1)!!}{(2n)!!},
\qquad\text{where}\qquad
(2n-1)!!=1\cdot 3\cdot 5\cdots(2n-1),\quad (2n)!!=2\cdot 4\cdot 6\cdots(2n).
\]
Prove that the sequence $\{x_{n}\}$ converges to zero.
\end{problem}

\begin{problem}{Example 2.1.15}\label{e:2.1.15}
Prove that
\[
\lim_{n\to\infty}n^{\frac{1}{n}}=1 .
\]
\end{problem}


\srcline{Worked examples 2.1.17--2.1.21}

\begin{problem}{Example 2.1.17}\label{e:2.1.17}
Find the limit
\[
\lim_{n\to\infty}\sqrt{n}\,\bigl(\sqrt{n+1}-\sqrt{n}\bigr).
\]
\end{problem}

\begin{problem}{Example 2.1.18}\label{e:2.1.18}
Prove that if $a>0$, then
\[
\lim_{n\to\infty}\sqrt[n]{a}=1 .
\]
\end{problem}

\begin{problem}{Example 2.1.19}\label{e:2.1.19}
Prove that the sequence
\[
\left\{\frac{a^{n}}{n!}\right\}
\]
is an infinitesimal, where $a>0$.
\end{problem}

\begin{problem}{Example 2.1.20}\label{e:2.1.20}
Suppose that
\[
\lim_{n\to\infty}\bigl(x_{n}-x_{n-2}\bigr)=0 .
\]
Prove that
\[
\left\{\frac{x_{n}-x_{n-1}}{n}\right\}
\]
is an infinitesimal.
\end{problem}

\begin{problem}{Example 2.1.21}\label{e:2.1.21}
Suppose that
\[
\lim_{n\to\infty}x_{n}=a,\qquad \lim_{n\to\infty}y_{n}=b .
\]
Prove that
\[
\lim_{n\to\infty}\frac{x_{n}y_{1}+x_{n-1}y_{2}+\cdots+x_{1}y_{n}}{n}=ab .
\]
\end{problem}


\srcline{Worked examples 2.1.25--2.1.26}

\begin{problem}{Example 2.1.25}\label{e:2.1.25}
Let $x_n=\sin\dfrac{n\pi}{4}$. Discuss the convergence of the sequence $\{x_n\}$.
\end{problem}

\begin{problem}{Example 2.1.26}\label{e:2.1.26}
Prove that an unbounded sequence must have a subsequence which is infinitely large.
\end{problem}


\srcline{Worked examples 2.2.2--2.2.3}

\begin{problem}{Example 2.2.2}\label{e:2.2.2}
Suppose that $\lim\limits_{n\to\infty} a_n = a$ (with $a$ a real number). Prove that
\[
\lim_{n\to\infty}\frac{a_1 + 2^2 a_2 + \cdots + n^2 a_n}{n^3} = \frac{a}{3}.
\]
\end{problem}

\begin{problem}{Example 2.2.3}\label{e:2.2.3}
Compute the limit
\[
\lim_{n\to\infty} n\left(\frac{\sum\limits_{k=1}^{n}(2k+1)^2}{n^3} - \frac{4}{3}\right).
\]
\end{problem}


\srcline{Worked examples 2.3.2 and 2.3.6--2.3.9}

\begin{problem}{Example 2.3.2}\label{e:2.3.2}
Let $x_1$ be a given positive number. Prove that the sequence $\{x_n\}$ defined by
\[
x_{n+1}=\frac{3(1+x_n)}{3+x_n},\qquad n=1,2,\dots,
\]
converges, and find its limit.
\end{problem}

\begin{problem}{Example 2.3.6}\label{e:2.3.6}
Let
\[
x_n=1+\frac{1}{2^2}+\frac{1}{3^2}+\cdots+\frac{1}{n^2}.
\]
Prove that $\{x_n\}$ is a fundamental (Cauchy) sequence.
\end{problem}

\begin{problem}{Example 2.3.8}\label{e:2.3.8}
Let $x_0$ and $a$ be arbitrary constants. Set
\[
x_1=q\sin x_0+a,\quad x_2=q\sin x_1+a,\quad\cdots,\quad x_{n+1}=q\sin x_n+a,\quad\cdots,
\]
where $0<q<1$ is a constant. Prove that the sequence $\{x_n\}$ converges.
\end{problem}

\begin{problem}{Example 2.3.9}\label{e:2.3.9}
Let
\[
a_n=\sum_{i=1}^{n}\frac{1}{i}.
\]
Prove that the sequence $\{a_n\}$ diverges.
\end{problem}


\section*{B.\quad The exercise sets}

\addcontentsline{toc}{section}{B. The exercise sets}


\srcline{Exercise set 2.1, problems 1--17}

\begin{problem}{Exercise 2.1.1}\label{x:2.1.1}
Decide whether each of the following propositions is true or false. If it is
true, prove it; otherwise give a counterexample.
\begin{enumerate}[(1)]
  \item $\displaystyle\lim_{n\to\infty}x_n=a\ (a\in\mathbb{R})\iff
    \forall\,0<\varepsilon<1,\ \exists\,N\in\mathbb{N},\ \forall\,n\ge N:\ |x_n-a|\le\varepsilon$;
  \item $\displaystyle\lim_{n\to\infty}x_n=a\ (a\in\mathbb{R})\iff
    \forall\,0<\varepsilon<1,\ \exists\,N\in\mathbb{N},\ \forall\,n>N:\ |x_n-a|<\varepsilon$;
  \item $\displaystyle\lim_{n\to\infty}x_n=a\ (a\in\mathbb{R})\iff
    \forall\,\varepsilon>0,\ \exists\,N\in\mathbb{N},\ \forall\,n>N:\ |x_n-a|<M\varepsilon$
    \ ($M$ is a positive constant);
  \item $\displaystyle\lim_{n\to\infty}x_n=a\ (a\in\mathbb{R})\iff
    \forall\,k\in\mathbb{N},\ \exists\,N\in\mathbb{N},\ \forall\,n>N:\ |x_n-a|<\frac1k$;
  \item $\displaystyle\lim_{n\to\infty}x_n=a\ (a\in\mathbb{R})\iff
    \forall\,\varepsilon>0$, the interval $(a-\varepsilon,a+\varepsilon)$ contains
    infinitely many terms of the sequence $\{x_n\}$.
\end{enumerate}
\end{problem}

\begin{problem}{Exercise 2.1.2}\label{x:2.1.2}
Decide whether each of the following propositions is true or false. If it is
true, prove it; otherwise give a counterexample.
\begin{enumerate}[(1)]
  \item If $a_n\le b_n\le c_n\ (n\in\mathbb{N})$ and
    $\displaystyle\lim_{n\to\infty}(c_n-a_n)=0$, then the sequence $\{a_n\}$ converges;
  \item If $a_n\le b_n\le c_n\ (n\in\mathbb{N})$, the sequence $\{b_n\}$ converges and
    $\displaystyle\lim_{n\to\infty}(c_n-a_n)=0$, then the sequence $\{a_n\}$ converges;
  \item If $\{a_n\}$ converges and $a_n\ne 0$, $n\in\mathbb{N}$, then
    $\displaystyle\lim_{n\to\infty}\frac{a_{n+1}}{a_n}=1$;
  \item If $\displaystyle\lim_{n\to\infty}a_nb_n=0$, then
    $\displaystyle\lim_{n\to\infty}a_n=0$ or $\displaystyle\lim_{n\to\infty}b_n=0$;
  \item If $\displaystyle\lim_{n\to\infty}\frac{b_n}{a_n}=1$ and the sequence
    $\{a_n\}$ converges, then the sequence $\{b_n\}$ converges and
    $\displaystyle\lim_{n\to\infty}b_n=\lim_{n\to\infty}a_n$;
  \item If the sequences $\{a_{2k}\}$ and $\{a_{2k-1}\}$ both converge, then the
    sequence $\{a_n\}$ converges.
\end{enumerate}
\end{problem}

\begin{problem}{Exercise 2.1.3}\label{x:2.1.3}
Using the $\varepsilon$-$N$ definition, prove that:
\begin{enumerate}[(1)]
  \item $\displaystyle\lim_{n\to\infty}\frac{1}{1+\sqrt{n}}=0$;
  \item $\displaystyle\lim_{n\to\infty}\frac{3n^2+n}{2n^2-1}=\frac32$;
  \item $\displaystyle\lim_{n\to\infty}\bigl(\sqrt{n^2+n}-n\bigr)=\frac12$;
  \item $\displaystyle\lim_{n\to\infty}\arctan n=\frac{\pi}{2}$;
  \item $\displaystyle\lim_{n\to\infty}\frac{n^2\arctan n}{1+n^2}=\frac{\pi}{2}$.
\end{enumerate}
\end{problem}

\begin{problem}{Exercise 2.1.4}\label{x:2.1.4}
Find the following limits:
\begin{enumerate}[(1)]
  \item $\displaystyle\lim_{n\to\infty}\frac{(-2)^n+3^n}{(-2)^{n+1}+3^{n+1}}$;
  \item $\displaystyle\lim_{n\to\infty}
    \frac{\dfrac12+\dfrac{1}{2^2}+\cdots+\dfrac{1}{2^n}}
         {\dfrac13+\dfrac{1}{3^2}+\cdots+\dfrac{1}{3^n}}$;
  \item $\displaystyle\lim_{n\to\infty}
    \left(\frac{1}{n^2}+\frac{3}{n^2}+\cdots+\frac{2n-1}{n^2}\right)$;
  \item $\displaystyle\lim_{n\to\infty}
    \left(1+\frac12\right)\left(1+\frac{1}{2^2}\right)\left(1+\frac{1}{2^4}\right)
    \cdots\left(1+\frac{1}{2^{2^n}}\right)$;
  \item $\displaystyle\lim_{n\to\infty}
    \left(\frac{1}{\sqrt{n^2+1}}+\frac{1}{\sqrt{n^2+2}}+\cdots+\frac{1}{\sqrt{n^2+n}}\right)$;
  \item $\displaystyle\lim_{n\to\infty}\left(1-\frac{1}{\sqrt[n]{2}}\right)\cos n$;
  \item $\displaystyle\lim_{n\to\infty}\sqrt[n]{n^2-n+2}$;
  \item $\displaystyle\lim_{n\to\infty}\sqrt[n]{n\ln n}$;
  \item $\displaystyle\lim_{n\to\infty}\sqrt[n]{\arctan n}$;
  \item $\displaystyle\lim_{n\to\infty}
    \left(\sqrt[n]{1}+\sqrt[n]{2}+\cdots+\sqrt[n]{10}\right)$;
  \item $\displaystyle\lim_{n\to\infty}\sqrt[n]{2\sin^2 n+\cos^2 n}$;
  \item $\displaystyle\lim_{n\to\infty}\bigl((n+1)^{\alpha}-n^{\alpha}\bigr)$,
    $\alpha\in(0,1)$;
  \item $\displaystyle\lim_{n\to\infty}\frac{\lfloor na_n\rfloor}{n}$.
    Here $\displaystyle\lim_{n\to\infty}a_n=a$.
\end{enumerate}
\end{problem}

\begin{problem}{Exercise 2.1.5}\label{x:2.1.5}
Prove that the sequence
\[
  \left\{\frac{2n+(-1)^n n}{3n+1}\right\},\qquad n\in\mathbb{N},
\]
diverges.
\end{problem}

\begin{problem}{Exercise 2.1.6}\label{x:2.1.6}
Suppose that $a_n\ne 0$,
\[
\frac{a_{n+1}}{a_n}>0\qquad(n\in\mathbb{N}),
\]
and that
\[
\lim_{n\to\infty}\frac{a_{n+1}}{a_n}=0 .
\]
Prove that for all sufficiently large $n$ the sequence $\{a_n\}$ is monotone.
\end{problem}

\begin{problem}{Exercise 2.1.7}\label{x:2.1.7}
Suppose that the sequence $\{a_n\}$ converges. Prove that the sequence $\{a_n\}$ has either a largest term or a smallest term. Also ask: must it necessarily have both?
\end{problem}

\begin{problem}{Exercise 2.1.8}\label{x:2.1.8}
Given that $\lim\limits_{n\to\infty}x_n=a$ with $a\ne 0$, use the $\varepsilon$-$N$ language to prove:
\begin{enumerate}[(1)]
  \item $\lim\limits_{n\to\infty}\sqrt{x_n}=\sqrt{a}$;
  \item $\lim\limits_{n\to\infty}\sin\dfrac{1}{x_n^2}=\sin\dfrac{1}{a^2}$.
\end{enumerate}
\end{problem}

\begin{problem}{Exercise 2.1.9}\label{x:2.1.9}
Suppose that $\lim\limits_{n\to\infty}(x_n-x_{n-1})=d$. Prove that
\[
\lim_{n\to\infty}\frac{x_n}{n}=d .
\]
\end{problem}

\begin{problem}{Exercise 2.1.10}\label{x:2.1.10}
Suppose that $\lim\limits_{n\to\infty}a_n=a$ ($a_n>0$, $n\in\mathbb{N}$). Prove that
\[
\lim_{n\to\infty}\sqrt[n]{a_1a_2\cdots a_n}=a ,
\]
and then prove successively:
\begin{enumerate}[(1)]
  \item if $\lim\limits_{n\to\infty}\dfrac{a_{n+1}}{a_n}=a$ ($a_n>0$, $n\in\mathbb{N}$), then $\lim\limits_{n\to\infty}\sqrt[n]{a_n}=a$;
  \item $\lim\limits_{n\to\infty}\dfrac{\sqrt[n]{n!}}{n}=\dfrac{1}{\mathrm{e}}$.
\end{enumerate}
\end{problem}

\begin{problem}{Exercise 2.1.11}\label{x:2.1.11}
Prove: if the sequences $\{a_{3k-2}\}$, $\{a_{3k-1}\}$ and $\{a_{3k}\}$ all converge and have the same limit, then the sequence $\{a_n\}$ converges.
\end{problem}

\begin{problem}{Exercise 2.1.12}\label{x:2.1.12}
Suppose that $\lim\limits_{n\to\infty}x_n=+\infty$. Prove that
\[
\lim_{n\to\infty}\frac{x_1+x_2+\cdots+x_n}{n}=+\infty .
\]
\end{problem}

\begin{problem}{Exercise 2.1.13}\label{x:2.1.13}
Suppose that $a_n>0$ for all $n\in\mathbb{N}$ and that
\[
\lim_{n\to\infty}\frac{a_n}{a_{n+1}+a_{n+2}}=0 .
\]
Prove that the sequence $\{a_n\}$ is unbounded.
\end{problem}

\begin{problem}{Exercise 2.1.14}\label{x:2.1.14}
Prove that a necessary and sufficient condition for a sequence to be an infinitely large quantity is that every subsequence of the sequence is an unbounded sequence.
\end{problem}

\begin{problem}{Exercise 2.1.15}\label{x:2.1.15}
Suppose that $a_n\ne 0$ for all $n\in\mathbb{N}$, that $\lim\limits_{n\to\infty}a_n=0$, and that $\lim\limits_{n\to\infty}\dfrac{a_{n+1}}{a_n}=l$. Prove that $|l|\le 1$.
\end{problem}

\begin{problem}{Exercise 2.1.16}\label{x:2.1.16}
Suppose that $f(x)$ is strictly monotonically increasing on $[a,b]$, and that $\{x_n\}\subset[a,b]$ satisfies
\[
\lim_{n\to\infty}f(x_n)=f(a).
\]
Prove that $\lim\limits_{n\to\infty}x_n=a$.
\end{problem}

\begin{problem}{Exercise 2.1.17}\label{x:2.1.17}
Suppose that the sequence $\{x_n\}$ is monotonically increasing and that
\[
\lim_{n\to\infty}\frac{x_1+x_2+\cdots+x_n}{n}=a .
\]
Prove that $\lim\limits_{n\to\infty}x_n=a$.
\end{problem}


\srcline{Exercise set 2.2, problems 1--6 --- printed page 38}

\begin{problem}{Exercise 2.2.1}\label{x:2.2.1}
Give an example showing that the converse of the Stolz theorem is false.
\end{problem}

\begin{problem}{Exercise 2.2.2}\label{x:2.2.2}
Give an example showing that the Stolz theorem fails when $a=\infty$.
\end{problem}

\begin{problem}{Exercise 2.2.3}\label{x:2.2.3}
Compute the following limits:
\begin{enumerate}[(1)]
  \item $\displaystyle\lim_{n\to\infty}\frac{1+\frac{1}{\sqrt{2}}+\cdots+\frac{1}{\sqrt{n}}}{\ln\sqrt{n}}$;
  \item $\displaystyle\lim_{n\to\infty}\frac{1+\sqrt{2}+\cdots+\sqrt{n}}{n\sqrt{n}}$;
  \item $\displaystyle\lim_{n\to\infty}\frac{1+\sqrt{2}+\sqrt[3]{3}+\cdots+\sqrt[n]{n}}{n}$;
  \item $\displaystyle\lim_{n\to\infty}\frac{1^{2}+3^{2}+\cdots+(2n-1)^{2}}{n^{3}}$;
  \item $\displaystyle\lim_{n\to\infty}\frac{a_{1}+2a_{2}+\cdots+na_{n}}{1+2+\cdots+n}\qquad\Bigl(\text{given that }\lim_{n\to\infty}a_{n}=a\Bigr)$.
\end{enumerate}
\end{problem}

\begin{problem}{Exercise 2.2.4}\label{x:2.2.4}
Compute the limit $\displaystyle\lim_{n\to\infty}(n!)^{\frac{1}{n^{2}}}$.
\end{problem}

\begin{problem}{Exercise 2.2.5}\label{x:2.2.5}
Let $\displaystyle x_{n}=\frac{1}{n^{2}}\sum_{k=0}^{n}\ln C_{n}^{k}$, $n=1,2,\cdots$. Find the limit $\displaystyle\lim_{n\to\infty}x_{n}$.
\end{problem}

\begin{problem}{Exercise 2.2.6}\label{x:2.2.6}
Suppose that $\displaystyle\lim_{n\to\infty}n(A_{n}-A_{n-1})=0$. Prove that if the limit
$\displaystyle\lim_{n\to\infty}\frac{A_{1}+A_{2}+\cdots+A_{n}}{n}$ exists, then
\[
\lim_{n\to\infty}A_{n}=\lim_{n\to\infty}\frac{A_{1}+A_{2}+\cdots+A_{n}}{n}.
\]
\end{problem}


\srcline{Exercise set 2.3, problems 1--2}

\begin{problem}{Exercise 2.3.1}\label{x:2.3.1}
Prove that each of the following sequences has a limit and find its value:
\begin{enumerate}[(1)]
  \item $x_1=\sqrt{2}$, $\displaystyle x_{n+1}=\sqrt{2x_n}$, $n=1,2,\dots$;
  \item $x_1=\sqrt{c}$ $(c>0)$, $\displaystyle x_{n+1}=\sqrt{c+x_n}$, $n=1,2,\dots$;
  \item $x_1>1$, $\displaystyle x_{n+1}=\frac{3x_n+1}{x_n+3}$, $n=1,2,\dots$.
\end{enumerate}
\end{problem}

\begin{problem}{Exercise 2.3.2}\label{x:2.3.2}
Let $0<a_1<b_1$, and set
\[
a_{n+1}=\sqrt{a_nb_n},\qquad b_{n+1}=\frac{a_n+b_n}{2},\qquad n\in\mathbb N_+.
\]
Prove that $\{a_n\}$ and $\{b_n\}$ converge to the same limit.
\end{problem}

\begin{problem}{Exercise 2.3.3}\label{x:2.3.3}
Let the sequence $\{x_n\}$ satisfy $0<x_n<1$ and
\[
(1-x_n)x_{n+1}>\frac14,\qquad n=1,2,3,\dots .
\]
Prove that
\[
\lim_{n\to\infty}x_n=\frac12 .
\]
\end{problem}

\begin{problem}{Exercise 2.3.4}\label{x:2.3.4}
Prove: if a monotone sequence $\{x_n\}$ contains a convergent subsequence, then $\{x_n\}$ converges.
\end{problem}

\begin{problem}{Exercise 2.3.5}\label{x:2.3.5}
Let $x_1\in(0,1)$ and $x_{n+1}=x_n(1-x_n)$, $n\in\mathbb N_+$. Prove that the sequence $\{nx_n\}$ converges and find its limit.
\end{problem}

\begin{problem}{Exercise 2.3.6}\label{x:2.3.6}
Prove: if $a_n>0$ and $\displaystyle\lim_{n\to\infty}\frac{a_n}{a_{n+1}}=l>1$, then $\displaystyle\lim_{n\to\infty}a_n=0$.
\end{problem}

\begin{problem}{Exercise 2.3.7}\label{x:2.3.7}
Prove that
\[
\frac12+\frac13+\dots+\frac1n<\ln(n+1)<1+\frac12+\dots+\frac1n .
\]
\end{problem}

\begin{problem}{Exercise 2.3.8}\label{x:2.3.8}
Prove that the sequence $\{x_n\}$ converges, where
\[
x_n=1+\frac12+\frac13+\dots+\frac1n-\ln n,\qquad n\in\mathbb N_+ .
\]
\end{problem}

\begin{problem}{Exercise 2.3.9}\label{x:2.3.9}
Find the limit
\[
\lim_{n\to\infty}\bigl(n!\,e-[n!\,e]\bigr),
\]
where $[\,\cdot\,]$ denotes the greatest-integer (floor) function.
\end{problem}

\begin{problem}{Exercise 2.3.10}\label{x:2.3.10}
Find the limit
\[
\lim_{n\to\infty}\left(\frac{1}{n+1}+\frac{1}{n+2}+\dots+\frac{1}{n+n}\right).
\]
\end{problem}

\begin{problem}{Exercise 2.3.11}\label{x:2.3.11}
Let
\[
x_n=\left(1+\frac12\right)\left(1+\frac{1}{2^2}\right)\dots\left(1+\frac{1}{2^n}\right).
\]
Prove that $\displaystyle\lim_{n\to\infty}x_n$ exists.
\end{problem}

\begin{problem}{Exercise 2.3.12}\label{x:2.3.12}
Can the following statements serve as necessary and sufficient conditions for the convergence of a sequence?
\begin{enumerate}[(1)]
  \item $\forall\,\varepsilon>0$, $\forall\,p\in\mathbb N$, $\exists\,N\in\mathbb N$, $\forall\,n>N$: $|x_{n+p}-x_n|<\varepsilon$;
  \item $\forall\,\varepsilon>0$, $\exists\,N,p\in\mathbb N$, $\forall\,n>N$: $|x_{n+p}-x_n|<\varepsilon$;
  \item $\forall\,n,p\in\mathbb N_+$, $\displaystyle|x_{n+p}-x_n|<\frac{p}{n}$; \ or \ $\forall\,p\in\mathbb N_+$, $\displaystyle\lim_{n\to\infty}(x_{n+p}-x_n)=0$;
  \item $\forall\,n,p\in\mathbb N_+$, $\displaystyle|x_{n+p}-x_n|<\frac{p}{n^2}$.
\end{enumerate}
\end{problem}

\begin{problem}{Exercise 2.3.13}\label{x:2.3.13}
Apply the Cauchy convergence criterion to prove that each of the following sequences $\{a_n\}$ converges:
\begin{enumerate}[(1)]
  \item $\displaystyle a_n = 1 - \frac{1}{2^2} + \frac{1}{3^2} - \cdots + (-1)^{n-1}\frac{1}{n^2}\qquad (n \in \mathbb{N}_+)$;
  \item $\displaystyle a_n = 1 + \frac{\sin x}{1^2} + \cdots + \frac{\sin nx}{n^2}\qquad (x \in \mathbb{R})$;
  \item $\displaystyle a_n = \frac{\sin 2x}{2(2+\sin 2x)} + \frac{\sin 3x}{3(3+\sin 3x)} + \cdots + \frac{\sin nx}{n(n+\sin nx)}\qquad (n \in \mathbb{N}_+,\ x \in \mathbb{R})$.
\end{enumerate}
\end{problem}

\begin{problem}{Exercise 2.3.14}\label{x:2.3.14}
Suppose the sequence
\[
\bigl\{\,|a_2 - a_1| + |a_3 - a_2| + \cdots + |a_n - a_{n-1}|\,\bigr\}
\]
is bounded. Prove that $\{a_n\}$ converges.
\end{problem}

\begin{problem}{Exercise 2.3.15}\label{x:2.3.15}
Prove the $0/0$ form of the Stolz theorem: suppose that
\[
\lim_{n\to\infty} a_n = 0, \qquad \lim_{n\to\infty} b_n = 0,
\qquad b_1 > b_2 > b_3 > \cdots .
\]
If
\[
\lim_{n\to\infty} \frac{a_n - a_{n-1}}{b_n - b_{n-1}} = A,
\]
then
\[
\lim_{n\to\infty} \frac{a_n}{b_n} = A .
\]
Here $A$ is either a finite number, or positive or negative infinity.
\end{problem}

\begin{problem}{Exercise 2.3.16}\label{x:2.3.16}
Let $a_n \in [a,b]$ $(n \in \mathbb{N}_+)$. Prove: if $\{a_n\}$ diverges, then $\{a_n\}$ must have two subsequences converging to different numbers.
\end{problem}

\begin{problem}{Exercise 2.3.17}\label{x:2.3.17}
Use the Cauchy convergence criterion to prove that a monotone bounded sequence must converge.
\end{problem}

\begin{problem}{Exercise 2.3.18}\label{x:2.3.18}
Use the nested interval theorem to prove the Bolzano--Weierstrass theorem.
\end{problem}

\begin{problem}{Exercise 2.3.19}\label{x:2.3.19}
Use the nested interval theorem to prove the Cauchy convergence criterion.
\end{problem}

\begin{problem}{Exercise 2.3.20}\label{x:2.3.20}
Let the sequence $\{x_n\}$ satisfy
\[
\lim_{n\to\infty} x_n \sum_{k=1}^{n} x_k^2 = 1 .
\]
Prove that
\[
\lim_{n\to\infty} \sqrt[3]{3n}\, x_n = 1 .
\]
\end{problem}

\begin{problem}{Exercise 2.3.21}\label{x:2.3.21}
Let $a_0 = 3$, $a_n = a_{n-1}^2 - 2$. Prove:
\begin{enumerate}[(1)]
  \item $\displaystyle \lim_{n\to\infty} a_n = +\infty$;
  \item $\displaystyle \lim_{n\to\infty} \frac{a_n}{a_0 a_1 \cdots a_{n-1}} = \sqrt{5}$.
\end{enumerate}
\end{problem}


\srcline{Exercise set 2.4, problems 1--10}

\begin{problem}{Exercise 2.4.1}\label{x:2.4.1}
Prove the uniqueness of the supremum and of the infimum.
\end{problem}

\begin{problem}{Exercise 2.4.2}\label{x:2.4.2}
Let $A$ be a nonempty set of real numbers. Prove:
\begin{enumerate}[(1)]
  \item $\sup A \in A$ if and only if $A$ has a maximum, and in that case $\sup A = \max A$;
  \item $\inf A \in A$ if and only if $A$ has a minimum, and in that case $\inf A = \min A$.
\end{enumerate}
\end{problem}

\begin{problem}{Exercise 2.4.3}\label{x:2.4.3}
Suppose that $x < a$ holds for every $x \in S$. Which of the two conclusions $\sup S < a$ and $\sup S \le a$ is the correct one?
\end{problem}

\begin{problem}{Exercise 2.4.4}\label{x:2.4.4}
Find the supremum and the infimum of each of the following sets $S$:
\begin{enumerate}[(1)]
  \item $S = \{x \in \mathbb{Q} \mid x > \sqrt{2}\}$;
  \item $S = \{y \in \mathbb{R} \mid y = x^2,\ x \in (-\sqrt{3}, 10)\}$;
  \item $S = \{y \in \mathbb{R} \mid y = \arctan x,\ x \in \mathbb{Q}\}$;
  \item $S = \left\{1, \dfrac{1}{2}, \dfrac{2}{3}, \dfrac{3}{4}, \cdots \right\}$.
\end{enumerate}
\end{problem}

\begin{problem}{Exercise 2.4.5}\label{x:2.4.5}
Let $A$ be a nonempty set of real numbers and write $-A = \{x \mid -x \in A\}$. Then
\begin{enumerate}[(1)]
  \item $\sup(-A) = -\inf A$;
  \item $\inf(-A) = -\sup A$.
\end{enumerate}
\end{problem}

\begin{problem}{Exercise 2.4.6}\label{x:2.4.6}
Let the sets of real numbers $A$ and $B$ be bounded above. If a set $S$ satisfies
\[
S \subset \{x + y \mid x \in A,\ y \in B\},
\]
prove that $\sup S \le \sup A + \sup B$. Give an example showing that the strict inequality need not hold.
\end{problem}

\begin{problem}{Exercise 2.4.7}\label{x:2.4.7}
Let the sets of real numbers $A$ and $B$ be bounded above. If a set $S$ satisfies
\[
S \supset \{x + y \mid x \in A,\ y \in B\},
\]
prove that $\sup S \ge \sup A + \sup B$. Give an example showing that the strict inequality need not hold.
\end{problem}

\begin{problem}{Exercise 2.4.8}\label{x:2.4.8}
Let $A, B$ be two nonempty sets of real numbers and let $\alpha \ge 0$. Write
\[
A + \alpha B = \{x \mid x = a + \alpha b,\ a \in A,\ b \in B\}.
\]
Prove:
\begin{enumerate}[(1)]
  \item $\sup(A + \alpha B) = \sup A + \alpha \sup B$;
  \item $\inf(A + \alpha B) = \inf A + \alpha \inf B$.
\end{enumerate}
\end{problem}

\begin{problem}{Exercise 2.4.9}\label{x:2.4.9}
Let $A$ and $B$ be two sets of real numbers with $A \subset B$. Prove that $\sup A \le \sup B$ and $\inf A \ge \inf B$.
\end{problem}

\begin{problem}{Exercise 2.4.10}\label{x:2.4.10}
Let $A$ and $B$ be two sets of real numbers such that the inequality $x \le y$ holds for every number $x$ in the set $A$ and every number $y$ in the set $B$. Prove that $\sup A \le \inf B$.
\end{problem}

\clearpage

\part*{Part II\quad Solutions}

\addcontentsline{toc}{part}{Part II --- Solutions}


\section*{A.\quad Worked examples from the text}

\addcontentsline{toc}{section}{A. Worked examples from the text}


\srcline{Worked examples 2.1.2--2.1.5}

\begin{solution}{Solution 2.1.2}\label{es:2.1.2}
Recall the definition: $\displaystyle\lim_{n\to\infty}a_{n}=a$ means that for every $\varepsilon>0$ there is a natural number $N$ such that $\lvert a_{n}-a\rvert<\varepsilon$ for every $n>N$.

Let $0<\lvert q\rvert<1$ and let $\varepsilon>0$ be arbitrary. Since $\lvert q^{n}-0\rvert=\lvert q\rvert^{n}$ and $\lvert q\rvert^{n}>0$, the desired inequality
\[
\lvert q^{n}-0\rvert=\lvert q\rvert^{n}<\varepsilon
\]
is, after taking logarithms of both sides, equivalent to
\[
n\ln\lvert q\rvert<\ln\varepsilon .
\]
Because $0<\lvert q\rvert<1$ we have $\ln\lvert q\rvert<0$, so dividing by $\ln\lvert q\rvert$ reverses the inequality and we obtain the sufficient condition
\[
n>\frac{\ln\varepsilon}{\ln\lvert q\rvert}.
\]
Let
\[
N=\max\left\{\left[\frac{\ln\varepsilon}{\ln\lvert q\rvert}\right],\,1\right\},
\]
where $[\,x\,]$ denotes the integer part of $x$, i.e.\ the greatest integer not exceeding $x$; then $N$ is a natural number with $N\ge\left[\dfrac{\ln\varepsilon}{\ln\lvert q\rvert}\right]$. Hence for every $n>N$ the index $n$ is an integer with $n\ge N+1>\left[\dfrac{\ln\varepsilon}{\ln\lvert q\rvert}\right]$, and since the integer part is the greatest integer not exceeding the number, this forces $n>\dfrac{\ln\varepsilon}{\ln\lvert q\rvert}$, that is, $n\ln\lvert q\rvert<\ln\varepsilon$, and therefore
\[
\lvert q^{n}-0\rvert=\lvert q\rvert^{n}<\varepsilon .
\]
Thus for every $\varepsilon>0$ there exists $N=\max\left\{\left[\dfrac{\ln\varepsilon}{\ln\lvert q\rvert}\right],1\right\}\in\mathbb{N}$ such that $\lvert q^{n}-0\rvert<\varepsilon$ for all $n>N$, which is precisely the statement
\[
\lim_{n\to\infty}q^{n}=0 .
\]
\end{solution}

\begin{solution}{Solution 2.1.3}\label{es:2.1.3}
We use again the definition: $\displaystyle\lim_{n\to\infty}a_{n}=a$ means that for every $\varepsilon>0$ there is a natural number $N$ with $\lvert a_{n}-a\rvert<\varepsilon$ for all $n>N$.

First simplify the difference. For every $n$ with $3n^{2}+2n-4\neq 0$,
\[
\left\lvert\frac{n^{2}-n+2}{3n^{2}+2n-4}-\frac{1}{3}\right\rvert
=\left\lvert\frac{3(n^{2}-n+2)-(3n^{2}+2n-4)}{3(3n^{2}+2n-4)}\right\rvert
=\left\lvert\frac{-5n+10}{3(3n^{2}+2n-4)}\right\rvert .
\]
For $n\ge 2$ we have $5n-10\ge 0$ and $3n^{2}+2n-4\ge 3n^{2}>0$, so the absolute values may be removed and
\[
\left\lvert\frac{n^{2}-n+2}{3n^{2}+2n-4}-\frac{1}{3}\right\rvert
=\frac{5n-10}{3(3n^{2}+2n-4)}
\le\frac{5n}{3\cdot 3n^{2}}
=\frac{5}{9n}
<\frac{1}{n},
\]
where in the first inequality we used $5n-10\le 5n$ and $3n^{2}+2n-4\ge 3n^{2}$ (valid for $n\ge 2$).

Now let $\varepsilon>0$ be arbitrary. In view of the estimate above, the inequality
\[
\left\lvert\frac{n^{2}-n+2}{3n^{2}+2n-4}-\frac{1}{3}\right\rvert<\varepsilon
\]
is certainly satisfied as soon as $\dfrac{1}{n}<\varepsilon$, that is, as soon as
\[
n>\frac{1}{\varepsilon}.
\]
Take
\[
N=\max\left\{\left[\frac{1}{\varepsilon}\right],\,2\right\},
\]
where $[\,x\,]$ denotes the integer part. Then $N\in\mathbb{N}$, $N\ge \left[\dfrac{1}{\varepsilon}\right]$ and $N\ge 2$. Hence for every $n>N$ we have on the one hand $n\ge 3>2$, so that the estimate above applies, and on the other hand, $n$ being an integer with $n\ge N+1>\left[\dfrac{1}{\varepsilon}\right]$, we have $n>\dfrac{1}{\varepsilon}$, so that $\dfrac{1}{n}<\varepsilon$. Consequently
\[
\left\lvert\frac{n^{2}-n+2}{3n^{2}+2n-4}-\frac{1}{3}\right\rvert<\frac{1}{n}<\varepsilon
\qquad\text{for every } n>N .
\]
Therefore for every $\varepsilon>0$ there exists $N=\max\left\{\left[\dfrac{1}{\varepsilon}\right],2\right\}\in\mathbb{N}$ such that the last inequality holds whenever $n>N$, and so
\[
\lim_{n\to\infty}\frac{n^{2}-n+2}{3n^{2}+2n-4}=\frac{1}{3}.
\]
(Remark: the enlargement of the inequality requires the condition $n\ge 2$, so when choosing $N$ we must ensure simultaneously that $N\ge\left[\dfrac{1}{\varepsilon}\right]$ and $N\ge 2$; this is exactly what the maximum above does.)
\end{solution}

\begin{solution}{Solution 2.1.4}\label{es:2.1.4}
Again, $\displaystyle\lim_{n\to\infty}a_{n}=a$ means that for every $\varepsilon>0$ there is a natural number $N$ such that $\lvert a_{n}-a\rvert<\varepsilon$ for every $n>N$.

Set $y_{k}=x_{k}-a$ for $k=1,2,\dots$. Then $\displaystyle\lim_{n\to\infty}y_{n}=0$, and
\[
\frac{x_{1}+x_{2}+\cdots+x_{n}}{n}-a
=\frac{(x_{1}-a)+(x_{2}-a)+\cdots+(x_{n}-a)}{n}
=\frac{y_{1}+y_{2}+\cdots+y_{n}}{n}.
\]
Hence it suffices to prove the assertion in the case $a=0$, i.e.\ to show that $\displaystyle\lim_{n\to\infty}y_{n}=0$ implies $\displaystyle\lim_{n\to\infty}\frac{y_{1}+\cdots+y_{n}}{n}=0$; the general case follows by adding $a$ back.

So assume $\displaystyle\lim_{n\to\infty}x_{n}=0$. Let $\varepsilon>0$ be given. By the definition of limit applied with the positive number $\varepsilon/2$, there is $N_{1}\in\mathbb{N}$ such that
\[
\lvert x_{n}\rvert<\frac{\varepsilon}{2}\qquad\text{for every } n>N_{1}.
\]
Fix this $N_{1}$ and put $C=\lvert x_{1}\rvert+\lvert x_{2}\rvert+\cdots+\lvert x_{N_{1}}\rvert$; this is a fixed real number, independent of $n$. For $n>N_{1}$, using the triangle inequality,
\[
\left\lvert\frac{x_{1}+x_{2}+\cdots+x_{n}}{n}\right\rvert
\le\frac{\lvert x_{1}+x_{2}+\cdots+x_{N_{1}}\rvert}{n}
+\frac{\lvert x_{N_{1}+1}+x_{N_{1}+2}+\cdots+x_{n}\rvert}{n}
\]
\[
\le\frac{\lvert x_{1}\rvert+\lvert x_{2}\rvert+\cdots+\lvert x_{N_{1}}\rvert}{n}
+\frac{\lvert x_{N_{1}+1}\rvert+\lvert x_{N_{1}+2}\rvert+\cdots+\lvert x_{n}\rvert}{n}
<\frac{C}{n}+\frac{n-N_{1}}{n}\cdot\frac{\varepsilon}{2}
<\frac{C}{n}+\frac{\varepsilon}{2},
\]
because the sum $\lvert x_{N_{1}+1}\rvert+\cdots+\lvert x_{n}\rvert$ contains exactly $n-N_{1}<n$ terms, each of them less than $\varepsilon/2$.

Since $C$ is a constant, $\displaystyle\lim_{n\to\infty}\frac{C}{n}=0$; equivalently, there exists $N_{2}\in\mathbb{N}$ such that
\[
\frac{C}{n}=\left\lvert\frac{C}{n}-0\right\rvert<\frac{\varepsilon}{2}\qquad\text{for every } n>N_{2}.
\]
Now take $N=\max\{N_{1},N_{2}\}$. For every $n>N$ we have $n>N_{1}$ and $n>N_{2}$, hence the two estimates above hold simultaneously and
\[
\left\lvert\frac{x_{1}+x_{2}+\cdots+x_{n}}{n}\right\rvert<\frac{\varepsilon}{2}+\frac{\varepsilon}{2}=\varepsilon .
\]
Thus, in the case $a=0$, for every $\varepsilon>0$ there is $N=\max\{N_{1},N_{2}\}\in\mathbb{N}$ with $\left\lvert\frac{x_{1}+\cdots+x_{n}}{n}-0\right\rvert<\varepsilon$ for all $n>N$, that is,
\[
\lim_{n\to\infty}\frac{x_{1}+x_{2}+\cdots+x_{n}}{n}=0 .
\]
Returning to the general case by means of the reduction at the beginning, we obtain
\[
\lim_{n\to\infty}\frac{x_{1}+x_{2}+\cdots+x_{n}}{n}=a .
\]
\end{solution}

\begin{solution}{Solution 2.1.5}\label{es:2.1.5}
We must correctly negate the definition of limit. By definition, $\displaystyle\lim_{n\to\infty}a_{n}=a$ means
\[
\forall\,\varepsilon>0,\ \exists\,N\in\mathbb{N},\ \forall\,n>N:\quad \lvert a_{n}-a\rvert<\varepsilon .
\]
Negating this statement (each ``for all'' becomes ``there exists'', each ``there exists'' becomes ``for all'', and the final inequality is negated), we obtain: $\displaystyle\lim_{n\to\infty}a_{n}\neq a$ means
\[
\exists\,\varepsilon_{0}>0,\ \forall\,N\in\mathbb{N},\ \exists\,n>N:\quad \lvert a_{n}-a\rvert\ge\varepsilon_{0}.
\]
Take $a_{n}=\dfrac{1}{n^{2}}$ and $a=1$, and choose
\[
\varepsilon_{0}=\frac{3}{4}.
\]
Let $N\in\mathbb{N}$ be arbitrary and take
\[
n=\max\{N+1,\,2\},
\]
so that $n>N$ and $n\ge 2$. Since $n\ge 2$ we have $n^{2}\ge 4$, hence $n^{2}-1\ge n^{2}-\dfrac{n^{2}}{4}=\dfrac{3}{4}n^{2}>0$, and therefore
\[
\left\lvert\frac{1}{n^{2}}-1\right\rvert=\frac{n^{2}-1}{n^{2}}\ge\frac{\frac{3}{4}n^{2}}{n^{2}}=\frac{3}{4}=\varepsilon_{0}.
\]
So there exists $\varepsilon_{0}=\dfrac{3}{4}>0$ such that for every natural number $N$ there is an index $n=\max\{N+1,2\}>N$ with $\left\lvert\dfrac{1}{n^{2}}-1\right\rvert\ge\varepsilon_{0}$. By the negation of the definition of limit this says exactly that
\[
\lim_{n\to\infty}\frac{1}{n^{2}}\neq 1 .
\]
\end{solution}


\srcline{Worked examples 2.1.12--2.1.15}

\begin{solution}{Solution 2.1.12}\label{es:2.1.12}
Recall the definition: $\displaystyle\lim_{n\to\infty}x_{n}=a$ means that for every $\varepsilon>0$ there is a natural number $N$ such that
\[
\lvert x_{n}-a\rvert<\varepsilon\qquad\text{for every }n>N .
\]
We prove the assertion by distinguishing the two cases $a=0$ and $a\neq 0$.

\medskip
\noindent\textbf{Case 1: $a=0$.}
Let $\varepsilon>0$ be arbitrary. Since $\varepsilon^{3}>0$ and $\displaystyle\lim_{n\to\infty}x_{n}=0$, the definition supplies a natural number $N$ such that
\[
\lvert x_{n}-0\rvert=\lvert x_{n}\rvert<\varepsilon^{3}\qquad\text{for every }n>N .
\]
Because the cube root is an odd function, $\lvert x_{n}^{1/3}\rvert=\lvert x_{n}\rvert^{1/3}$; hence for every $n>N$
\[
\bigl\lvert x_{n}^{\frac{1}{3}}-0^{\frac{1}{3}}\bigr\rvert=\lvert x_{n}^{\frac{1}{3}}\rvert=\lvert x_{n}\rvert^{\frac{1}{3}}<\bigl(\varepsilon^{3}\bigr)^{\frac{1}{3}}=\varepsilon .
\]
As $\varepsilon>0$ was arbitrary, this proves
\[
\lim_{n\to\infty}x_{n}^{\frac{1}{3}}=0=0^{\frac{1}{3}} .
\]

\medskip
\noindent\textbf{Case 2: $a\neq 0$.}
Here $a^{2/3}=\lvert a\rvert^{2/3}>0$, so $\varepsilon' :=a^{2/3}\varepsilon$ is a positive number whenever $\varepsilon>0$.

We first record a consequence of the sign-preserving property of limits: since $\displaystyle\lim_{n\to\infty}x_{n}=a$ with $a\neq 0$, applying the definition with the positive number $\lvert a\rvert$ gives a natural number $N_{2}$ such that $\lvert x_{n}-a\rvert<\lvert a\rvert$ for all $n>N_{2}$. If $a>0$ this gives $x_{n}>a-\lvert a\rvert=0$, while if $a<0$ it gives $x_{n}<a+\lvert a\rvert=0$; in either case
\[
x_{n}a>0\qquad\text{for every }n>N_{2}.
\]

Now let $\varepsilon>0$ be arbitrary. Applying the definition of $\displaystyle\lim_{n\to\infty}x_{n}=a$ to the positive number $\varepsilon'=a^{2/3}\varepsilon$, we obtain a natural number $N_{1}$ such that
\[
\lvert x_{n}-a\rvert<a^{\frac{2}{3}}\varepsilon\qquad\text{for every }n>N_{1}.
\]
Put $N=\max\{N_{1},N_{2}\}$ and fix any $n>N$. Then both estimates above hold. We use the algebraic identity
\[
u^{3}-v^{3}=(u-v)\bigl(u^{2}+uv+v^{2}\bigr),
\]
valid for all real $u,v$; taking $u=x_{n}^{1/3}$ and $v=a^{1/3}$ (so that $u^{3}=x_{n}$ and $v^{3}=a$) it yields
\[
x_{n}-a=\bigl(x_{n}^{\frac{1}{3}}-a^{\frac{1}{3}}\bigr)\bigl(x_{n}^{\frac{2}{3}}+x_{n}^{\frac{1}{3}}a^{\frac{1}{3}}+a^{\frac{2}{3}}\bigr),
\]
that is,
\[
x_{n}^{\frac{1}{3}}-a^{\frac{1}{3}}
=\frac{x_{n}-a}{x_{n}^{\frac{2}{3}}+x_{n}^{\frac{1}{3}}a^{\frac{1}{3}}+a^{\frac{2}{3}}} .
\]
Since $x_{n}a>0$, each of the three terms in the denominator is strictly positive: $x_{n}^{2/3}>0$, $x_{n}^{1/3}a^{1/3}=(x_{n}a)^{1/3}>0$ and $a^{2/3}>0$. Consequently the denominator is strictly larger than $a^{2/3}$, so
\[
\frac{1}{x_{n}^{\frac{2}{3}}+x_{n}^{\frac{1}{3}}a^{\frac{1}{3}}+a^{\frac{2}{3}}}<a^{-\frac{2}{3}} .
\]
Taking absolute values in the identity above and combining the two estimates, we get
\[
\bigl\lvert x_{n}^{\frac{1}{3}}-a^{\frac{1}{3}}\bigr\rvert
=\frac{\lvert x_{n}-a\rvert}{x_{n}^{\frac{2}{3}}+x_{n}^{\frac{1}{3}}a^{\frac{1}{3}}+a^{\frac{2}{3}}}
\le\lvert x_{n}-a\rvert\,a^{-\frac{2}{3}}
<a^{\frac{2}{3}}\varepsilon\cdot a^{-\frac{2}{3}}=\varepsilon .
\]
Thus for every $\varepsilon>0$ there is a natural number $N=\max\{N_{1},N_{2}\}$ with $\lvert x_{n}^{1/3}-a^{1/3}\rvert<\varepsilon$ for all $n>N$, which is exactly the statement
\[
\lim_{n\to\infty}x_{n}^{\frac{1}{3}}=a^{\frac{1}{3}} .
\]
\end{solution}

\begin{solution}{Solution 2.1.13}\label{es:2.1.13}
Write $y_{n}=\dfrac{10^{n}}{n!}$; clearly $y_{n}>0$ for every $n$. We first bound $y_{n}$ from above.

Let $n>10$. Splitting off the first ten factors of $n!=1\cdot 2\cdots n$ and of $10^{n}$,
\[
y_{n}=\frac{10^{n}}{n!}
=\frac{10^{10}}{10!}\cdot\frac{10}{11}\cdot\frac{10}{12}\cdots\frac{10}{n}
=\frac{10^{10}}{10!}\prod_{k=11}^{n}\frac{10}{k} .
\]
For every $k\ge 11$ we have $0<\dfrac{10}{k}\le\dfrac{10}{11}<1$, and the product contains exactly $n-10$ factors; therefore
\[
0<y_{n}\le\frac{10^{10}}{10!}\left(\frac{10}{11}\right)^{n-10}
\qquad (n>10),
\]
where the inequality is an equality exactly when all the factors $\frac{10}{k}$ equal $\frac{10}{11}$, that is, exactly when $n=11$. For every $n>11$ it is strict. In any case the upper bound tends to $0$, as we now check.

Since $0<\frac{10}{11}<1$, the geometric sequence with ratio $q=\frac{10}{11}$ tends to $0$: indeed $q=\frac{1}{1+h}$ with $h=\frac{1}{10}>0$, and Bernoulli's inequality gives $(1+h)^{m}\ge 1+mh$ for every $m\ge 1$, whence
\[
0<q^{m}=\frac{1}{(1+h)^{m}}\le\frac{1}{1+mh}\longrightarrow 0 ,
\]
so that $\displaystyle\lim_{m\to\infty}q^{m}=0$ by the squeeze theorem. Taking $m=n-10$ and multiplying by the constant $\dfrac{10^{10}}{10!}$,
\[
\lim_{n\to\infty}\frac{10^{10}}{10!}\left(\frac{10}{11}\right)^{n-10}
=\frac{10^{10}}{10!}\cdot 0=0 .
\]

We may now apply the squeeze theorem for sequences (if $x_{n}\le y_{n}\le z_{n}$ for all sufficiently large $n$ and $\displaystyle\lim_{n\to\infty}x_{n}=\lim_{n\to\infty}z_{n}=a$, then $\displaystyle\lim_{n\to\infty}y_{n}=a$) to the three sequences
\[
x_{n}=0,\qquad y_{n}=\frac{10^{n}}{n!},\qquad
z_{n}=\frac{10^{10}}{10!}\left(\frac{10}{11}\right)^{n-10},
\]
which satisfy $0<y_{n}\le z_{n}$ for all $n>10$, while $\displaystyle\lim_{n\to\infty}x_{n}=0$ and $\displaystyle\lim_{n\to\infty}z_{n}=0$. We conclude that
\[
\lim_{n\to\infty}\frac{10^{n}}{n!}=0 .
\]
\end{solution}

\begin{solution}{Solution 2.1.14}\label{es:2.1.14}
For a natural number $m$ we have, by cross-multiplication with the positive denominators,
\[
\frac{m}{m+1}<\frac{m+1}{m+2}
\quad\Longleftrightarrow\quad
m(m+2)<(m+1)^{2}
\quad\Longleftrightarrow\quad
m^{2}+2m<m^{2}+2m+1 ,
\]
which is true. Since
\[
x_{n}=\frac{(2n-1)!!}{(2n)!!}
=\frac{1\cdot 3\cdot 5\cdots(2n-1)}{2\cdot 4\cdot 6\cdots(2n)}>0 ,
\]
the sequence $\{x_{n}\}$ is well defined and positive for every $n\ge 1$.

Apply the inequality above to the $n$ natural numbers $m=1,3,5,\dots,2n-1$:
\[
\frac{1}{2}<\frac{2}{3},\qquad
\frac{3}{4}<\frac{4}{5},\qquad
\frac{5}{6}<\frac{6}{7},\qquad\dots,\qquad
\frac{2n-1}{2n}<\frac{2n}{2n+1} .
\]
All the numbers involved are positive, so we may multiply these $n$ inequalities side by side; the left-hand sides multiply to $x_{n}$ and we obtain
\[
\frac{1\cdot 3\cdot 5\cdots(2n-1)}{2\cdot 4\cdot 6\cdots(2n)}
<\frac{2\cdot 4\cdot 6\cdots(2n)}{3\cdot 5\cdot 7\cdots(2n+1)} .
\]
On the right-hand side we factor $3\cdot 5\cdot 7\cdots(2n+1)=(2n+1)\bigl(3\cdot 5\cdots(2n-1)\bigr)$ and then cancel the common factor $3\cdot 5\cdots(2n-1)$:
\[
\frac{2\cdot 4\cdot 6\cdots(2n)}{3\cdot 5\cdot 7\cdots(2n+1)}
=\frac{1}{2n+1}\cdot\frac{2\cdot 4\cdot 6\cdots(2n)}{3\cdot 5\cdots(2n-1)}
=\frac{1}{(2n+1)\;x_{n}} ,
\]
because $\dfrac{2\cdot 4\cdot 6\cdots(2n)}{3\cdot 5\cdots(2n-1)}=\dfrac{2\cdot 4\cdots(2n)}{1\cdot 3\cdot 5\cdots(2n-1)}=\dfrac{1}{x_{n}}$. Therefore
\[
x_{n}<\frac{1}{(2n+1)x_{n}} .
\]
Multiplying by the positive number $x_{n}$ we find $x_{n}^{2}<\dfrac{1}{2n+1}$, and since $x_{n}>0$,
\[
0<x_{n}<\frac{1}{\sqrt{2n+1}}\qquad (n\ge 1).
\]

Finally, $\displaystyle\lim_{n\to\infty}\frac{1}{\sqrt{2n+1}}=0$. Indeed, let $\varepsilon>0$; by the Archimedean property choose a natural number $N$ with $2N+1\ge\dfrac{1}{\varepsilon^{2}}$. Then for every $n>N$ we have $2n+1>2N+1\ge\dfrac{1}{\varepsilon^{2}}$, hence $\sqrt{2n+1}>\dfrac{1}{\varepsilon}$ and therefore
\[
\left\lvert\frac{1}{\sqrt{2n+1}}-0\right\rvert=\frac{1}{\sqrt{2n+1}}<\varepsilon .
\]

Now the squeeze theorem for sequences applies to
\[
0<x_{n}<\frac{1}{\sqrt{2n+1}}\qquad (n\ge 1),
\]
where the left-hand sequence is identically $0$ and the right-hand sequence tends to $0$. Consequently
\[
\lim_{n\to\infty}x_{n}=\lim_{n\to\infty}\frac{(2n-1)!!}{(2n)!!}=0 ,
\]
so the sequence $\{x_{n}\}$ converges, and its limit is $0$.
\end{solution}

\begin{solution}{Solution 2.1.15}\label{es:2.1.15}
We use the inequality between the geometric and arithmetic means: for positive real numbers $a_{1},a_{2},\dots,a_{n}$,
\[
(a_{1}a_{2}\cdots a_{n})^{\frac{1}{n}}\le\frac{a_{1}+a_{2}+\cdots+a_{n}}{n},
\]
with equality if and only if $a_{1}=a_{2}=\cdots=a_{n}$.

Let $n\ge 2$ and apply this inequality to the $n$ positive numbers consisting of $n-2$ copies of $1$ and two copies of $\sqrt{n}$. Their product equals
\[
\underbrace{1\cdot 1\cdots 1}_{n-2\text{ factors}}\cdot\sqrt{n}\cdot\sqrt{n}=n,
\]
and their sum equals $n-2+2\sqrt{n}$. Hence
\[
n^{\frac{1}{n}}
=\bigl(1\cdot 1\cdots 1\cdot\sqrt{n}\cdot\sqrt{n}\bigr)^{\frac{1}{n}}
\le\frac{n-2+2\sqrt{n}}{n}
=1+\frac{2\sqrt{n}-2}{n}
\le 1+\frac{2\sqrt{n}}{n}
=1+\frac{2}{\sqrt{n}} .
\]
On the other hand $n\ge 1$ gives $n^{1/n}\ge 1^{1/n}=1$ (the map $t\mapsto t^{1/n}$ is increasing on $[0,+\infty)$ for every $n\ge 1$). Thus, for every $n\ge 1$,
\[
1\le n^{\frac{1}{n}}\le 1+\frac{2}{\sqrt{n}} .
\]
(For $n=1,2$ the estimate is immediate: $1^{1}=1\le 3$ and $2^{1/2}=\sqrt{2}\le 1+\sqrt{2}$.)

It remains to note that $\displaystyle\lim_{n\to\infty}\left(1+\frac{2}{\sqrt{n}}\right)=1$, which follows from $\displaystyle\lim_{n\to\infty}\frac{1}{\sqrt{n}}=0$: given $\varepsilon>0$, the Archimedean property provides a natural number $N$ with $N>\dfrac{1}{\varepsilon^{2}}$, and then for every $n>N$ we have $\sqrt{n}>\dfrac{1}{\varepsilon}$, that is, $\left\lvert\dfrac{1}{\sqrt{n}}-0\right\rvert=\dfrac{1}{\sqrt{n}}<\varepsilon$.

The constant sequence $1$ tends to $1$, and the sequence $n^{1/n}$ is squeezed between the constant sequence $1$ and the sequence $1+\dfrac{2}{\sqrt{n}}$, which also tends to $1$. By the squeeze theorem for sequences,
\[
\lim_{n\to\infty}n^{\frac{1}{n}}=1 .
\]
\end{solution}


\srcline{Worked examples 2.1.17--2.1.21}

\begin{solution}{Solution 2.1.17}\label{es:2.1.17}
Rationalize the factor $\sqrt{n+1}-\sqrt{n}$ by multiplying and dividing by
$\sqrt{n+1}+\sqrt{n}$:
\[
\sqrt{n}\,\bigl(\sqrt{n+1}-\sqrt{n}\bigr)
=\frac{\sqrt{n}\,\bigl(\sqrt{n+1}-\sqrt{n}\bigr)\bigl(\sqrt{n+1}+\sqrt{n}\bigr)}
       {\sqrt{n+1}+\sqrt{n}}
=\frac{\sqrt{n}\,\bigl((n+1)-n\bigr)}{\sqrt{n+1}+\sqrt{n}}
=\frac{\sqrt{n}}{\sqrt{n+1}+\sqrt{n}} .
\]
Dividing numerator and denominator by $\sqrt{n}$ (which is legitimate for
$n\geq 1$),
\[
\frac{\sqrt{n}}{\sqrt{n+1}+\sqrt{n}}
=\frac{1}{\sqrt{1+\frac1n}+1}.
\]
Now $\frac1n\to 0$, hence $1+\frac1n\to 1$. For the square root we have, for every
$n\geq 1$,
\[
0\leq \sqrt{1+\tfrac1n}-1
=\frac{\bigl(\sqrt{1+\frac1n}-1\bigr)\bigl(\sqrt{1+\frac1n}+1\bigr)}
        {\sqrt{1+\frac1n}+1}
=\frac{\frac1n}{\sqrt{1+\frac1n}+1}
\leq \frac1n ,
\]
and $\frac1n\to 0$; by the squeeze theorem (sandwich theorem)
$\sqrt{1+\frac1n}-1\to 0$, that is, $\sqrt{1+\frac1n}\to 1$. Since the denominator
$\sqrt{1+\frac1n}+1$ has limit $1+1=2\neq 0$, the quotient rule for limits gives
\[
\lim_{n\to\infty}\sqrt{n}\,\bigl(\sqrt{n+1}-\sqrt{n}\bigr)
=\lim_{n\to\infty}\frac{1}{\sqrt{1+\frac1n}+1}=\frac12 .
\]
\end{solution}

\begin{solution}{Solution 2.1.18}\label{es:2.1.18}
We treat the three cases $a=1$, $a>1$ and $0<a<1$.

\emph{Case 1: $a=1$.} Then $\sqrt[n]{a}=\sqrt[n]{1}=1$ for every $n$, so
$\lim_{n\to\infty}\sqrt[n]{a}=1$.

\emph{Case 2: $a>1$.} For every $n\geq 1$ we have $\sqrt[n]{a}>1$. Put
\[
h_{n}=\sqrt[n]{a}-1>0,\qquad\text{so that}\qquad \sqrt[n]{a}=1+h_{n}.
\]
Raising the last equality to the $n$-th power and applying the binomial theorem,
\[
a=(1+h_{n})^{n}
=1+\binom{n}{1}h_{n}+\binom{n}{2}h_{n}^{2}+\cdots+\binom{n}{n}h_{n}^{n}
\geq 1+n\,h_{n},
\]
because $h_{n}>0$ makes every term $\binom{n}{k}h_{n}^{k}$ with $k\geq 1$
nonnegative (in particular $\binom{n}{1}h_{n}=n h_{n}$). Hence
\[
0<h_{n}\leq \frac{a-1}{n}.
\]
Since $a-1$ is a fixed constant and $\lim_{n\to\infty}\frac{a-1}{n}=0$, the squeeze
theorem gives $\lim_{n\to\infty}h_{n}=0$. Therefore, by the sum rule for limits,
\[
\lim_{n\to\infty}\sqrt[n]{a}=\lim_{n\to\infty}(1+h_{n})=1+0=1 .
\]

\emph{Case 3: $0<a<1$.} Then $\frac1a>1$, so Case 2 applied to the number
$\frac1a>1$ yields
\[
\lim_{n\to\infty}\sqrt[n]{\frac1a}=1 .
\]
For every $n\geq 1$,
\[
\sqrt[n]{a}=\frac{1}{\sqrt[n]{\frac1a}},
\]
and the denominator has limit $1\neq 0$; by the quotient rule for limits,
\[
\lim_{n\to\infty}\sqrt[n]{a}
=\frac{1}{\lim_{n\to\infty}\sqrt[n]{\frac1a}}=\frac11=1 .
\]
\end{solution}

\begin{solution}{Solution 2.1.19}\label{es:2.1.19}
Recall that a sequence $\{z_{n}\}$ is called an \emph{infinitesimal} (a null
sequence) if $\lim_{n\to\infty}z_{n}=0$. Let $x_{n}=\frac{a^{n}}{n!}>0$. We prove
that $x_{n}\to 0$, distinguishing two cases according to the integer part
$\lfloor a\rfloor$ of $a$.

\emph{Case 1: $\lfloor a\rfloor=0$} (that is, $0<a<1$). For every $n\geq 1$ we have
$n!\geq n$ and $a^{n}<1$, whence
\[
0<x_{n}=\frac{a^{n}}{n!}\leq \frac{a^{n}}{n}<\frac1n .
\]
Since $\frac1n\to 0$, the squeeze theorem yields $\lim_{n\to\infty}x_{n}=0$, i.e.
$\left\{\frac{a^{n}}{n!}\right\}$ is an infinitesimal.

\emph{Case 2: $\lfloor a\rfloor\neq 0$} (that is, $a\geq 1$). Let
$m=\lfloor a\rfloor\geq 1$, so that $m\leq a<m+1$. Fix $n>m$; then
\[
\frac{a^{n}}{n!}
=\frac{a}{1}\cdot\frac{a}{2}\cdots\frac{a}{m}\cdot\frac{a}{m+1}\cdots\frac{a}{n}
\leq \frac{a^{m}}{m!}\cdot\frac{a}{n},
\]
because the first $m$ factors have product $\frac{a^{m}}{m!}$, while each of the
remaining factors satisfies $\frac{a}{k}\leq\frac{a}{n}$ for $m+1\leq k\leq n$
(the larger the denominator, the smaller the fraction, and $a>0$). Thus, whenever
$n>m$ and $\frac{a}{n}\leq 1$ (which holds for all $n\geq a$),
\[
0<x_{n}\leq \frac{a^{m}}{m!}\cdot\frac{a}{n},
\]
and $\frac{a^{m}}{m!}\cdot\frac{a}{n}\to 0$ as $n\to\infty$, since
$\frac{a^{m+1}}{m!}$ is a fixed constant. By the squeeze theorem,
\[
\lim_{n\to\infty}\frac{a^{n}}{n!}=0,
\]
so $\left\{\frac{a^{n}}{n!}\right\}$ is an infinitesimal.

Combining the two cases, $\left\{\frac{a^{n}}{n!}\right\}$ is an infinitesimal for
every $a>0$.
\end{solution}

\begin{solution}{Solution 2.1.20}\label{es:2.1.20}
By hypothesis $\lim_{n\to\infty}(x_{n}-x_{n-2})=0$: for every $\varepsilon>0$ there
exists a natural number $N_{1}$ such that
\[
| x_{n}-x_{n-2}|<\frac{\varepsilon}{2}\qquad\text{for all } n>N_{1}.
\]
We may assume $N_{1}\geq 2$. Fix $n>\max\{2,N_{1}\}$. Using the triangle
inequality $|x_{k}-x_{k-1}|\leq|x_{k}-x_{k-2}|+|x_{k-1}-x_{k-2}|$ and applying it
recursively, we obtain
\[
\begin{aligned}
|x_{n}-x_{n-1}|
&\leq |x_{n}-x_{n-2}|+|x_{n-1}-x_{n-2}|\\
&\leq |x_{n}-x_{n-2}|+|x_{n-1}-x_{n-3}|+|x_{n-2}-x_{n-3}|\\
&\leq \cdots\\
&\leq |x_{n}-x_{n-2}|+|x_{n-1}-x_{n-3}|+\cdots+|x_{N_{1}+1}-x_{N_{1}-1}|
   +|x_{N_{1}}-x_{N_{1}-1}|\\
&=\sum_{k=N_{1}+1}^{n}|x_{k}-x_{k-2}|+|x_{N_{1}}-x_{N_{1}-1}|\\
&<\frac{(n-N_{1})\varepsilon}{2}+|x_{N_{1}}-x_{N_{1}-1}|,
\end{aligned}
\]
where in the last step we used that all indices $k$ occurring in the sum satisfy
$k\geq N_{1}+1>N_{1}$, so each of the $n-N_{1}$ summands is less than
$\frac{\varepsilon}{2}$.

The integer $N_{1}$ is now fixed, so $|x_{N_{1}}-x_{N_{1}-1}|$ is a constant and
therefore
\[
\left\{\frac{1}{n}\,|x_{N_{1}}-x_{N_{1}-1}|\right\}
\]
is an infinitesimal: indeed its term is a fixed constant divided by $n$. Hence for
the given $\varepsilon>0$ there exists a natural number $N_{2}$ such that
\[
\frac{1}{n}\,|x_{N_{1}}-x_{N_{1}-1}|<\frac{\varepsilon}{2}\qquad\text{for all }
n>N_{2}.
\]
Put $N=\max\{2,N_{1},N_{2}\}$. Then for every $n>N$ the two estimates above hold
simultaneously, and
\[
\left|\frac{x_{n}-x_{n-1}}{n}\right|
=\frac{|x_{n}-x_{n-1}|}{n}
<\frac{(n-N_{1})\varepsilon}{2n}+\frac{|x_{N_{1}}-x_{N_{1}-1}|}{n}
<\frac{\varepsilon}{2}+\frac{\varepsilon}{2}=\varepsilon .
\]
Therefore
\[
\lim_{n\to\infty}\frac{x_{n}-x_{n-1}}{n}=0,
\]
that is, $\left\{\frac{x_{n}-x_{n-1}}{n}\right\}$ is an infinitesimal.
\end{solution}

\begin{solution}{Solution 2.1.21}\label{es:2.1.21}
Recall that $\{z_{n}\}$ is an \emph{infinitesimal} if $\lim_{n\to\infty}z_{n}=0$;
equivalently, $\lim_{n\to\infty}z_{n}=c$ if and only if $z_{n}=c+\gamma_{n}$ with
$\{\gamma_{n}\}$ an infinitesimal.

Set
\[
\alpha_{n}=x_{n}-a,\qquad \beta_{n}=y_{n}-b .
\]
Since $\lim_{n\to\infty}x_{n}=a$ and $\lim_{n\to\infty}y_{n}=b$, both
$\{\alpha_{n}\}$ and $\{\beta_{n}\}$ are infinitesimals, so that
\[
x_{n}=a+\alpha_{n},\qquad y_{n}=b+\beta_{n}.
\]
Every convergent sequence is bounded; hence there is a constant $M>0$ with
\[
|\alpha_{n}|\leq M,\qquad |\beta_{n}|\leq M\qquad\text{for all } n\geq 1 .
\]
Write
\[
I_{n}=\frac{x_{n}y_{1}+x_{n-1}y_{2}+\cdots+x_{1}y_{n}}{n}
=\frac{1}{n}\sum_{k=1}^{n}x_{n-k+1}y_{k}.
\]
Substituting $x_{n-k+1}=a+\alpha_{n-k+1}$ and $y_{k}=b+\beta_{k}$ and expanding,
\[
\begin{aligned}
I_{n}
&=\frac{1}{n}\sum_{k=1}^{n}\bigl(a+\alpha_{n-k+1}\bigr)\bigl(b+\beta_{k}\bigr)\\
&=ab+\frac{b}{n}\sum_{k=1}^{n}\alpha_{k}+\frac{a}{n}\sum_{k=1}^{n}\beta_{k}
   +\frac{1}{n}\sum_{k=1}^{n}\alpha_{n-k+1}\beta_{k},
\end{aligned}
\]
where in the second term we used the reindexing
$\sum_{k=1}^{n}\alpha_{n-k+1}=\sum_{j=1}^{n}\alpha_{j}$.

\emph{Step 1 (Ces\`aro mean lemma).} If $z_{n}\to 0$, then
$\frac{1}{n}\sum_{j=1}^{n}z_{j}\to 0$.

Indeed, let $\varepsilon>0$. Choose $N_{1}$ such that $|z_{j}|<\frac{\varepsilon}{2}$
for all $j>N_{1}$, and put $C=\sum_{j=1}^{N_{1}}|z_{j}|$ (a fixed constant; if
$N_{1}=0$ take $C=0$). For $n>N_{1}$,
\[
\left|\frac{1}{n}\sum_{j=1}^{n}z_{j}\right|
\leq \frac{C}{n}+\frac{1}{n}\sum_{j=N_{1}+1}^{n}|z_{j}|
<\frac{C}{n}+\frac{(n-N_{1})\varepsilon}{2n}
<\frac{C}{n}+\frac{\varepsilon}{2}.
\]
Choose $N_{2}\geq N_{1}$ with $\frac{C}{n}<\frac{\varepsilon}{2}$ for all $n>N_{2}$.
Then $\left|\frac{1}{n}\sum_{j=1}^{n}z_{j}\right|<\varepsilon$ for all $n>N_{2}$,
which proves the lemma. Applying it to the infinitesimals $\{\alpha_{n}\}$ and
$\{\beta_{n}\}$,
\[
\lim_{n\to\infty}\frac{b}{n}\sum_{k=1}^{n}\alpha_{k}=0,
\qquad
\lim_{n\to\infty}\frac{a}{n}\sum_{k=1}^{n}\beta_{k}=0 .
\]

\emph{Step 2.} $\displaystyle\lim_{n\to\infty}\frac{1}{n}\sum_{k=1}^{n}
\alpha_{n-k+1}\beta_{k}=0$.

Let $\varepsilon>0$. Since $\beta_{k}\to 0$, choose $N_{1}$ such that
$|\beta_{k}|<\frac{\varepsilon}{2M}$ for all $k>N_{1}$. Since $\alpha_{j}\to 0$ and
the indices $n-k+1$, for $k=1,\dots,N_{1}$, tend to $+\infty$ as $n\to\infty$,
there exists $N_{2}$ such that
\[
|\alpha_{n-k+1}|<\frac{\varepsilon}{2MN_{1}}\qquad
\text{for all } n>N_{2}\ \text{ and all } 1\leq k\leq N_{1}
\]
(if $N_{1}=0$ the block $1\leq k\leq N_{1}$ is empty and no such choice is needed).
Then, for $n>\max\{N_{1},N_{2}\}$ (so that the block $1\leq k\leq N_{1}$ lies inside
the range $1\leq k\leq n$),
\[
\begin{aligned}
\left|\frac{1}{n}\sum_{k=1}^{n}\alpha_{n-k+1}\beta_{k}\right|
&\leq \frac{1}{n}\sum_{k=1}^{N_{1}}|\alpha_{n-k+1}|\,|\beta_{k}|
   +\frac{1}{n}\sum_{k=N_{1}+1}^{n}|\alpha_{n-k+1}|\,|\beta_{k}|\\
&\leq \frac{1}{n}\cdot N_{1}\cdot\frac{\varepsilon}{2MN_{1}}\cdot M
   +\frac{1}{n}\cdot\frac{\varepsilon}{2M}\sum_{k=N_{1}+1}^{n}|\alpha_{n-k+1}|\\
&\leq \frac{\varepsilon}{2n}+\frac{\varepsilon}{2M}\cdot M
\;\leq\; \varepsilon ,
\end{aligned}
\]
where we used $\sum_{k=N_{1}+1}^{n}|\alpha_{n-k+1}|\leq\sum_{j=1}^{n}|\alpha_{j}|
\leq nM$. This proves Step 2.

Combining the two steps with the expansion of $I_{n}$,
\[
\lim_{n\to\infty}I_{n}
=ab+0+0+0=ab,
\]
that is,
\[
\lim_{n\to\infty}\frac{x_{n}y_{1}+x_{n-1}y_{2}+\cdots+x_{1}y_{n}}{n}=ab .
\]
\end{solution}


\srcline{Worked examples 2.1.25--2.1.26}

\begin{solution}{Solution 2.1.25}\label{es:2.1.25}
We use the following theorem: if a sequence $\{x_n\}$ converges to $a$, then every
subsequence of $\{x_n\}$ also converges to $a$. Indeed, let $\{x_{n_k}\}$ be any
subsequence of $\{x_n\}$; given $\varepsilon>0$, choose $N$ such that
$|x_n-a|<\varepsilon$ for all $n>N$. Since $n_k\ge k$, for every $k>N$ we have
$n_k\ge k>N$ and therefore $|x_{n_k}-a|<\varepsilon$. Hence
$\lim\limits_{k\to\infty}x_{n_k}=a$. Consequently, if two subsequences of one and the
same sequence have different limits, that sequence cannot converge.

Now consider the two subsequences of $\{x_n\}$ indexed by the even multiples of $4$
and by the numbers congruent to $2$ modulo $8$.

First, for $n=4k$, $k\in\mathbb{N}_+$:
\[
x_{4k}=\sin\frac{4k\pi}{4}=\sin k\pi=0
\qquad\text{for every } k\in\mathbb{N}_+,
\]
so the subsequence $\{x_{4k}\}$ is the constant sequence $0,0,0,\dots$, whence
\[
\lim_{k\to\infty}x_{4k}=0 .
\]

Second, for $n=8k+2$, $k\in\mathbb{N}_+$:
\[
x_{8k+2}=\sin\frac{(8k+2)\pi}{4}=\sin\left(2k+\frac12\right)\pi
=\sin\left(2k\pi+\frac{\pi}{2}\right)=1
\qquad\text{for every } k\in\mathbb{N}_+,
\]
so the subsequence $\{x_{8k+2}\}$ is the constant sequence $1,1,1,\dots$, whence
\[
\lim_{k\to\infty}x_{8k+2}=1 .
\]

The two subsequences $\{x_{4k}\}$ and $\{x_{8k+2}\}$ are subsequences of $\{x_n\}$ and
have the different limits $0$ and $1$. If $\{x_n\}$ converged, say to $a$, then by the
theorem recalled above both subsequences would converge to $a$, forcing $a=0$ and
$a=1$, a contradiction. Therefore the sequence
$\{x_n\}=\left\{\sin\dfrac{n\pi}{4}\right\}$ diverges.
\end{solution}

\begin{solution}{Solution 2.1.26}\label{es:2.1.26}
Recall the definitions involved. A sequence $\{x_n\}$ is \emph{bounded} if there is a
constant $M>0$ such that $|x_n|\le M$ for all $n\in\mathbb{N}_+$; it is
\emph{unbounded} if no such constant exists. Thus $\{x_n\}$ is unbounded if and only if
\begin{equation}\label{eq:unb}
\forall\, G>0,\ \exists\, n\in\mathbb{N}_+:\ |x_n|>G .
\end{equation}
Indeed, if for some $G>0$ no such $n$ existed, then $|x_n|\le G$ for all $n$ and the
sequence would be bounded; conversely, if the sequence is unbounded, then for every
$G>0$ the choice $M=G$ fails, so some $n$ satisfies $|x_n|>G$. A sequence
$\{y_k\}$ is \emph{infinitely large}, written $\lim\limits_{k\to\infty}y_k=\infty$, if
\begin{equation}\label{eq:inf}
\forall\, G>0,\ \exists\, K\in\mathbb{N}_+,\ \forall\, k>K:\ |y_k|>G .
\end{equation}

Now let $\{x_n\}$ be an unbounded sequence. We construct indices
$n_1<n_2<n_3<\cdots$ recursively.

\emph{Step 1.} Put $M_1=1$. By \eqref{eq:unb} with $G=M_1=1$ there exists
$n_1\in\mathbb{N}_+$ such that
\[
|x_{n_1}|>M_1=1 .
\]

\emph{Inductive step.} Suppose $k\in\mathbb{N}_+$ and $n_k\in\mathbb{N}_+$ have been
chosen with $|x_{n_k}|>k$. Put
\[
M_{k+1}=1+\sum_{i=1}^{n_k}|x_i| .
\]
By \eqref{eq:unb} with $G=M_{k+1}$ there exists $n\in\mathbb{N}_+$ such that
$|x_n|>M_{k+1}$. For every index $i\le n_k$ we have
\[
|x_i|\le\sum_{j=1}^{n_k}|x_j|<1+\sum_{j=1}^{n_k}|x_j|=M_{k+1}<|x_n| ,
\]
so $n\ne i$ for all $i\le n_k$, that is, $n>n_k$. Set $n_{k+1}=n$. Then
$n_{k+1}>n_k$ and
\[
|x_{n_{k+1}}|>M_{k+1}=1+\sum_{i=1}^{n_k}|x_i|\ge 1+|x_{n_k}|>1+k=k+1 .
\]
Thus the induction produces a strictly increasing sequence of indices
$n_1<n_2<n_3<\cdots$ in $\mathbb{N}_+$ such that
\begin{equation}\label{eq:key}
|x_{n_k}|>k\qquad\text{for every } k\in\mathbb{N}_+ .
\end{equation}
In particular $\{x_{n_k}\}$ is a subsequence of $\{x_n\}$.

It remains to check \eqref{eq:inf} for $\{x_{n_k}\}$. Let $G>0$ be arbitrary. By the
Archimedean property of the real numbers there is $K\in\mathbb{N}_+$ with $K>G$. Then
for every $k>K$, using \eqref{eq:key},
\[
|x_{n_k}|>k>K>G .
\]
Hence $\lim\limits_{k\to\infty}x_{n_k}=\infty$, i.e. the subsequence $\{x_{n_k}\}$ is
infinitely large.

Therefore every unbounded sequence has a subsequence which is infinitely large.
\end{solution}


\srcline{Worked examples 2.2.2--2.2.3}

\begin{solution}{Solution 2.2.2}\label{es:2.2.2}
We use the Stolz theorem (Theorem 2.2.1) in the following form: if $\{c_n\}$ is a strictly
increasing sequence with $c_n \to +\infty$, and if
\[
\lim_{n\to\infty}\frac{b_n - b_{n-1}}{c_n - c_{n-1}} = l
\]
for some real number $l$ (with $b_0$ arbitrary), then $\displaystyle\lim_{n\to\infty}\frac{b_n}{c_n} = l$.

Set
\[
c_n = n^3, \qquad b_n = \sum_{k=1}^{n} k^2 a_k \quad (n \ge 1), \qquad b_0 = 0 .
\]
The sequence $\{c_n\}$ is strictly increasing and $c_n \to +\infty$. For every $n \ge 1$ we have
\[
b_n - b_{n-1} = n^2 a_n,
\]
and
\[
c_n - c_{n-1} = n^3 - (n-1)^3 = 3n^2 - 3n + 1 .
\]
Hence
\[
\frac{b_n - b_{n-1}}{c_n - c_{n-1}}
= \frac{n^2 a_n}{3n^2 - 3n + 1}
= a_n \cdot \frac{n^2}{3n^2 - 3n + 1} .
\]
Now
\[
\frac{n^2}{3n^2 - 3n + 1} = \frac{1}{3 - \dfrac{3}{n} + \dfrac{1}{n^2}} \longrightarrow \frac13
\qquad (n \to \infty),
\]
and $a_n \to a$ by hypothesis. The limit of a product of two convergent sequences is the product
of their limits, so
\[
\lim_{n\to\infty}\frac{b_n - b_{n-1}}{c_n - c_{n-1}} = a \cdot \frac13 = \frac{a}{3}.
\]
The hypotheses of the Stolz theorem are therefore satisfied with $l = a/3$, and consequently
\[
\lim_{n\to\infty}\frac{a_1 + 2^2 a_2 + \cdots + n^2 a_n}{n^3}
= \lim_{n\to\infty}\frac{b_n}{c_n} = \frac{a}{3},
\]
which is the required identity.
\end{solution}

\begin{solution}{Solution 2.2.3}\label{es:2.2.3}
We again apply the Stolz theorem: if $\{c_n\}$ is strictly increasing with $c_n \to +\infty$ and
$\displaystyle\lim_{n\to\infty}\frac{b_n - b_{n-1}}{c_n - c_{n-1}} = l$ exists in $\mathbb{R}$, then
$\displaystyle\lim_{n\to\infty}\frac{b_n}{c_n} = l$.

First we simplify the expression inside the limit. Bringing the bracketed difference to a common
denominator,
\[
n\left(\frac{\sum\limits_{k=1}^{n}(2k+1)^2}{n^3} - \frac{4}{3}\right)
= \frac{\sum\limits_{k=1}^{n}(2k+1)^2 - \dfrac{4}{3}n^3}{n^2}.
\]
Accordingly, set
\[
c_n = n^2, \qquad b_n = \sum_{k=1}^{n}(2k+1)^2 - \frac{4}{3}n^3 \quad (n \ge 0),
\]
where the empty sum is $0$, so that $b_0 = 0$. The sequence $\{c_n\}$ is strictly increasing and
$c_n = n^2 \to +\infty$.

For $n \ge 1$,
\[
b_n - b_{n-1} = (2n+1)^2 - \frac{4}{3}\left[n^3 - (n-1)^3\right],
\qquad
c_n - c_{n-1} = n^2 - (n-1)^2 = 2n - 1 .
\]
Since
\[
n^3 - (n-1)^3 = 3n^2 - 3n + 1
\qquad\text{and}\qquad
(2n+1)^2 = 4n^2 + 4n + 1,
\]
we get
\[
b_n - b_{n-1}
= (4n^2 + 4n + 1) - \frac{4}{3}\left(3n^2 - 3n + 1\right)
= 4n^2 + 4n + 1 - 4n^2 + 4n - \frac{4}{3}
= 8n - \frac13 = \frac{24n - 1}{3}.
\]
Therefore
\[
\frac{b_n - b_{n-1}}{c_n - c_{n-1}}
= \frac{\dfrac{24n-1}{3}}{2n-1}
= \frac{24n-1}{3(2n-1)}
= \frac{24 - \dfrac{1}{n}}{3\left(2 - \dfrac{1}{n}\right)}
\longrightarrow \frac{24}{6} = 4
\qquad (n \to \infty).
\]
By the Stolz theorem,
\[
\lim_{n\to\infty} n\left(\frac{\sum\limits_{k=1}^{n}(2k+1)^2}{n^3} - \frac{4}{3}\right)
= \lim_{n\to\infty}\frac{b_n}{c_n} = 4 .
\]

For completeness, here is a direct check of the same value. Using
$\sum_{k=1}^{n} k = \frac{n(n+1)}{2}$ and $\sum_{k=1}^{n} k^2 = \frac{n(n+1)(2n+1)}{6}$,
\[
\sum_{k=1}^{n}(2k+1)^2 = 4\sum_{k=1}^{n}k^2 + 4\sum_{k=1}^{n}k + n
= \frac{4}{3}n^3 + 4n^2 + \frac{11}{3}n ,
\]
so that
\[
n\left(\frac{\sum\limits_{k=1}^{n}(2k+1)^2}{n^3} - \frac{4}{3}\right)
= n\left(\frac{4}{n} + \frac{11}{3n^2}\right)
= 4 + \frac{11}{3n} \longrightarrow 4 ,
\]
in agreement with the Stolz-theorem computation.
\end{solution}


\srcline{Worked examples 2.3.2 and 2.3.6--2.3.9}

\begin{solution}{Solution 2.3.2}\label{es:2.3.2}
We use the monotone bounded convergence theorem: every monotone sequence of real
numbers that is bounded is convergent.

\emph{The sequence is bounded.} Since $x_1>0$, and since $x_n>0$ implies
$x_{n+1}=3(1+x_n)/(3+x_n)>0$, induction gives $x_n>0$ for every $n\ge1$. Moreover
\[
x_{n+1}=\frac{3(1+x_n)}{3+x_n}=\frac{3(3+x_n)-6}{3+x_n}=3-\frac{6}{3+x_n}<3,
\]
because $3+x_n>0$. Hence $0<x_n<3$ for all $n\ge1$, so $\{x_n\}$ is bounded.

\emph{The sequence is monotone.} For $n\ge2$,
\[
x_{n+1}-x_n
=\frac{3(1+x_n)}{3+x_n}-\frac{3(1+x_{n-1})}{3+x_{n-1}}
=\frac{3(1+x_n)(3+x_{n-1})-3(1+x_{n-1})(3+x_n)}
{(3+x_n)(3+x_{n-1})}
=\frac{6\,(x_n-x_{n-1})}{(3+x_n)(3+x_{n-1})}.
\]
Since $3+x_n>0$ and $3+x_{n-1}>0$, the denominator is positive, so $x_{n+1}-x_n$
has the same sign as $x_n-x_{n-1}$. Iterating, $x_{n+1}-x_n$ has the same sign as
$x_2-x_1$ for every $n\ge1$. Consequently, if $x_2\ge x_1$ then $x_{n+1}\ge x_n$
for all $n$, and if $x_2\le x_1$ then $x_{n+1}\le x_n$ for all $n$; in either case
$\{x_n\}$ is monotone (nondecreasing, respectively nonincreasing).

\emph{Conclusion.} Being monotone and bounded, $\{x_n\}$ converges by the monotone
bounded convergence theorem. Write $a=\lim_{n\to\infty}x_n$.

\emph{Computation of the limit.} Since $x_n>0$ for all $n$, the limit satisfies
$a\ge0$. Passing to the limit in the recurrence and using the algebra of limits
(the denominator is nonzero in the limit, as $3+a\ge3>0$),
\[
a=\lim_{n\to\infty}x_{n+1}=\lim_{n\to\infty}\frac{3(1+x_n)}{3+x_n}
=\frac{3(1+a)}{3+a}.
\]
Thus $a(3+a)=3(1+a)$, that is, $a^2+3a=3+3a$, hence $a^2=3$ and $a=\pm\sqrt3$.
Since $a\ge0$, we must have $a=\sqrt3$. Therefore
\[
\lim_{n\to\infty}x_n=\sqrt3 .
\]
\end{solution}

\begin{solution}{Solution 2.3.6}\label{es:2.3.6}
Recall Definition 2.3.5: a sequence $\{x_n\}$ is a fundamental (Cauchy) sequence
if for every given $\varepsilon>0$ there is a natural number $N$ such that
$|x_m-x_n|<\varepsilon$ whenever $m,n>N$.

Let $m,n$ be positive integers with $m>n$; the case $n\ge m$ is identical because
$|x_m-x_n|=|x_n-x_m|$. Then
\[
x_m-x_n=\frac{1}{(n+1)^2}+\frac{1}{(n+2)^2}+\cdots+\frac{1}{m^2}>0 .
\]
For every integer $k\ge2$ we have $k^2=k(k-1)+k>k(k-1)$, hence
$\dfrac{1}{k^2}<\dfrac{1}{k(k-1)}=\dfrac{1}{k-1}-\dfrac{1}{k}$. Since the indices
$k=n+1,\dots,m$ all satisfy $k\ge2$, we obtain
\[
0<x_m-x_n<\sum_{k=n+1}^{m}\left(\frac{1}{k-1}-\frac{1}{k}\right)
=\left(\frac1n-\frac{1}{n+1}\right)+\cdots+\left(\frac{1}{m-1}-\frac1m\right)
=\frac1n-\frac1m<\frac1n .
\]
Now let $\varepsilon>0$ be given and choose
\[
N=\left[\frac{1}{\varepsilon}\right]+1,
\]
where $[\cdot]$ denotes the greatest-integer function. Then $N>\frac1\varepsilon$,
hence $\frac1N<\varepsilon$. If $m,n>N$, then, writing $r=\min\{m,n\}>N$ (and
noting that $|x_m-x_n|=0$ when $m=n$, while $|x_m-x_n|=x_{\max}-x_{\min}<\frac1r$
when $m\neq n$),
\[
|x_m-x_n|<\frac{1}{r}<\frac{1}{N}<\varepsilon .
\]
Thus the definition is satisfied, that is, $\{x_n\}$ is a fundamental (Cauchy)
sequence.
\end{solution}

\begin{solution}{Solution 2.3.8}\label{es:2.3.8}
We use the Cauchy convergence criterion (Theorem 2.3.7): a sequence $\{x_n\}$ of
real numbers converges if and only if it is a fundamental (Cauchy) sequence, that
is, if and only if for every $\varepsilon>0$ there is $N\in\mathbb{N}$ such that
\[
|x_{n+p}-x_n|<\varepsilon\qquad\text{for all }n>N\text{ and all }p\in\mathbb{N}.
\]
We also use the elementary estimates $|\sin t|\le|t|$ and $|\cos t|\le1$ for all
real $t$, together with
\[
\sin u-\sin v=2\sin\frac{u-v}{2}\cos\frac{u+v}{2}.
\]

Set $\lambda=|x_1-x_0|\ge0$, which is a constant depending only on $x_0$, $x_1$
and $q$.

\emph{Step 1: successive differences decay geometrically.} For every $n\ge1$,
since $x_{n+1}=q\sin x_n+a$ and $x_n=q\sin x_{n-1}+a$,
\[
|x_{n+1}-x_n|
=q\,|\sin x_n-\sin x_{n-1}|
=2q\left|\sin\frac{x_n-x_{n-1}}{2}\cos\frac{x_n+x_{n-1}}{2}\right|
\le 2q\cdot\frac{|x_n-x_{n-1}|}{2}
=q\,|x_n-x_{n-1}| .
\]
Applying this repeatedly, starting from $|x_2-x_1|=q|\sin x_1-\sin x_0|\le
q|x_1-x_0|=q\lambda$, we get by induction on $n$ that
\[
|x_{n+1}-x_n|\le q^n\lambda\qquad(n\ge1).
\]

\emph{Step 2: a Cauchy-type estimate.} Let $n\ge1$ and $p\in\mathbb{N}$. By the
triangle inequality and Step 1,
\[
|x_{n+p}-x_n|
\le\sum_{k=1}^{p}|x_{n+k}-x_{n+k-1}|
\le\sum_{k=1}^{p}q^{\,n+k-1}\lambda
=q^n\bigl(1+q+\cdots+q^{p-1}\bigr)\lambda .
\]
Since $0<q<1$, the finite geometric sum is bounded by the sum of the full
geometric series: $1+q+\cdots+q^{p-1}<\sum_{k=0}^{\infty}q^k=\dfrac{1}{1-q}$.
Hence
\[
|x_{n+p}-x_n|<\frac{q^n}{1-q}\,\lambda .
\]

\emph{Step 3: the criterion applies.} As $0<q<1$, the geometric sequence
$\{q^n\}$ tends to $0$, so
\[
\lim_{n\to\infty}\frac{q^n}{1-q}\,\lambda=0 .
\]
Therefore, given any $\varepsilon>0$, there exists $N\in\mathbb{N}$ such that
$\dfrac{q^n}{1-q}\lambda<\varepsilon$ for all $n>N$. For such $n$, Step 2 gives
\[
|x_{n+p}-x_n|<\frac{q^n}{1-q}\lambda<\varepsilon
\qquad\text{for every }p\in\mathbb{N}.
\]
Thus $\{x_n\}$ is a fundamental (Cauchy) sequence, and by the Cauchy convergence
criterion $\{x_n\}$ converges.
\end{solution}

\begin{solution}{Solution 2.3.9}\label{es:2.3.9}
Recall the negation of the Cauchy condition. A sequence $\{a_n\}$ is \emph{not} a
fundamental (Cauchy) sequence precisely when there exists an $\varepsilon_0>0$
such that for every $N\in\mathbb{N}$ one can find $n>N$ and $p\in\mathbb{N}$ with
\[
|a_{n+p}-a_n|\ge\varepsilon_0 .
\]
By the Cauchy convergence criterion (Theorem 2.3.7), a sequence of real numbers
converges if and only if it is a fundamental sequence; hence a sequence that is
not fundamental diverges. We shall show that $\{a_n\}$ fails the Cauchy condition
with the single choice $\varepsilon_0=\frac12$.

Let $N\ge2$ be a positive integer and take $n=N$ and $p=N$, so that $n+p=2N$.
Since the sum defining $a_n$ has positive terms and $a_{2N}>a_N$,
\[
|a_{n+p}-a_n|=|a_{2N}-a_N|=\sum_{i=N+1}^{2N}\frac1i .
\]
This sum contains exactly $N$ terms, and for each index $i$ with
$N+1\le i\le 2N$ we have $\frac1i\ge\frac{1}{2N}$. Therefore
\[
|a_{n+p}-a_n|=\sum_{i=N+1}^{2N}\frac1i\ge\sum_{i=N+1}^{2N}\frac{1}{2N}
=\frac{N}{2N}=\frac12=\varepsilon_0 .
\]

Suppose, for contradiction, that $\{a_n\}$ were a fundamental sequence. Then for
$\varepsilon_0=\frac12$ there would exist $N_0\in\mathbb{N}$ such that
$|a_{n+p}-a_n|<\frac12$ for all $n>N_0$ and all $p\in\mathbb{N}$. Choose a
positive integer $N>N_0$ with $N\ge2$ and put $n=p=N$; then $n=N>N_0$, yet the
computation above gives
$|a_{n+p}-a_n|\ge\frac12$, a contradiction. Hence $\{a_n\}$ is not a fundamental
sequence.

Since $\{a_n\}$ is not a fundamental sequence, the Cauchy convergence criterion
shows that $\{a_n\}$ does not converge; that is, the sequence
$a_n=\sum_{i=1}^{n}\frac1i$ diverges.
\end{solution}


\section*{B.\quad The exercise sets}

\addcontentsline{toc}{section}{B. The exercise sets}


\srcline{Exercise set 2.1, problems 1--17}

\begin{solution}{Solution 2.1.1}\label{xs:2.1.1}
Throughout we use the definition of convergence in the form
\[
  \lim_{n\to\infty}x_n=a\ (a\in\mathbb{R})
  \iff \forall\,\varepsilon>0,\ \exists\,N\in\mathbb{N},\ \forall\,n\ge N:\
  |x_n-a|<\varepsilon .
\]

\medskip
\noindent\textbf{(1) True.}
Let $(P_1)$ denote the statement on the right.

Assume first that $x_n\to a$ and let $0<\varepsilon<1$. By the definition there is
$N\in\mathbb{N}$ with $|x_n-a|<\varepsilon$ for all $n\ge N$; hence
$|x_n-a|\le\varepsilon$ for all $n\ge N$, so $(P_1)$ holds.

Conversely assume $(P_1)$ and let $\varepsilon>0$ be arbitrary. If
$\varepsilon<1$, then $(P_1)$ applied to this $\varepsilon$ gives
$N\in\mathbb{N}$ with $|x_n-a|\le\varepsilon$ for all $n\ge N$. If
$\varepsilon\ge1$, apply $(P_1)$ with $\varepsilon_1=\tfrac12<1$; there is
$N\in\mathbb{N}$ with
\[
  |x_n-a|\le\tfrac12<\varepsilon\qquad(n\ge N).
\]
In both cases find $N\in\mathbb{N}$ with $|x_n-a|<\varepsilon$ for all $n\ge N$;
hence $x_n\to a$. The equivalence is therefore true.

\medskip
\noindent\textbf{(2) True.}
Let $(P_2)$ denote the statement on the right; compared with the definition, it
restricts $\varepsilon$ to $(0,1)$ and requires $n>N$ instead of $n\ge N$.

If $x_n\to a$, the definition immediately gives $(P_2)$.

Conversely assume $(P_2)$ and let $\varepsilon>0$. Put
$\varepsilon_1=\min\{\varepsilon,\tfrac12\}\in(0,1)$ and apply $(P_2)$: there is
$N\in\mathbb{N}$ with $|x_n-a|<\varepsilon_1\le\varepsilon$ for all $n>N$. The
set of such $n$ is exactly $\{N+1,N+2,\dots\}$, so with $N_1=N+1$ we have
$|x_n-a|<\varepsilon$ for all $n\ge N_1$. Hence $x_n\to a$ and the equivalence is
true. (The change from $n\ge N$ to $n>N$ is immaterial, since $N$ is an arbitrary
positive integer and may be replaced by $N+1$.)

\medskip
\noindent\textbf{(3) True.}
Fix a positive constant $M>0$ and let $(P_3)$ denote the statement on the right.

Assume $(P_3)$ and let $\varepsilon>0$. Apply $(P_3)$ with
$\varepsilon_1=\varepsilon/M>0$: there is $N\in\mathbb{N}$ with
$|x_n-a|<M\varepsilon_1=\varepsilon$ for all $n>N$. Hence $x_n\to a$.

Conversely assume $x_n\to a$ and let $\varepsilon>0$. Apply the definition with
$\eta=\varepsilon/M>0$: there is $N\in\mathbb{N}$ with
$|x_n-a|<\eta=\varepsilon/M$ for all $n\ge N$, whence
$|x_n-a|<M\varepsilon$ for all $n\ge N$, and a fortiori for all $n>N$. Thus
$(P_3)$ holds. The equivalence is true: the constant $M$ can always be absorbed
into $\varepsilon$.

\medskip
\noindent\textbf{(4) True.}
Here $k$ ranges over $\mathbb{N}=\{1,2,3,\dots\}$, so that $1/k>0$. Let $(P_4)$
denote the statement on the right.

Assume $x_n\to a$ and let $k\in\mathbb{N}$. Apply the definition with
$\varepsilon=1/k>0$: there is $N\in\mathbb{N}$ with
$|x_n-a|<1/k$ for all $n\ge N$, hence for all $n>N$. Thus $(P_4)$ holds.

Conversely assume $(P_4)$ and let $\varepsilon>0$. By the Archimedean property of
$\mathbb{R}$ there exists $k\in\mathbb{N}$ with $1/k<\varepsilon$ (for instance
$k=\lfloor 1/\varepsilon\rfloor+1$, which satisfies $k>1/\varepsilon$, i.e.
$1/k<\varepsilon$). For this $k$, property $(P_4)$ gives $N\in\mathbb{N}$ with
$|x_n-a|<1/k<\varepsilon$ for all $n>N$. Hence $x_n\to a$, and the equivalence is
true.

\medskip
\noindent\textbf{(5) False.}
The condition on the right says only that $a$ is a \emph{limit point} (cluster
point) of the sequence: it requires the interval $(a-\varepsilon,a+\varepsilon)$
to catch infinitely many terms, but not to catch all but finitely many of them,
and the latter is what convergence needs.

Take $a=0$ and
\[
  x_n=\begin{cases}0, & n \text{ even},\\[2pt] 1, & n \text{ odd}.\end{cases}
\]
Then for every $\varepsilon>0$ the interval $(-\varepsilon,\varepsilon)$ contains
the terms $x_2=x_4=x_6=\cdots=0$, that is, infinitely many terms of the
sequence, so the statement on the right holds for $a=0$. However $\{x_n\}$ does
not converge to $0$: with $\varepsilon_0=\tfrac12$, for every $N\in\mathbb{N}$
there is an odd index $n\ge N$ with $|x_n-0|=1\ge\varepsilon_0$. (Equivalently,
$x_{2k}=0$ while $x_{2k-1}=1$, so the sequence has the two distinct subsequential
limits $0$ and $1$, and hence diverges.) Therefore the proposition is false.

The implication $\Leftarrow$ fails because from $\lnot(x_n\to a)$ one can only
conclude that for some $\varepsilon_0>0$ the set
$S=\{n\in\mathbb{N}:x_n\in(a-\varepsilon_0,a+\varepsilon_0)\}$ is not cofinite,
i.e.\ that its complement is infinite; this is much weaker than $S$ being
finite. In the example above, relative to $a=0$ and $\varepsilon_0=\tfrac12$,
this set is $S=\{2,4,6,\dots\}$, which is infinite but not cofinite.

\medskip
\noindent\textbf{Conclusion.} (1) True, (2) True, (3) True, (4) True,
(5) False.
\end{solution}

\begin{solution}{Solution 2.1.2}\label{xs:2.1.2}
We use the following standard facts. (a) \emph{Squeeze theorem}: if
$u_n\le w_n\le v_n$ for all $n$ and $u_n\to L$, $v_n\to L$, then $w_n\to L$.
(b) \emph{Algebra of limits}: convergent sequences are added, multiplied and
divided (when the denominator has nonzero limit and does not vanish) termwise,
and the limit of the result is the corresponding combination of the limits.
(c) \emph{Subsequences}: if $x_n\to L$ and $n_1<n_2<\cdots$ is strictly
increasing, then $x_{n_k}\to L$; conversely, if two subsequences of $\{x_n\}$
converge to different limits, then $\{x_n\}$ diverges.

\medskip
\noindent\textbf{(1) False.}
Take
\[
  a_n=b_n=c_n=n\qquad(n\in\mathbb{N}).
\]
Then $a_n\le b_n\le c_n$ for every $n$, and
\[
  c_n-a_n=0\qquad(n\in\mathbb{N}),
\]
so $\lim_{n\to\infty}(c_n-a_n)=0$. Nevertheless $\{a_n\}=\{n\}$ is unbounded and
therefore does not converge. The point is that $a_n\le b_n\le c_n$ together with
$\lim(c_n-a_n)=0$ only makes $a_n$ and $c_n$ asymptotically equal; they may both
tend to $+\infty$, so neither boundedness nor convergence of $\{a_n\}$ can be
concluded. Hence the proposition is false.

\medskip
\noindent\textbf{(2) True.}
Let $a_n\le b_n\le c_n$ for all $n$, let $b_n\to L$, and let
$\lim_{n\to\infty}(c_n-a_n)=0$.

From $a_n\le b_n\le c_n$ we get, for every $n$,
\[
  0\le b_n-a_n\le c_n-a_n
  \qquad\text{and}\qquad
  0\le c_n-b_n\le c_n-a_n .
\]
Let $\varepsilon>0$. Choose $N_1$ with $|b_n-L|<\varepsilon$ for $n\ge N_1$ and
$N_2$ with $|c_n-a_n|<\varepsilon$ for $n\ge N_2$, and put
$N=\max\{N_1,N_2\}$. Then for all $n\ge N$,
\[
  |a_n-b_n|=b_n-a_n\le c_n-a_n<\varepsilon,
  \qquad
  |c_n-b_n|=c_n-b_n\le c_n-a_n<\varepsilon .
\]
Hence, by the triangle inequality,
\[
  |a_n-L|\le|a_n-b_n|+|b_n-L|<\varepsilon+\varepsilon=2\varepsilon,
  \qquad
  |c_n-L|\le|c_n-b_n|+|b_n-L|<2\varepsilon
  \qquad(n\ge N).
\]
Since the splitting into $\varepsilon/2$ is available (replace $\varepsilon$ by
$\varepsilon/2$ in the choices of $N_1,N_2$), we may conclude
$|a_n-L|<\varepsilon$ for all $n\ge N$. As $\varepsilon>0$ is arbitrary,
\[
  \lim_{n\to\infty}a_n=L=\lim_{n\to\infty}b_n .
\]
Thus $\{a_n\}$ converges, and the proposition is true. (The hypothesis that
$\{b_n\}$ converges is essential: this is the squeeze theorem for a triple
$u_n\le w_n\le v_n$ in which only the middle sequence is known to converge and
the two outer ones are known to be asymptotically equal.)

\medskip
\noindent\textbf{(3) False.}
Let $a_n=\dfrac{(-1)^{n}}{n}$ for $n\in\mathbb{N}$. Then $a_n\ne0$ for every $n$
and $\{a_n\}$ converges (to $0$). But
\[
  \frac{a_{n+1}}{a_n}
  =\frac{(-1)^{n+1}/(n+1)}{(-1)^{n}/n}
  =-\frac{n}{n+1}\longrightarrow-1\ne1 .
\]
Hence the proposition is false. (If it is additionally assumed that
$\lim_{n\to\infty}a_n=L\ne0$, then the algebra of limits gives
\[
  \lim_{n\to\infty}\frac{a_{n+1}}{a_n}=\frac{L}{L}=1,
\]
so the statement becomes true; the hypothesis $a_n\ne0$ for all $n$ alone is not
enough, as the example above shows.)

\medskip
\noindent\textbf{(4) False.}
Take
\[
  a_n=\begin{cases}1, & n \text{ even},\\[2pt] 0, & n \text{ odd},\end{cases}
  \qquad
  b_n=\begin{cases}0, & n \text{ even},\\[2pt] 1, & n \text{ odd}.\end{cases}
\]
Then $a_nb_n=0$ for every $n$, so $\lim_{n\to\infty}a_nb_n=0$. However
$a_{2k}=1\to1$ and $a_{2k-1}=0\to0$, so the subsequences $\{a_{2k}\}$ and
$\{a_{2k-1}\}$ have different limits and $\{a_n\}$ does not converge (in
particular it does not converge to $0$); likewise $b_{2k}=0$ and
$b_{2k-1}=1$, so $\{b_n\}$ does not converge to $0$ either. Hence the
proposition is false. (Here neither $\lim a_n$ nor $\lim b_n$ even exists; the
hypothesis $\lim a_nb_n=0$ alone carries no information about the individual
limits.)

\medskip
\noindent\textbf{(5) True.}
Assume $\displaystyle\lim_{n\to\infty}\frac{b_n}{a_n}=1$ and that $\{a_n\}$
converges, say $a_n\to L\in\mathbb{R}$. For every $n$ the quotient $b_n/a_n$ is
defined, so $a_n\ne0$; moreover
\[
  b_n=a_n\cdot\frac{b_n}{a_n}\qquad(n\in\mathbb{N}).
\]
Both factors on the right converge, to $L$ and to $1$ respectively, so by the
product rule for limits $\{b_n\}$ converges and
\[
  \lim_{n\to\infty}b_n=\Bigl(\lim_{n\to\infty}a_n\Bigr)
  \Bigl(\lim_{n\to\infty}\frac{b_n}{a_n}\Bigr)=L\cdot1=L=\lim_{n\to\infty}a_n .
\]
Hence the proposition is true.

\medskip
\noindent\textbf{(6) False.}
Take
\[
  a_n=\begin{cases}1, & n \text{ even},\\[2pt] 0, & n \text{ odd}.\end{cases}
\]
Then $a_{2k}=1$ for every $k$, so $\{a_{2k}\}$ converges (to $1$), and
$a_{2k-1}=0$ for every $k$, so $\{a_{2k-1}\}$ converges (to $0$). But $\{a_n\}$
does not converge, since its even and odd subsequences have different limits.
Hence the proposition is false. (The correct criterion is: $\{a_n\}$ converges
if and only if $\{a_{2k}\}$ and $\{a_{2k-1}\}$ both converge and have the same
limit.)

\medskip
\noindent\textbf{Summary.} (1) False. (2) True. (3) False. (4) False.
(5) True. (6) False.
\end{solution}

\begin{solution}{Solution 2.1.3}\label{xs:2.1.3}
Recall the definition: $\displaystyle\lim_{n\to\infty}x_n=a$ means that for every
$\varepsilon>0$ there exists $N\in\mathbb{N}$ such that $|x_n-a|<\varepsilon$ for
every $n\ge N$.

\medskip
\noindent\textbf{(1)} We prove that
$\displaystyle\lim_{n\to\infty}\frac{1}{1+\sqrt n}=0$.

Let $\varepsilon>0$. By the Archimedean property of $\mathbb{R}$ choose
$N\in\mathbb{N}$ with $N>\dfrac{1}{\varepsilon^{2}}$. Then $\sqrt N>\dfrac1\varepsilon$.
For every $n\ge N$ we have $\sqrt n\ge\sqrt N$, hence
\[
  \left|\frac{1}{1+\sqrt n}-0\right|=\frac{1}{1+\sqrt n}
  \le\frac{1}{\sqrt n}\le\frac{1}{\sqrt N}<\varepsilon .
\]
Therefore $\displaystyle\lim_{n\to\infty}\frac{1}{1+\sqrt n}=0$.

\medskip
\noindent\textbf{(2)} We prove that
$\displaystyle\lim_{n\to\infty}\frac{3n^{2}+n}{2n^{2}-1}=\frac32$.

For $n\ge2$ (so that $2n^{2}-1>0$),
\[
  \left|\frac{3n^{2}+n}{2n^{2}-1}-\frac32\right|
  =\left|\frac{2(3n^{2}+n)-3(2n^{2}-1)}{2(2n^{2}-1)}\right|
  =\frac{2n+3}{4n^{2}-2} .
\]
For $n\ge2$ we have $2n+3\le2n+2n=4n$ and $4n^{2}-2\ge4n^{2}-n^{2}=3n^{2}$, hence
\[
  \left|\frac{3n^{2}+n}{2n^{2}-1}-\frac32\right|\le\frac{4n}{3n^{2}}
  =\frac{4}{3n}\le\frac{4}{3}\cdot\frac1n .
\]
Let $\varepsilon>0$ and choose $N\in\mathbb{N}$ with $N\ge2$ and
$N>\dfrac{4}{3\varepsilon}$. Then for every $n\ge N$,
\[
  \left|\frac{3n^{2}+n}{2n^{2}-1}-\frac32\right|\le\frac{4}{3n}
  \le\frac{4}{3N}<\varepsilon .
\]
Hence $\displaystyle\lim_{n\to\infty}\frac{3n^{2}+n}{2n^{2}-1}=\frac32$.

\medskip
\noindent\textbf{(3)} We prove that
$\displaystyle\lim_{n\to\infty}\bigl(\sqrt{n^{2}+n}-n\bigr)=\frac12$.

For every $n\ge1$, multiplying by the conjugate,
\[
  \sqrt{n^{2}+n}-n
  =\frac{\bigl(\sqrt{n^{2}+n}-n\bigr)\bigl(\sqrt{n^{2}+n}+n\bigr)}
        {\sqrt{n^{2}+n}+n}
  =\frac{n}{\sqrt{n^{2}+n}+n},
\]
so, using this same identity for the numerator,
\[
  0\le\frac12-\bigl(\sqrt{n^{2}+n}-n\bigr)
  =\frac12-\frac{n}{\sqrt{n^{2}+n}+n}
  =\frac{\sqrt{n^{2}+n}-n}{2\bigl(\sqrt{n^{2}+n}+n\bigr)}
  =\frac{n}{2\bigl(\sqrt{n^{2}+n}+n\bigr)^{2}} .
\]
Since $\sqrt{n^{2}+n}+n\ge2n$, we get
\[
  0\le\frac12-\bigl(\sqrt{n^{2}+n}-n\bigr)
  \le\frac{n}{2(2n)^{2}}=\frac{1}{8n} .
\]
Let $\varepsilon>0$ and choose $N\in\mathbb{N}$ with $N>\dfrac{1}{8\varepsilon}$.
Then for every $n\ge N$,
\[
  \left|\bigl(\sqrt{n^{2}+n}-n\bigr)-\frac12\right|
  =\frac12-\bigl(\sqrt{n^{2}+n}-n\bigr)\le\frac{1}{8n}\le\frac{1}{8N}<\varepsilon .
\]
Hence $\displaystyle\lim_{n\to\infty}\bigl(\sqrt{n^{2}+n}-n\bigr)=\frac12$.

\medskip
\noindent\textbf{(4)} We prove that $\displaystyle\lim_{n\to\infty}\arctan n=\frac\pi2$.

Since $\arctan n\in\bigl(0,\tfrac\pi2\bigr)$ for $n\ge1$, the number
$\dfrac\pi2-\arctan n$ lies in $\bigl(0,\tfrac\pi2\bigr)$. For
$\theta\in\bigl(0,\tfrac\pi2\bigr)$ the cotangent identity
$\tan\bigl(\tfrac\pi2-\theta\bigr)=\cot\theta=\dfrac{1}{\tan\theta}$ holds
(it follows from $\sin\bigl(\tfrac\pi2-\theta\bigr)=\cos\theta$ and
$\cos\bigl(\tfrac\pi2-\theta\bigr)=\sin\theta$). Taking $\theta=\arctan n$,
\[
  \tan\Bigl(\frac\pi2-\arctan n\Bigr)=\frac{1}{\tan(\arctan n)}=\frac1n .
\]
and since both angles lie in $\bigl(0,\tfrac\pi2\bigr)$, on which $\tan$ is
injective,
\[
  \frac\pi2-\arctan n=\arctan\frac1n .
\]
Finally $0<\arctan x<x$ for every $x>0$, so
\[
  0<\frac\pi2-\arctan n=\arctan\frac1n<\frac1n\qquad(n\ge1).
\]
Let $\varepsilon>0$ and choose $N\in\mathbb{N}$ with $N>\dfrac1\varepsilon$. For
every $n\ge N$,
\[
  \left|\arctan n-\frac\pi2\right|=\frac\pi2-\arctan n<\frac1n\le\frac1N<\varepsilon .
\]
Hence $\displaystyle\lim_{n\to\infty}\arctan n=\frac\pi2$.

\medskip
\noindent\textbf{(5)} We prove that
$\displaystyle\lim_{n\to\infty}\frac{n^{2}\arctan n}{1+n^{2}}=\frac\pi2$.

Write
\[
  \frac{n^{2}\arctan n}{1+n^{2}}-\frac\pi2
  =\frac{n^{2}\arctan n}{1+n^{2}}-\frac{n^{2}}{1+n^{2}}\cdot\frac\pi2
   -\frac{1}{1+n^{2}}\cdot\frac\pi2
  =\frac{n^{2}}{1+n^{2}}\Bigl(\arctan n-\frac\pi2\Bigr)-\frac{\pi}{2(1+n^{2})} .
\]
By part (4), $0<\dfrac\pi2-\arctan n<\dfrac1n$ for every $n\ge1$, and also
$0<\dfrac{n^{2}}{1+n^{2}}<1$. Hence, for every $n\ge1$,
\[
  \left|\frac{n^{2}\arctan n}{1+n^{2}}-\frac\pi2\right|
  \le\frac{n^{2}}{1+n^{2}}\Bigl(\frac\pi2-\arctan n\Bigr)
   +\frac{\pi}{2(1+n^{2})}
  <\frac1n+\frac{\pi}{2n}=\frac{2+\pi}{2n} .
\]
Let $\varepsilon>0$ and choose $N\in\mathbb{N}$ with
$N>\dfrac{2+\pi}{2\varepsilon}$. Then for every $n\ge N$,
\[
  \left|\frac{n^{2}\arctan n}{1+n^{2}}-\frac\pi2\right|
  <\frac{2+\pi}{2n}\le\frac{2+\pi}{2N}<\varepsilon .
\]
Hence $\displaystyle\lim_{n\to\infty}\frac{n^{2}\arctan n}{1+n^{2}}=\frac\pi2$.
\end{solution}

\begin{solution}{Solution 2.1.4}\label{xs:2.1.4}
We use the algebra of limits (sum, product and quotient rules), the squeeze
theorem, the standard limits
\[
  q^{1/n}\to1\ (q>0),\qquad n^{1/n}\to1,\qquad \frac1{n^{k}}\to0\ (k>0),
\]
and the geometric sum formula
$\displaystyle\sum_{i=1}^{n}q^{i}=\frac{q\,(1-q^{n})}{1-q}$ for $q\ne1$.

\medskip
\noindent\textbf{(1)} Divide numerator and denominator by $3^{n}$ and use
$3^{n+1}=3\cdot3^{n}$ together with $\bigl(-\tfrac23\bigr)^{n}\to0$ (because
$\bigl|\tfrac23\bigr|<1$):
\[
  \frac{(-2)^{n}+3^{n}}{(-2)^{n+1}+3^{n+1}}
  =\frac{\bigl(-\tfrac23\bigr)^{n}+1}{(-2)\bigl(-\tfrac23\bigr)^{n}+3}
  \longrightarrow\frac{0+1}{0+3}=\frac13 .
\]
The limit is $\boxed{\dfrac13}$.

\medskip
\noindent\textbf{(2)} By the geometric sum formula with $q=\tfrac12$ and with
$q=\tfrac13$,
\[
  \sum_{i=1}^{n}\frac{1}{2^{i}}
  =\frac{\tfrac12\bigl(1-2^{-n}\bigr)}{1-\tfrac12}=1-2^{-n},
  \qquad
  \sum_{i=1}^{n}\frac{1}{3^{i}}
  =\frac{\tfrac13\bigl(1-3^{-n}\bigr)}{1-\tfrac13}
  =\frac{1-3^{-n}}{2} .
\]
Since $2^{-n}\to0$ and $3^{-n}\to0$, the quotient of the two sums satisfies
\[
  \frac{\displaystyle\sum_{i=1}^{n}2^{-i}}{\displaystyle\sum_{i=1}^{n}3^{-i}}
  =\frac{1-2^{-n}}{\tfrac12\bigl(1-3^{-n}\bigr)}
  \longrightarrow\frac{1}{\tfrac12}=2 .
\]
The limit is $\boxed{2}$.

\medskip
\noindent\textbf{(3)} The sum in the parentheses consists of the first $n$ odd
numbers, and
\[
  1+3+5+\cdots+(2n-1)=n^{2} .
\]
Therefore
\[
  \frac{1}{n^{2}}+\frac{3}{n^{2}}+\cdots+\frac{2n-1}{n^{2}}
  =\frac{1+3+\cdots+(2n-1)}{n^{2}}=\frac{n^{2}}{n^{2}}=1
  \qquad(n\in\mathbb{N}),
\]
so the sequence is constant and the limit is $\boxed{1}$.

\medskip
\noindent\textbf{(4)} Denote the product by
\[
  P_n=\prod_{i=0}^{n}\Bigl(1+\frac{1}{2^{2^{i}}}\Bigr)
  =\Bigl(1+\frac12\Bigr)\Bigl(1+\frac{1}{2^{2}}\Bigr)
   \Bigl(1+\frac{1}{2^{4}}\Bigr)\cdots\Bigl(1+\frac{1}{2^{2^{n}}}\Bigr).
\]
The identity $(1-x)(1+x)=1-x^{2}$ gives, for each $i\ge0$,
\[
  \Bigl(1-\frac{1}{2^{2^{i}}}\Bigr)\Bigl(1+\frac{1}{2^{2^{i}}}\Bigr)
  =1-\frac{1}{2^{2^{i+1}}} .
\]
Apply this identity successively to the product $P_n$, starting with the single
factor $1-\tfrac12=1-\dfrac{1}{2^{2^{0}}}$: each step absorbs the next factor
$\Bigl(1+\dfrac{1}{2^{2^{i}}}\Bigr)$ of $P_n$ and doubles the exponent of the
remaining factor, since
\[
  \Bigl(1-\frac12\Bigr)\Bigl(1+\frac12\Bigr)=1-\frac{1}{2^{2}},\qquad
  \Bigl(1-\frac{1}{2^{2}}\Bigr)\Bigl(1+\frac{1}{2^{2}}\Bigr)=1-\frac{1}{2^{4}},
  \qquad\dots
\]
After all $n+1$ factors of $P_n$ have been absorbed, the exponent has been
doubled $n+1$ times, so
\[
  \Bigl(1-\frac12\Bigr)P_n=1-\frac{1}{2^{2^{n+1}}} .
\]
Since $1-\tfrac12=\tfrac12$,
\[
  P_n=2\Bigl(1-\frac{1}{2^{2^{n+1}}}\Bigr)
  \longrightarrow 2\,(1-0)=2 .
\]
The limit is $\boxed{2}$.

\medskip
\noindent\textbf{(5)} Let
$S_n=\displaystyle\sum_{k=1}^{n}\frac{1}{\sqrt{n^{2}+k}}$. For $1\le k\le n$ we
have $n^{2}<n^{2}+k\le n^{2}+n$, hence
\[
  \frac{1}{\sqrt{n^{2}+n}}\le\frac{1}{\sqrt{n^{2}+k}}\le\frac{1}{\sqrt{n^{2}+1}} .
\]
Summing these $n$ inequalities,
\[
  \frac{n}{\sqrt{n^{2}+n}}\le S_n\le\frac{n}{\sqrt{n^{2}+1}} .
\]
Now
\[
  \frac{n}{\sqrt{n^{2}+n}}=\frac{1}{\sqrt{1+\tfrac1n}}\longrightarrow1,
  \qquad
  \frac{n}{\sqrt{n^{2}+1}}=\frac{1}{\sqrt{1+\tfrac{1}{n^{2}}}}\longrightarrow1 .
\]
By the squeeze theorem $S_n\to1$. The limit is $\boxed{1}$.

\medskip
\noindent\textbf{(6)} Since $2>0$ we have $\sqrt[n]{2}=2^{1/n}\to1$, so
\[
  1-\frac{1}{\sqrt[n]{2}}=1-2^{-1/n}\longrightarrow1-1=0 .
\]
Also $|\cos n|\le1$ for every $n$, so $\{\cos n\}$ is bounded. By the theorem
that the product of a null sequence and a bounded sequence is a null sequence,
\[
  \Bigl(1-\frac{1}{\sqrt[n]{2}}\Bigr)\cos n\longrightarrow0 .
\]
The limit is $\boxed{0}$.

\medskip
\noindent\textbf{(7)} For $n\ge2$ the radicand satisfies
$1\le n^{2}-n+2\le n^{2}$, hence
\[
  1\le\sqrt[n]{n^{2}-n+2}\le\sqrt[n]{n^{2}}=\bigl(n^{1/n}\bigr)^{2} .
\]
Since $n^{1/n}\to1$, also $\bigl(n^{1/n}\bigr)^{2}\to1$; the squeeze theorem gives
\[
  \lim_{n\to\infty}\sqrt[n]{n^{2}-n+2}=1 .
\]
The limit is $\boxed{1}$.

\medskip
\noindent\textbf{(8)} For $n\ge3$ we have $\ln n\le n$, hence
$1\le n\ln n\le n^{2}$ and therefore
\[
  1\le\sqrt[n]{n\ln n}\le\sqrt[n]{n^{2}}=\bigl(n^{1/n}\bigr)^{2}\longrightarrow1 .
\]
By the squeeze theorem, $\displaystyle\lim_{n\to\infty}\sqrt[n]{n\ln n}=1$. The
limit is $\boxed{1}$.

\medskip
\noindent\textbf{(9)} Since $\arctan$ is increasing and $\arctan 2>1$ while
$\arctan n<\dfrac\pi2$, we have, for every $n\ge2$,
\[
  1\le\arctan n<\frac\pi2,
  \qquad\text{hence}\qquad
  1\le\sqrt[n]{\arctan n}\le\sqrt[n]{\frac\pi2}=\Bigl(\frac\pi2\Bigr)^{1/n} .
\]
(The inequality $1\le\arctan n$ does fail for the single index $n=1$, where
$\arctan1=\pi/4<1$; this affects neither the bounds for $n\ge2$ nor the limit.)
Since $\pi/2>0$, we have $(\pi/2)^{1/n}\to1$, and the squeeze theorem gives
$\displaystyle\lim_{n\to\infty}\sqrt[n]{\arctan n}=1$. The limit is
$\boxed{1}$.

\medskip
\noindent\textbf{(10)} The expression
\[
  T_n=\sqrt[n]{1}+\sqrt[n]{2}+\cdots+\sqrt[n]{10}
\]
is a sum of the ten fixed numbers $1,2,\dots,10$, each raised to the power
$1/n$. For every $n\ge1$ and $k=1,\dots,10$ we have $1\le\sqrt[n]{k}\le\sqrt[n]{10}$,
hence
\[
  10\le T_n\le 10\,\sqrt[n]{10}=10\cdot10^{1/n} .
\]
Since $10^{1/n}\to1$, the squeeze theorem gives $T_n\to10$. The limit is
$\boxed{10}$; equivalently,
\[
  \lim_{n\to\infty}\sum_{k=1}^{10}\sqrt[n]{k}=\sum_{k=1}^{10}1=10 .
\]

\medskip
\noindent\textbf{(11)} For every $n$,
\[
  1=\sin^{2}n+\cos^{2}n\le 2\sin^{2}n+\cos^{2}n=\sin^{2}n+1\le2,
\]
so
\[
  1\le\sqrt[n]{2\sin^{2}n+\cos^{2}n}\le\sqrt[n]{2}=2^{1/n}\longrightarrow1 .
\]
By the squeeze theorem,
$\displaystyle\lim_{n\to\infty}\sqrt[n]{2\sin^{2}n+\cos^{2}n}=1$. The limit is
$\boxed{1}$.

\medskip
\noindent\textbf{(12)} Let $\alpha\in(0,1)$ and put
$u_n=(n+1)^{\alpha}-n^{\alpha}$, $n\ge1$. Since $0<\alpha<1$, the function
$f(t)=t^{\alpha}$ is increasing on $[0,\infty)$, so $u_n>0$. By the mean value
theorem applied to $f$ on the interval $[n,n+1]$ there exists
$\xi_n\in(n,n+1)$ with
\[
  u_n=f(n+1)-f(n)=f'(\xi_n)=\alpha\,\xi_n^{\alpha-1} .
\]
Because $\xi_n>n>0$ and $\alpha-1<0$, we have
$\xi_n^{\alpha-1}<n^{\alpha-1}$, so
\[
  0<u_n<\alpha\,n^{\alpha-1}=\frac{\alpha}{n^{1-\alpha}} .
\]
Since $1-\alpha>0$, we have $n^{\alpha-1}=n^{-(1-\alpha)}\to0$; by the squeeze
theorem $u_n\to0$. The limit is $\boxed{0}$.

\emph{Elementary alternative without the mean value theorem.} For $n\ge1$ and
$0<\alpha<1$, the subadditivity-type inequality
$(1+x)^{\alpha}\le 1+x$ for $x\ge0$ gives
\[
  0<(n+1)^{\alpha}-n^{\alpha}
  =n^{\alpha}\Bigl[\Bigl(1+\frac1n\Bigr)^{\alpha}-1\Bigr]
  \le n^{\alpha}\cdot\frac1n=n^{\alpha-1}\longrightarrow0 ,
\]
and again $u_n\to0$ by the squeeze theorem.

\medskip
\noindent\textbf{(13)} Let $a_n\to a$ and put
$x_n=\dfrac{\lfloor na_n\rfloor}{n}$, where $\lfloor t\rfloor$ denotes the
greatest integer not exceeding $t$ (the book writes $[\,t\,]$ for this function).
For every real $t$ one has
\[
  t-1<\lfloor t\rfloor\le t .
\]
Applying this with $t=na_n$ and dividing by $n>0$ yields
\[
  a_n-\frac1n<\frac{\lfloor na_n\rfloor}{n}\le a_n
  \qquad(n\in\mathbb{N}).
\]
Now $a_n-\dfrac1n\to a-0=a$ and $a_n\to a$. By the squeeze theorem (the two
outer sequences converge to the same limit $a$),
\[
  \lim_{n\to\infty}\frac{\lfloor na_n\rfloor}{n}=a .
\]

\emph{The same conclusion in $\varepsilon$-$N$ language.} Let $\varepsilon>0$.
Choose $N_1$ with $|a_n-a|<\dfrac\varepsilon2$ for $n\ge N_1$ and $N_2$ with
$\dfrac1n<\dfrac\varepsilon2$ for $n\ge N_2$, and put $N=\max\{N_1,N_2\}$. For
$n\ge N$, the inequalities above give
\[
  a-\frac\varepsilon2-\frac\varepsilon2<a_n-\frac1n<x_n\le a_n<a+\frac\varepsilon2,
\]
that is, $a-\varepsilon<x_n<a+\varepsilon$; hence $|x_n-a|<\varepsilon$ for all
$n\ge N$. Therefore $\displaystyle\lim_{n\to\infty}\frac{\lfloor na_n\rfloor}{n}=a$.
\end{solution}

\begin{solution}{Solution 2.1.5}\label{xs:2.1.5}
Let
\[
  x_n=\frac{2n+(-1)^{n}n}{3n+1}\qquad(n\in\mathbb{N}),
\]
and note that
\[
  x_n=\frac{n\bigl(2+(-1)^{n}\bigr)}{3n+1}
  =\frac{2+(-1)^{n}}{3+\tfrac1n} .
\]

\medskip
\noindent\emph{Step 1: the even-indexed subsequence.}
For $n=2k$ we have $(-1)^{n}=1$, so
\[
  x_{2k}=\frac{2(2k)+2k}{3(2k)+1}=\frac{6k}{6k+1}
  =\frac{1}{1+\dfrac{1}{6k}} .
\]
Since $\dfrac{1}{6k}\to0$ as $k\to\infty$,
\[
  \lim_{k\to\infty}x_{2k}=1 .
\]

\medskip
\noindent\emph{Step 2: the odd-indexed subsequence.}
For $n=2k-1$ we have $(-1)^{n}=-1$, so
\[
  x_{2k-1}=\frac{2(2k-1)-(2k-1)}{3(2k-1)+1}
  =\frac{2k-1}{6k-2}
  =\frac{2-\dfrac1k}{6-\dfrac2k},
\]
where the last expression is obtained by dividing numerator and denominator by
$k$. Since $\dfrac1k\to0$ and $\dfrac2k\to0$,
\[
  \lim_{k\to\infty}x_{2k-1}=\frac{2}{6}=\frac13 .
\]

\medskip
\noindent\emph{Step 3: conclusion.}
Suppose, for contradiction, that $\{x_n\}$ converges, say
$\displaystyle\lim_{n\to\infty}x_n=L\in\mathbb{R}$. We use the theorem: \emph{if a
sequence converges to $L$, then every subsequence of it converges to $L$.}

Indeed, let $n_1<n_2<\cdots$ be a strictly increasing sequence of positive
integers and let $\varepsilon>0$. Choose $N$ with $|x_n-L|<\varepsilon$ for all
$n\ge N$. Since $n_k\ge k$ for every $k$ (proved by induction: $n_1\ge1$ and
$n_{k+1}>n_k\ge k$, so $n_{k+1}\ge k+1$), we have $n_k\ge N$ whenever $k\ge N$,
and therefore $|x_{n_k}-L|<\varepsilon$ for all $k\ge N$. Hence $x_{n_k}\to L$.

Applying this with the strictly increasing sequences $(2k)$ and $(2k-1)$, and
using Steps 1 and 2, we obtain
\[
  1=\lim_{k\to\infty}x_{2k}=L=\lim_{k\to\infty}x_{2k-1}=\frac13 ,
\]
which is impossible because $1\ne\tfrac13$. This contradiction shows that
$\{x_n\}$ does not converge; that is, the sequence diverges.

\medskip
\noindent\emph{Remark (direct $\varepsilon$ argument).} The same conclusion can
be reached without the subsequence theorem. Suppose $x_n\to L$ and put
$\varepsilon_0=\dfrac16>0$. Choose $N$ with $|x_n-L|<\varepsilon_0$ for $n\ge N$.
For $k$ large enough that $2k-1\ge N$ we have, by the triangle inequality,
\[
  \bigl|x_{2k}-x_{2k-1}\bigr|
  \le|x_{2k}-L|+|L-x_{2k-1}|<\frac16+\frac16=\frac13 .
\]
On the other hand
\[
  x_{2k}-x_{2k-1}=\frac{6k}{6k+1}-\frac{2k-1}{6k-2}
  =\frac{24k^{2}-8k+1}{(6k+1)(6k-2)}\longrightarrow\frac{24}{36}=\frac23 ,
\]
so for all sufficiently large $k$ we have $x_{2k}-x_{2k-1}>\dfrac13$, i.e.
$|x_{2k}-x_{2k-1}|>\dfrac13$. This is a contradiction. Hence $\{x_n\}$ diverges.

\medskip
\noindent Equivalently: the sequence $\{x_n\}$ has the two distinct
subsequential limits $1$ and $\dfrac13$, so it cannot converge.
\end{solution}

\begin{solution}{Solution 2.1.6}\label{xs:2.1.6}
We use the three hypotheses
\[
a_n\ne 0\quad(n\in\mathbb N),\qquad \frac{a_{n+1}}{a_n}>0\quad(n\in\mathbb N),
\qquad \lim_{n\to\infty}\frac{a_{n+1}}{a_n}=0 .
\]

\emph{Step 1: all terms have the same sign.} Since $\dfrac{a_{n+1}}{a_n}>0$ for
every $n$, two consecutive terms $a_n$ and $a_{n+1}$ always have the same sign.
By induction $a_n$ therefore has the same sign as $a_1$, for every
$n\in\mathbb N$; equivalently,
\[
a_n=a_1\prod_{k=1}^{n-1}\frac{a_{k+1}}{a_k},
\]
and this product is positive because every factor in it is positive.

\emph{Step 2: the ratios are eventually smaller than $1$.} Applying the
definition of $\displaystyle\lim_{n\to\infty}\frac{a_{n+1}}{a_n}=0$ with
$\varepsilon=1$, we obtain $N\in\mathbb N$ such that
\[
\left|\frac{a_{n+1}}{a_n}\right|<1\qquad(n>N).
\]
Together with the positivity of the ratios this says
\[
0<\frac{a_{n+1}}{a_n}<1\qquad(n>N).
\]

\emph{Step 3: the tail is monotone.} Fix $n>N$.
\begin{itemize}
  \item If $a_n>0$, then Step 1 gives $a_{n+1}>0$; multiplying the inequality
        $0<\dfrac{a_{n+1}}{a_n}<1$ by the positive number $a_n$ preserves it, so
        $0<a_{n+1}<a_n$.
  \item If $a_n<0$, then Step 1 gives $a_{n+1}<0$; multiplying by the negative
        number $a_n$ reverses the inequality, so $a_n<a_{n+1}<0$.
\end{itemize}
By Step 1 the sign is the same for every $n$, so one and the same case applies to
the whole tail: if $a_1>0$ then $a_{n+1}<a_n$ for every $n>N$, while if $a_1<0$
then $a_{n+1}>a_n$ for every $n>N$. Hence the tail $\{a_n\}_{n>N}$ is strictly
decreasing in the first case and strictly increasing in the second; in either
case there is an index $N$ such that the sequence $\{a_n\}$ is monotone for all
sufficiently large $n$, which is the required conclusion.
\end{solution}

\begin{solution}{Solution 2.1.7}\label{xs:2.1.7}
\emph{Part 1: a convergent sequence has a largest term or a smallest term.}

Let $\lim\limits_{n\to\infty}a_n=L$. By the definition of convergence with $\varepsilon=1$ there exists $N\in\mathbb{N}$ such that
\[
L-1<a_n<L+1\qquad(n>N).
\]
The set $S=\{a_n:n\in\mathbb{N}\}$ is therefore nonempty and bounded: if $N=0$ then $L-1<a_n<L+1$ for all $n$; if $N\ge 1$ then the finitely many numbers $a_1,\dots,a_N$ are bounded and the tail lies in $(L-1,L+1)$. Hence
\[
s=\sup S,\qquad i=\inf S
\]
are finite real numbers.

First note that $s\in S$ means exactly that $s=a_m$ for some index $m$, and then $a_n\le s=a_m$ for every $n$, so that $a_m$ is the largest term of the sequence; similarly, $i\in S$ means that $i$ is the smallest term. So it suffices to prove that $s\in S$ or $i\in S$.

If $s\in S$ there is nothing to prove. Assume therefore that $s\notin S$; we show that then $i\in S$. Because $s=\sup S$ is not an element of $S$, we have $a_n<s$ for all $n$. We prove that $s=L$: let $\delta>0$; since $s-\delta<s=\sup S$ and $s-\delta$ is not an upper bound of $S$, there is an index $n$ with $a_n>s-\delta$. Moreover, for each $k$ the set $\{n:a_n>s-\frac1k\}$ is infinite: if it were finite, then choosing $N$ larger than all of its elements we would have $a_n\le s-\frac1k$ for every $n\ge N$, so that $\max\{a_1,\dots,a_{N-1},s-\frac1k\}$ (the maximum over an empty list of the $a_j$ being read as $s-\frac1k$) would be an upper bound of $S$ strictly smaller than $s$, contradicting $s=\sup S$. We may therefore choose indices $n_1<n_2<\cdots$ with $a_{n_k}>s-\frac1k$ for every $k$; since also $a_{n_k}<s$, we get $|a_{n_k}-s|<\frac1k$, so $a_{n_k}\to s$. Every subsequence of the convergent sequence $\{a_n\}$ converges to the same limit $L$, so $L=s$. In particular, for every $\delta>0$ there are indices $n$ with $|a_n-L|<\delta$, that is, $a_n>L-\delta$.

Now suppose, for contradiction, that also $i\notin S$. Since $i$ is the greatest lower bound and no term equals $i$, for every $n$ we have $a_n>i$; moreover $i+\delta$ is not a lower bound of $S$ for any $\delta>0$, so there is an index $n$ with $a_n<i+\delta$, and a subsequence argument as before shows that $i=L$. Hence $L=i$ and $L=s$, so for every $n$,
\[
i=L\le a_n\le L=s ,
\]
which forces $a_n=L$ for all $n$. But then $L=a_1$ belongs to $S$ and $s=L$, so $s\in S$, contradicting the assumption $s\notin S$. Therefore $i\in S$, and the sequence has a smallest term.

Consequently, for every convergent sequence, $s\in S$ or $i\in S$: the sequence has a largest term or a smallest term.

\emph{Part 2: must it have both?} No. Consider the convergent sequence
\[
a_n=\frac1n\qquad(n\in\mathbb{N}),\qquad \lim_{n\to\infty}a_n=0 .
\]
It has the largest term $a_1=1$. It has no smallest term: for every $n$,
\[
a_{n+1}=\frac{1}{n+1}<\frac1n=a_n ,
\]
so no term is $\le$ all the others; equivalently, $\inf_n a_n=0$ is not attained, because $a_n>0$ for every $n$. Hence a convergent sequence need not have both a largest and a smallest term.
\end{solution}

\begin{solution}{Solution 2.1.8}\label{xs:2.1.8}
We use throughout the definition of convergence: $\lim\limits_{n\to\infty}x_n=a$ means that for every $\eta>0$ there exists $N_0\in\mathbb{N}$ such that $|x_n-a|<\eta$ whenever $n>N_0$.

\emph{(1)} We distinguish two cases according to the value of $a$.

\emph{Case $a>0$.} Taking $\eta=\dfrac a2$ in the definition of convergence, we get $N_1\in\mathbb{N}$ such that
\[
|x_n-a|<\frac a2\qquad(n>N_1),
\]
hence
\[
x_n>a-\frac a2=\frac a2>0\qquad(n>N_1),
\]
so in particular $x_n\ge 0$ and $\sqrt{x_n}$ is real for all $n>N_1$. Now let $\varepsilon>0$ be arbitrary. By the definition of convergence with $\eta=\varepsilon\sqrt a$ (which is positive because $a>0$) there is $N_2\in\mathbb{N}$ such that
\[
|x_n-a|<\varepsilon\sqrt a\qquad(n>N_2).
\]
Put $N=\max\{N_1,N_2\}$. For every $n>N$ the number $\sqrt{x_n}$ is real and, using the identity $\bigl(\sqrt{x_n}-\sqrt a\bigr)\bigl(\sqrt{x_n}+\sqrt a\bigr)=x_n-a$ and $\sqrt{x_n}+\sqrt a\ge\sqrt a>0$,
\[
\bigl|\sqrt{x_n}-\sqrt a\bigr|
=\frac{|x_n-a|}{\sqrt{x_n}+\sqrt a}
\le\frac{|x_n-a|}{\sqrt a}
<\frac{\varepsilon\sqrt a}{\sqrt a}=\varepsilon .
\]
Since $\varepsilon>0$ was arbitrary, $\lim\limits_{n\to\infty}\sqrt{x_n}=\sqrt a$.

\emph{Case $a=0$.} Let $\varepsilon>0$. Applying the definition of convergence with $\eta=\varepsilon^2$ we obtain $N\in\mathbb{N}$ such that $|x_n|<\varepsilon^2$ for all $n>N$. For such $n$ we have $0\le x_n<\varepsilon^2$, hence
\[
\bigl|\sqrt{x_n}-\sqrt 0\bigr|=\sqrt{x_n}<\varepsilon .
\]
Therefore $\lim\limits_{n\to\infty}\sqrt{x_n}=0=\sqrt a$.

(The case $a<0$ cannot occur here, because then $x_n<0$ for all large $n$ and $\sqrt{x_n}$ would not be real. Observe that the hypothesis $a\ne 0$ is stronger than necessary: the assertion is true for every $a\ge 0$.)

\emph{(2)} We need the elementary inequality for the sine function,
\[
|\sin u-\sin v|\le|u-v|\qquad(u,v\in\mathbb{R}),
\]
which follows from the product formula $\sin u-\sin v=2\cos\dfrac{u+v}{2}\sin\dfrac{u-v}{2}$ combined with $|\cos t|\le 1$ and $|\sin t|\le|t|$.

Let $\varepsilon>0$ be arbitrary. Since $a\ne 0$, taking $\eta=\dfrac{|a|}{2}$ in the definition of convergence we get $N_1\in\mathbb{N}$ such that
\[
|x_n-a|<\frac{|a|}{2}\qquad(n>N_1),
\]
so that for $n>N_1$,
\[
|x_n|>|a|-\frac{|a|}{2}=\frac{|a|}{2},\qquad
|x_n+a|\le|x_n|+|a|<\frac{|a|}{2}+|a|=\frac{3|a|}{2}.
\]
In particular $x_n\ne 0$ for $n>N_1$, so $\dfrac{1}{x_n^2}$ is defined for those $n$. For $n>N_1$,
\[
\left|\frac{1}{x_n^2}-\frac{1}{a^2}\right|
=\frac{|a^2-x_n^2|}{x_n^2a^2}
=\frac{|x_n-a|\,|x_n+a|}{x_n^2a^2}
<|x_n-a|\cdot\frac{\frac{3|a|}{2}}{\frac{|a|^2}{4}\cdot a^2}
=|x_n-a|\cdot\frac{6}{|a|^3},
\]
where we used $a^2=|a|^2$ and $x_n^2>\dfrac{|a|^2}{4}$. Choose $N_2\in\mathbb{N}$ such that
\[
|x_n-a|<\frac{|a|^3}{6}\,\varepsilon\qquad(n>N_2)
\]
(this is possible by the definition of convergence, since $\dfrac{|a|^3}{6}\varepsilon>0$). Put $N=\max\{N_1,N_2\}$. For every $n>N$ we have $x_n\ne 0$ and, by the sine inequality with $u=\dfrac{1}{x_n^2}$ and $v=\dfrac{1}{a^2}$,
\[
\left|\sin\frac{1}{x_n^2}-\sin\frac{1}{a^2}\right|
\le\left|\frac{1}{x_n^2}-\frac{1}{a^2}\right|
<|x_n-a|\cdot\frac{6}{|a|^3}
<\frac{|a|^3}{6}\varepsilon\cdot\frac{6}{|a|^3}
=\varepsilon .
\]
Since $\varepsilon>0$ was arbitrary, $\lim\limits_{n\to\infty}\sin\dfrac{1}{x_n^2}=\sin\dfrac{1}{a^2}$.
\end{solution}

\begin{solution}{Solution 2.1.9}\label{xs:2.1.9}
Let $\varepsilon>0$ be arbitrary. By the hypothesis $\lim\limits_{n\to\infty}(x_n-x_{n-1})=d$, applied with the positive number $\varepsilon$, there is $N_1\in\mathbb{N}$ such that
\[
\bigl|(x_n-x_{n-1})-d\bigr|<\varepsilon\qquad(n>N_1).
\]
Fix an index $n_0>N_1$, for instance $n_0=N_1+1$. For every $n>n_0$ the telescoping identity
\[
x_n=x_{n_0}+\sum_{k=n_0+1}^{n}(x_k-x_{k-1})
\]
holds. Subtracting $nd$, where $d$ is summed $n-n_0$ times over the indices $k=n_0+1,\dots,n$ and $n_0$ times separately, we get
\[
x_n-nd=x_{n_0}-n_0d+\sum_{k=n_0+1}^{n}\bigl[(x_k-x_{k-1})-d\bigr].
\]
Dividing by $n>0$ and taking absolute values,
\[
\left|\frac{x_n}{n}-d\right|
\le\frac{|x_{n_0}-n_0d|}{n}+\frac{1}{n}\sum_{k=n_0+1}^{n}\bigl|(x_k-x_{k-1})-d\bigr| .
\]
Every index $k$ occurring in the sum satisfies $k\ge n_0+1>N_1$, so each summand is $<\varepsilon$; there are $n-n_0<n$ of them. Therefore, for every $n>n_0$,
\[
\left|\frac{x_n}{n}-d\right|
<\frac{C}{n}+\frac{n-n_0}{n}\,\varepsilon
<\frac Cn+\varepsilon,
\qquad C:=|x_{n_0}-n_0d| .
\]
The number $C\ge 0$ is a fixed constant and $C/n\to 0$; hence there is $N_2\in\mathbb{N}$ with $C/n<\varepsilon$ for all $n>N_2$. Putting $N=\max\{n_0,N_2\}$, we obtain for every $n>N$
\[
\left|\frac{x_n}{n}-d\right|<\varepsilon+\varepsilon=2\varepsilon .
\]
Since $\varepsilon>0$ was arbitrary, replacing $\varepsilon$ by $\varepsilon/2$ at the beginning of the argument gives $\bigl|\frac{x_n}{n}-d\bigr|<\varepsilon$ for all large $n$. Hence
\[
\lim_{n\to\infty}\frac{x_n}{n}=d .
\]
\end{solution}

\begin{solution}{Solution 2.1.10}\label{xs:2.1.10}
\emph{Remark on the hypotheses.} Since $a_n>0$ for every $n$ and $a_n\to a$, we have $a\ge 0$. Moreover each product $a_1a_2\cdots a_n$ is positive, so all the $n$-th roots below are real and positive.

\emph{Main assertion: $\lim\limits_{n\to\infty}\sqrt[n]{a_1a_2\cdots a_n}=a$.} We split the proof into two cases.

\emph{Case 1: $a>0$.} Let $0<\varepsilon<a$ be arbitrary. Choose a number $M$ with $a+\varepsilon<M<a+2\varepsilon$. By the definition of convergence there is $N\in\mathbb{N}$ such that
\[
a-\varepsilon<a_n<M\qquad(n>N).
\]
Put $P=\displaystyle\prod_{k=1}^{N}a_k>0$ (if $N=0$, read $P=1$). For every $n>N$,
\[
a_1a_2\cdots a_n=P\prod_{k=N+1}^{n}a_k ,
\]
and each of the $n-N$ factors of the second product lies strictly between $a-\varepsilon$ and $M$, whence
\[
P\,(a-\varepsilon)^{\,n-N}<a_1a_2\cdots a_n<P\,M^{\,n-N}.
\]
Taking $n$-th roots (the map $t\mapsto t^{1/n}$ is increasing on $[0,+\infty)$) gives, for $n>N$,
\[
P^{\frac1n}(a-\varepsilon)^{1-\frac Nn}<\sqrt[n]{a_1a_2\cdots a_n}<P^{\frac1n}M^{1-\frac Nn}.
\]
Now $P^{\frac1n}\to 1$ for every $P>0$: if $P\ge 1$ then Bernoulli's inequality gives $1\le P^{\frac1n}\le 1+\frac{P-1}{n}$, and if $0<P<1$ the same bound applied to $1/P$ gives $P^{\frac1n}\ge\frac{1}{1+\frac{1/P-1}{n}}$. Also
\[
(a-\varepsilon)^{1-\frac Nn}\to a-\varepsilon,\qquad M^{1-\frac Nn}\to M .
\]
Hence there is $N_1>N$ such that for every $n>N_1$,
\[
(a-\varepsilon)-\varepsilon<P^{\frac1n}(a-\varepsilon)^{1-\frac Nn},
\qquad
P^{\frac1n}M^{1-\frac Nn}<M+\varepsilon .
\]
Combining the last three displays and using $M+\varepsilon<a+2\varepsilon$ and $(a-\varepsilon)-\varepsilon=a-2\varepsilon$, we conclude that for every $n>N_1$,
\[
a-2\varepsilon<\sqrt[n]{a_1a_2\cdots a_n}<a+2\varepsilon ,
\]
that is, $\bigl|\sqrt[n]{a_1a_2\cdots a_n}-a\bigr|<2\varepsilon$. Since $\varepsilon>0$ was arbitrary, the main assertion holds in this case.

\emph{Case 2: $a=0$.} Let $\varepsilon>0$. By the definition of convergence
there is $N\in\mathbb{N}$ with
\[
0<a_n<\frac\varepsilon2\qquad(n>N).
\]
Put $Q=\displaystyle\prod_{k=1}^{N}a_k>0$ (if $N=0$, $Q=1$). Since every factor
$a_k$ with $N<k\le n$ is $<\dfrac\varepsilon2$, we get for every $n>N$
\[
a_1a_2\cdots a_n=Q\prod_{k=N+1}^{n}a_k<Q\,\Bigl(\frac\varepsilon2\Bigr)^{\,n-N},
\]
hence
\[
0<\sqrt[n]{a_1a_2\cdots a_n}
<\Bigl(Q\,\Bigl(\frac\varepsilon2\Bigr)^{\,n-N}\Bigr)^{\frac1n}
=\frac\varepsilon2\cdot Q^{\frac1n}\Bigl(\frac\varepsilon2\Bigr)^{-\frac Nn}.
\]
As $n\to\infty$ we have $Q^{\frac1n}\to 1$ and
$\Bigl(\dfrac\varepsilon2\Bigr)^{-\frac Nn}\to 1$, hence
\[
Q^{\frac1n}\Bigl(\frac\varepsilon2\Bigr)^{-\frac Nn}\longrightarrow 1 ;
\]
in particular there is $N_1>N$ such that
$Q^{\frac1n}\bigl(\frac\varepsilon2\bigr)^{-\frac Nn}<2$ for every $n>N_1$
(convergence to $1$ gives, for instance, values $<2$ eventually), and then
\[
0<\sqrt[n]{a_1a_2\cdots a_n}<\frac\varepsilon2\cdot 2=\varepsilon .
\]
Thus $\bigl|\sqrt[n]{a_1a_2\cdots a_n}-0\bigr|<\varepsilon$ for all $n>N_1$. Since $\varepsilon>0$ was arbitrary,
\[
\lim_{n\to\infty}\sqrt[n]{a_1a_2\cdots a_n}=0=a .
\]
This proves the main assertion.

\emph{(1)} Assume $\lim\limits_{n\to\infty}\dfrac{a_{n+1}}{a_n}=a$, where $a_n>0$ for all $n\in\mathbb{N}$. Apply the main assertion to the positive sequence
\[
r_k=\frac{a_{k+1}}{a_k}\qquad(k\in\mathbb{N}),
\]
whose limit is $a$. This yields
\[
\lim_{m\to\infty}\sqrt[m]{r_1r_2\cdots r_m}=a .
\]
But the product telescopes:
\[
r_1r_2\cdots r_{n-1}
=\frac{a_2}{a_1}\cdot\frac{a_3}{a_2}\cdots\frac{a_n}{a_{n-1}}
=\frac{a_n}{a_1},
\]
and therefore
\[
\sqrt[n]{a_n}=\sqrt[n]{a_1\cdot\frac{a_n}{a_1}}
=a_1^{\frac1n}\bigl(r_1r_2\cdots r_{n-1}\bigr)^{\frac1n}.
\]
The sequence $\{(r_1r_2\cdots r_{n-1})^{1/n}\}_{n\ge 2}$ is the sequence $\{(r_1r_2\cdots r_m)^{1/m}\}_m$ with $m=n-1$, raised to the powers $\frac{n-1}{n}\to 1$; indeed,
\[
\bigl(r_1r_2\cdots r_{n-1}\bigr)^{\frac1n}
=\bigl[\bigl(r_1r_2\cdots r_{n-1}\bigr)^{\frac{1}{n-1}}\bigr]^{\frac{n-1}{n}},
\]
and $\frac{n-1}{n}\to 1$. Since $t\mapsto t^{\alpha}$ is continuous on $(0,+\infty)$ for every fixed real $\alpha$, and since the limit of a sequence and of any of its subsequences coincide, we conclude that
\[
\lim_{n\to\infty}\bigl(r_1r_2\cdots r_{n-1}\bigr)^{\frac1n}=a .
\]
Finally $a_1^{1/n}\to 1$, so
\[
\lim_{n\to\infty}\sqrt[n]{a_n}=1\cdot a=a .
\]

If $a=0$, the same computation applies: by Case 2 of the main assertion,
\[
\sqrt[m]{r_1r_2\cdots r_m}\longrightarrow 0 ,
\]
and $a_1^{1/n}\to 1$ still holds because $a_1>0$.

\emph{(2)} Since $n!=1\cdot2\cdots n>0$ for every $n$, define
\[
b_n=\frac{n!}{n^n}>0\qquad(n\in\mathbb{N}).
\]
Then
\[
\sqrt[n]{b_n}=\frac{\sqrt[n]{n!}}{n},
\]
because $\sqrt[n]{n^n}=n$. For the successive ratio of the sequence $\{b_n\}$ we compute
\[
\frac{b_{n+1}}{b_n}
=\frac{(n+1)!}{(n+1)^{n+1}}\cdot\frac{n^n}{n!}
=\frac{(n+1)\,n!}{(n+1)^{n+1}}\cdot\frac{n^n}{n!}
=\Bigl(\frac{n}{n+1}\Bigr)^{n}
=\frac{1}{\Bigl(1+\dfrac1n\Bigr)^{n}} .
\]
By the fundamental limit $\lim\limits_{n\to\infty}\Bigl(1+\dfrac1n\Bigr)^{n}=\mathrm{e}$ (where $\mathrm{e}$ denotes the base of the natural logarithm), we obtain
\[
\lim_{n\to\infty}\frac{b_{n+1}}{b_n}=\frac1{\mathrm{e}} .
\]
The sequence $\{b_n\}$ is positive, so part (1) with $a=\dfrac1{\mathrm{e}}$ gives
\[
\lim_{n\to\infty}\sqrt[n]{b_n}=\frac1{\mathrm{e}} .
\]
Substituting $\sqrt[n]{b_n}=\dfrac{\sqrt[n]{n!}}{n}$ yields
\[
\lim_{n\to\infty}\frac{\sqrt[n]{n!}}{n}=\frac1{\mathrm{e}} .
\]
\end{solution}

\begin{solution}{Solution 2.1.11}\label{xs:2.1.11}
Let $a$ be the common limit:
\[
\lim_{k\to\infty}a_{3k-2}=\lim_{k\to\infty}a_{3k-1}=\lim_{k\to\infty}a_{3k}=a .
\]
Let $\varepsilon>0$ be arbitrary. By the definition of convergence applied to each of the three subsequences, there exist $K_1,K_2,K_3\in\mathbb{N}$ such that
\[
|a_{3k-2}-a|<\varepsilon\ \ (k>K_1),\qquad
|a_{3k-1}-a|<\varepsilon\ \ (k>K_2),\qquad
|a_{3k}-a|<\varepsilon\ \ (k>K_3).
\]
Put
\[
N=\max\{3K_1-2,\;3K_2-1,\;3K_3\}.
\]
Every positive integer $n$ is of exactly one of the three forms $3k-2$, $3k-1$, $3k$ (division with remainder by $3$). Let $n>N$.

If $n=3k-2$, then $3k-2=n>N\ge 3K_1-2$, hence $k>K_1$ and therefore $|a_n-a|=|a_{3k-2}-a|<\varepsilon$.

If $n=3k-1$, then $3k-1=n>N\ge 3K_2-1$, hence $k>K_2$ and therefore $|a_n-a|=|a_{3k-1}-a|<\varepsilon$.

If $n=3k$, then $3k=n>N\ge 3K_3$, hence $k>K_3$ and therefore $|a_n-a|=|a_{3k}-a|<\varepsilon$.

Thus $|a_n-a|<\varepsilon$ for every $n>N$, that is, $\lim\limits_{n\to\infty}a_n=a$; the sequence $\{a_n\}$ converges.
\end{solution}

\begin{solution}{Solution 2.1.12}\label{xs:2.1.12}
Recall the definition: $\lim\limits_{n\to\infty}x_n=+\infty$ means that for every $M>0$ there is $N\in\mathbb{N}$ such that $x_n>M$ for all $n>N$.

Let $\sigma_n=\dfrac{x_1+x_2+\cdots+x_n}{n}$. Fix an arbitrary $M>0$. Applying the definition with the positive number $2M$, we obtain $N_1\in\mathbb{N}$ such that
\[
x_n>2M\qquad(n>N_1).
\]
Put $C=x_1+x_2+\cdots+x_{N_1}$ (a fixed real number). For every $n>N_1$,
\[
\sigma_n=\frac{1}{n}\Bigl(\sum_{k=1}^{N_1}x_k+\sum_{k=N_1+1}^{n}x_k\Bigr)
>\frac{1}{n}\bigl(C+(n-N_1)\cdot 2M\bigr)
=2M-\frac{2MN_1-C}{n}.
\]
Choose $N\ge N_1$ so large that $\dfrac{2MN_1-C}{n}<M$ for all $n>N$; this is possible because $\dfrac{2MN_1-C}{n}\to 0$ as $n\to\infty$ (if $2MN_1-C\le 0$, take $N=N_1$). Then for every $n>N$,
\[
\sigma_n>2M-M=M .
\]
Since $M>0$ was arbitrary, this proves $\lim\limits_{n\to\infty}\dfrac{x_1+x_2+\cdots+x_n}{n}=+\infty$.
\end{solution}

\begin{solution}{Solution 2.1.13}\label{xs:2.1.13}
Suppose, for contradiction, that $\{a_n\}$ is bounded. By hypothesis $\dfrac{a_n}{a_{n+1}+a_{n+2}}\to 0$. Taking $\varepsilon=\dfrac14$ in the definition of convergence (the limit is $0$), we obtain $N_0\in\mathbb{N}$ such that
\[
0<\frac{a_n}{a_{n+1}+a_{n+2}}<\frac14\qquad(n\ge N_0),
\]
that is,
\[
a_{n+1}+a_{n+2}>4a_n\qquad(n\ge N_0).\qquad(1)
\]
Because $\{a_n\}$ is bounded and $a_n>0$ for all $n$, the number
\[
s=\sup_{n\ge N_0}a_n
\]
is finite and positive. Since $s$ is the least upper bound of the nonempty set $\{a_n:n\ge N_0\}$ and $s>0$, there exists an index $m\ge N_0$ with $a_m>\dfrac{s}{2}$: indeed, otherwise $a_n\le \dfrac{s}{2}$ for all $n\ge N_0$, so $\dfrac{s}{2}<s$ would be a smaller upper bound, contradicting the definition of $s$.

Applying $(1)$ with $n=m$ and using $a_{m+1}\le s$, $a_{m+2}\le s$, we get
\[
2s\ \ge\ a_{m+1}+a_{m+2}\ >\ 4a_m\ >\ 4\cdot\frac{s}{2}=2s,
\]
which is impossible. Hence $\{a_n\}$ cannot be bounded; the sequence $\{a_n\}$ is unbounded.
\end{solution}

\begin{solution}{Solution 2.1.14}\label{xs:2.1.14}
Recall the definitions. A sequence $\{a_n\}$ is \emph{unbounded} if for every $M>0$ there exists $n\in\mathbb{N}$ with $|a_n|>M$. A sequence $\{a_n\}$ is an \emph{infinitely large quantity}, written $\lim\limits_{n\to\infty}a_n=\infty$, if for every $M>0$ there exists $N\in\mathbb{N}$ such that $|a_n|>M$ for all $n>N$. A subsequence of $\{a_n\}$ is a sequence $\{a_{n_k}\}$ where $n_1<n_2<\cdots$ are positive integers.

\emph{Necessity.} Assume that $\{a_n\}$ is an infinitely large quantity, and let $\{a_{n_k}\}$ be any subsequence. Given $M>0$, choose $N$ with $|a_n|>M$ whenever $n>N$. Since $n_k\ge k$, for every $k>N$ we have $n_k\ge k>N$, hence $|a_{n_k}|>M$. Therefore the subsequence $\{a_{n_k}\}$ is unbounded. As the subsequence was arbitrary, every subsequence of $\{a_n\}$ is unbounded.

\emph{Sufficiency.} Assume that every subsequence of $\{a_n\}$ is unbounded, and suppose, for contradiction, that $\{a_n\}$ is not an infinitely large quantity. Then the defining property fails: there exists $M>0$ such that for every $N\in\mathbb{N}$ there is some $n>N$ with $|a_n|\le M$.

We construct a subsequence recursively. Choose $n_1>0$ with $|a_{n_1}|\le M$ (take $N=0$). Having chosen $n_k$, take $N=n_k$ and pick $n_{k+1}>n_k$ with $|a_{n_{k+1}}|\le M$. This produces indices $n_1<n_2<\cdots$ with $|a_{n_k}|\le M$ for every $k$, so the subsequence $\{a_{n_k}\}$ is bounded (by $M$), contradicting the hypothesis that every subsequence is unbounded.

Therefore $\{a_n\}$ is an infinitely large quantity, and the stated condition is both necessary and sufficient.
\end{solution}

\begin{solution}{Solution 2.1.15}\label{xs:2.1.15}
Assume, to the contrary, that $|l|>1$. Choose a real number $c$ with
\[
1<c<|l|;
\]
such a $c$ exists because $|l|>1$ (take, for instance, $c=\dfrac{1+|l|}{2}$).

Since $\dfrac{a_{n+1}}{a_n}\to l$, we have $\left|\dfrac{a_{n+1}}{a_n}\right|\to |l|$ (the absolute value function is continuous). Choose $N\in\mathbb{N}$ such that
\[
\left|\frac{a_{n+1}}{a_n}\right|>c\qquad(n\ge N).
\]
Note that $a_N\ne 0$ by hypothesis, so $|a_N|>0$. For every $n>N$, writing the quotient as a telescoping product,
\[
|a_n|=|a_N|\cdot\prod_{k=N}^{n-1}\left|\frac{a_{k+1}}{a_k}\right|>|a_N|\,c^{\,n-N}.
\]
Since $c>1$, we have $c^{\,n-N}\to+\infty$ (for example, $c=1+r$ with $r>0$, and $c^m\ge 1+mr\to+\infty$ by Bernoulli's inequality). Hence, for every $M>0$, choosing $n$ so large that $|a_N|c^{\,n-N}>M$ gives $|a_n|>M$. Therefore $|a_n|\to+\infty$, and in particular the sequence $\{a_n\}$ does not converge to $0$, contradicting the hypothesis $\lim\limits_{n\to\infty}a_n=0$.

Consequently $|l|\le 1$.
\end{solution}

\begin{solution}{Solution 2.1.16}\label{xs:2.1.16}
Let $\varepsilon>0$ be arbitrary. We distinguish two cases.

If $\varepsilon\ge b-a$, then for every $n$ both $x_n$ and $a$ lie in $[a,b]$, so
\[
|x_n-a|\le b-a\le\varepsilon ,
\]
and the required inequality holds automatically for that $\varepsilon$.

Now suppose $0<\varepsilon<b-a$. Then $a<a+\varepsilon<b$, and since $f$ is strictly increasing on $[a,b]$,
\[
\delta:=f(a+\varepsilon)-f(a)>0 .
\]
By hypothesis $f(x_n)\to f(a)$, so there exists $N\in\mathbb{N}$ such that
\[
|f(x_n)-f(a)|<\delta\qquad(n>N),
\]
in particular $f(x_n)<f(a)+\delta=f(a+\varepsilon)$ for all $n>N$. Since $x_n\in[a,b]$ and $f$ is strictly increasing, $f(x_n)<f(a+\varepsilon)$ forces $x_n<a+\varepsilon$: indeed, $x_n\ge a+\varepsilon$ would imply $f(x_n)\ge f(a+\varepsilon)$. On the other hand $x_n\ge a$ because $\{x_n\}\subset[a,b]$. Hence
\[
0\le x_n-a<\varepsilon,\qquad\text{that is,}\qquad |x_n-a|<\varepsilon
\]
for all $n>N$.

Thus for every $\varepsilon>0$ there is $N$ with $|x_n-a|<\varepsilon$ whenever $n>N$, which is exactly the statement $\lim\limits_{n\to\infty}x_n=a$.
\end{solution}

\begin{solution}{Solution 2.1.17}\label{xs:2.1.17}
Write
\[
\sigma_n=\frac{x_1+x_2+\cdots+x_n}{n},\qquad \sigma_n\to a .
\]
Since $\{x_n\}$ is monotonically increasing, $x_k\le x_n$ for all $k\le n$, hence for every $n$
\[
\sigma_n=\frac{1}{n}\sum_{k=1}^{n}x_k\le\frac{1}{n}\sum_{k=1}^{n}x_n=x_n. \qquad(1)
\]
Also, for every fixed $m$ and every $n>m$, monotonicity gives $x_k\ge x_m$ for $m<k\le n$, so
\[
\sigma_n=\frac{1}{n}\Bigl(\sum_{k=1}^{m}x_k+\sum_{k=m+1}^{n}x_k\Bigr)
\ge \frac{1}{n}\sum_{k=1}^{m}x_k+\frac{n-m}{n}\,x_m .
\]
The first term on the right tends to $0$ and $\frac{n-m}{n}\to 1$ as $n\to\infty$, so the lower limit of the right-hand side equals $x_m$. Since $\sigma_n\to a$, we may pass to lower limits in the inequality above and obtain
\[
a\ \ge\ x_m\qquad\text{for every }m\in\mathbb{N}. \qquad(2)
\]

\emph{Case 1: $a$ is finite.} By $(2)$ the monotonically increasing sequence $\{x_n\}$ is bounded above by $a$; hence it converges to its supremum $L=\sup_n x_n$, and $L\le a$. On the other hand, $(1)$ gives $\sigma_n\le x_n$ for all $n$, so taking upper limits,
\[
a=\limsup_{n\to\infty}\sigma_n\le\limsup_{n\to\infty}x_n=L .
\]
Therefore $a\le L\le a$, that is $L=a$, and consequently $\lim\limits_{n\to\infty}x_n=a$.

\emph{Case 2: $a=+\infty$.} Then $(1)$ gives $x_n\ge\sigma_n$ for every $n$, hence
\[
\liminf_{n\to\infty}x_n\ge\lim_{n\to\infty}\sigma_n=+\infty,
\]
so $\lim\limits_{n\to\infty}x_n=+\infty=a$.

(The case $a=-\infty$ cannot occur, since $\sigma_n\ge x_1>-\infty$ for all $n$.) In every possible case, $\lim\limits_{n\to\infty}x_n=a$.
\end{solution}


\srcline{Exercise set 2.2, problems 1--6 --- printed page 38}

\begin{solution}{Solution 2.2.1}\label{xs:2.2.1}
The converse of the Stolz theorem (Theorem 2.2.1) would read as follows: if $\{a_n\}$ is strictly increasing with $a_n\to+\infty$, if $\{b_n\}$ is an arbitrary sequence and if
\[
\lim_{n\to\infty}\frac{b_n}{a_n}=a
\]
exists, then necessarily
\[
\lim_{n\to\infty}\frac{b_n-b_{n-1}}{a_n-a_{n-1}}=a .
\]
We show that this assertion is false.

Take
\[
a_n=n,\qquad b_n=(-1)^n,\qquad n\ge 1 .
\]
For every $n\ge2$ we have $a_n-a_{n-1}=1>0$, so $\{a_n\}$ is strictly increasing, and $a_n=n\to+\infty$; thus all the hypotheses on the denominator are satisfied. Moreover
\[
\frac{b_n}{a_n}=\frac{(-1)^n}{n},\qquad\text{hence}\qquad
\left|\frac{b_n}{a_n}-0\right|=\frac{1}{n}\longrightarrow 0 ,
\]
so that $\lim_{n\to\infty}b_n/a_n=0$; the hypothesis of the converse therefore holds with $a=0$.

On the other hand, since $(-1)^{n-1}=-(-1)^n$,
\[
\frac{b_n-b_{n-1}}{a_n-a_{n-1}}
=\frac{(-1)^n-(-1)^{n-1}}{1}
=(-1)^n\bigl(1+1\bigr)
=2(-1)^n .
\]
Thus this increment quotient equals $2$ when $n$ is even and $-2$ when $n$ is odd. Such a sequence has no limit: its even-indexed subsequence converges to $2$ and its odd-indexed subsequence converges to $-2$, while every subsequence of a convergent sequence must converge to the same limit. Consequently the increment quotient diverges, and in particular it does not equal $a=0$. Hence the converse of the Stolz theorem is false.
\end{solution}

\begin{solution}{Solution 2.2.2}\label{xs:2.2.2}
\begin{enumerate}[(i)]
\item We first record that, keeping \emph{all} the hypotheses of Theorem 2.2.1 (in particular the strict monotonicity of $\{a_n\}$), the case $a=+\infty$ does not fail. Suppose that $\{a_n\}$ is strictly increasing with $a_n\to+\infty$, that $\{b_n\}$ is a sequence, and that
\[
\lim_{n\to\infty}\frac{b_n-b_{n-1}}{a_n-a_{n-1}}=+\infty .
\]
Let $M>0$ be fixed. By the definition of a limit equal to $+\infty$ there is an $N$ such that
\[
\frac{b_n-b_{n-1}}{a_n-a_{n-1}}>M\qquad\text{for all }n>N .
\]
Since $\{a_n\}$ is strictly increasing we have $a_n-a_{n-1}>0$, so this is equivalent to
\[
b_n-b_{n-1}>M\,(a_n-a_{n-1})\qquad(n>N).
\]
Summing these inequalities for $n=N+1,N+2,\dots,n$ and telescoping both sides gives
\[
b_n-b_N>M\,(a_n-a_N),\qquad\text{that is,}\qquad b_n>M a_n+\bigl(b_N-M a_N\bigr).
\]
Now $a_n\to+\infty$, hence $a_n>0$ for all sufficiently large $n$, and dividing the last inequality by $a_n$ yields
\[
\frac{b_n}{a_n}>M+\frac{b_N-M a_N}{a_n}\longrightarrow M\qquad(n\to\infty),
\]
because $b_N-Ma_N$ is a constant. As $M>0$ was arbitrary, this says precisely that $\lim_{n\to\infty}b_n/a_n=+\infty$ (and the same argument with $-M$ in place of $M$ gives the case $a=-\infty$). Therefore, under the hypotheses of Theorem 2.2.1, the conclusion does hold in the infinite case as well; the failure asked for in this exercise must come from a genuinely weaker hypothesis, namely from dropping the strict monotonicity of the denominator.

\item The genuine failure occurs once the strict monotonicity of the denominator is weakened to $a_n\to+\infty$. Define
\[
a_n=n+(-1)^n,\qquad b_n=3\sum_{k=1}^{n}(-1)^k k,\qquad n\ge1 .
\]
We verify all the relevant properties.

\emph{The sequence $\{a_n\}$ tends to $+\infty$ but is not monotone.} Since $a_n\ge n-1$ we have $a_n\to+\infty$. On the other hand
\[
a_1=0,\quad a_2=3,\quad a_3=2,\quad a_4=5,\quad a_5=4,\ \dots,
\]
and $a_2=3>2=a_3$, so $\{a_n\}$ is not monotone; the hypothesis of Theorem 2.2.1 is indeed violated.

\emph{A closed form for $b_n$.} Pairing consecutive terms, for $n=2m$ we get
\[
\sum_{k=1}^{2m}(-1)^k k=(-1+2)+(-3+4)+\cdots+\bigl(-(2m-1)+2m\bigr)=m=\frac{n}{2},
\]
while for $n=2m-1$ the same pairing applied to the first $2m-2$ terms gives $m-1$, so that
\[
\sum_{k=1}^{2m-1}(-1)^k k=(m-1)-(2m-1)=-m=-\frac{n+1}{2}.
\]
Thus for every $n\ge1$
\[
b_n=3(-1)^n\left\lceil\frac{n}{2}\right\rceil,\qquad\text{so}\qquad |b_n|=3\left\lceil\frac{n}{2}\right\rceil ,
\]
and since $\frac{n}{2}\le\left\lceil\frac{n}{2}\right\rceil\le\frac{n+1}{2}$ we have $|b_n|\sim\frac{3n}{2}$ as $n\to\infty$.

\emph{The increment quotient tends to $+\infty$.} For $n\ge2$,
\[
a_n-a_{n-1}=\bigl(n+(-1)^n\bigr)-\bigl(n-1+(-1)^{n-1}\bigr)=1+2(-1)^n,
\]
which equals $3$ for even $n$ and $-1$ for odd $n$, and
\[
b_n-b_{n-1}=3(-1)^n n .
\]
Hence
\[
\frac{b_n-b_{n-1}}{a_n-a_{n-1}}=
\begin{cases}
\dfrac{3n}{3}=n, & n\ \text{even},\\[2ex]
\dfrac{-3n}{-1}=3n, & n\ \text{odd},
\end{cases}
\]
so this quotient equals $n$ at even indices and $3n$ at odd indices. In particular it is at least $n$ for every $n\ge2$, whence
\[
\lim_{n\to\infty}\frac{b_n-b_{n-1}}{a_n-a_{n-1}}=+\infty .
\]

\emph{The quotients $b_n/a_n$ have no limit.} Since $a_n>0$ for $n\ge2$, consider $n\ge2$. For even $n=2m$ we have $b_n=3m>0$ and $a_n=2m+1$, so
\[
\frac{b_n}{a_n}=\frac{3m}{2m+1}\longrightarrow\frac{3}{2},
\]
while for odd $n=2m+1$ we have $b_n=-3(m+1)<0$ and $a_n=2m$, so
\[
\frac{b_n}{a_n}=\frac{-3(m+1)}{2m}\longrightarrow-\frac{3}{2}.
\]
Thus $|b_n|/a_n\to\frac32$, but $b_n$ changes sign with $n$: along even indices the quotients tend to $+\frac32$ and along odd indices to $-\frac32$. Since these two subsequential limits differ, the sequence $\{b_n/a_n\}$ has no limit; in particular it does not tend to $+\infty$, being negative for all sufficiently large odd $n$.

We conclude that with $a_n=n+(-1)^n\to+\infty$ (so that the strict monotonicity required by Theorem 2.2.1 is dropped) and with $b_n=3\sum_{k=1}^{n}(-1)^k k$, the increment quotient tends to $+\infty$ whereas $b_n/a_n$ has no limit. Hence the conclusion of the Stolz theorem can fail in the infinite case as soon as the strict monotonicity of the denominator is weakened to $a_n\to+\infty$.
\end{enumerate}
\end{solution}

\begin{solution}{Solution 2.2.3}\label{xs:2.2.3}
We use the Stolz theorem (Theorem 2.2.1): if $\{a_n\}$ is strictly increasing with $a_n\to+\infty$, if $\{b_n\}$ is an arbitrary sequence, and if
\[
\lim_{n\to\infty}\frac{b_n-b_{n-1}}{a_n-a_{n-1}}=a,
\]
where $a$ may be a finite number, $+\infty$ or $-\infty$, then $\lim_{n\to\infty}b_n/a_n=a$. In each item we exhibit the sequences $a_n,b_n$, check the hypotheses and compute the increment quotient.

\begin{enumerate}[(1)]
\item Take $b_n=\sum_{k=1}^{n}\frac{1}{\sqrt{k}}$ and $a_n=\ln\sqrt{n}=\frac12\ln n$. The sequence $\{a_n\}$ is strictly increasing and $a_n\to+\infty$. For $n\ge2$,
\[
\frac{b_n-b_{n-1}}{a_n-a_{n-1}}
=\frac{\dfrac{1}{\sqrt{n}}}{\dfrac12\bigl(\ln n-\ln(n-1)\bigr)}
=\frac{2}{\sqrt{n}\,\ln\Bigl(1+\dfrac{1}{n-1}\Bigr)} .
\]
Since $\ln(1+t)\le t$ for every $t>0$, with $t=\frac{1}{n-1}$ we obtain
\[
\frac{2}{\sqrt{n}\,\ln\Bigl(1+\dfrac{1}{n-1}\Bigr)}\ \ge\ \frac{2(n-1)}{\sqrt{n}}\longrightarrow+\infty ,
\]
so the increment quotient tends to $+\infty$. By the Stolz theorem,
\[
\lim_{n\to\infty}\frac{1+\frac{1}{\sqrt{2}}+\cdots+\frac{1}{\sqrt{n}}}{\ln\sqrt{n}}=+\infty .
\]

\item Take $b_n=\sum_{k=1}^{n}\sqrt{k}$ and $a_n=n\sqrt{n}=n^{3/2}$. Again $\{a_n\}$ is strictly increasing with $a_n\to+\infty$. For $n\ge2$,
\[
\frac{b_n-b_{n-1}}{a_n-a_{n-1}}=\frac{\sqrt{n}}{n^{3/2}-(n-1)^{3/2}} .
\]
Writing $n^{3/2}-(n-1)^{3/2}=\bigl(\sqrt{n}-\sqrt{n-1}\bigr)\bigl(n+\sqrt{n(n-1)}+n-1\bigr)$ and using
$\sqrt{n}-\sqrt{n-1}=\dfrac{1}{\sqrt{n}+\sqrt{n-1}}$, we get
\[
\frac{b_n-b_{n-1}}{a_n-a_{n-1}}
=\frac{\sqrt{n}\,\bigl(\sqrt{n}+\sqrt{n-1}\bigr)}{2n-1+\sqrt{n(n-1)}}
=\frac{1+\sqrt{1-\frac1n}}{2-\frac1n+\sqrt{1-\frac1n}}\longrightarrow\frac{1+1}{2+1}=\frac23,
\]
the second equality resulting from dividing numerator and denominator by $n$. By the Stolz theorem,
\[
\lim_{n\to\infty}\frac{1+\sqrt{2}+\cdots+\sqrt{n}}{n\sqrt{n}}=\frac23 .
\]

\item Take $b_n=\sum_{k=1}^{n}\sqrt[k]{k}$ and $a_n=n$. Here $\{a_n\}$ is strictly increasing with $a_n\to+\infty$, and
\[
\frac{b_n-b_{n-1}}{a_n-a_{n-1}}=\frac{\sqrt[n]{n}}{1}=\sqrt[n]{n}=e^{\frac{\ln n}{n}} .
\]
Since $\frac{\ln n}{n}\to0$ and the exponential function is continuous, $\sqrt[n]{n}\to e^{0}=1$. By the Stolz theorem,
\[
\lim_{n\to\infty}\frac{1+\sqrt{2}+\sqrt[3]{3}+\cdots+\sqrt[n]{n}}{n}=1 .
\]

\item Take $b_n=\sum_{k=1}^{n}(2k-1)^{2}$ and $a_n=n^{3}$. The sequence $\{a_n\}$ is strictly increasing with $a_n\to+\infty$, and for $n\ge2$, since $n^{3}-(n-1)^{3}=3n^{2}-3n+1$,
\[
\frac{b_n-b_{n-1}}{a_n-a_{n-1}}
=\frac{(2n-1)^{2}}{n^{3}-(n-1)^{3}}
=\frac{4n^{2}-4n+1}{3n^{2}-3n+1}\longrightarrow\frac43 .
\]
By the Stolz theorem,
\[
\lim_{n\to\infty}\frac{1^{2}+3^{2}+\cdots+(2n-1)^{2}}{n^{3}}=\frac43 .
\]

\item Let $B_n=\sum_{k=1}^{n}k\,a_k$ and $A_n=1+2+\cdots+n=\frac{n(n+1)}{2}$, where $\{a_n\}$ is the given sequence with $\lim_{n\to\infty}a_n=a$ (a finite real number). The sequence $\{A_n\}$ is strictly increasing and $A_n\to+\infty$. Moreover
\[
B_n-B_{n-1}=n a_n,\qquad
A_n-A_{n-1}=\frac{n(n+1)}{2}-\frac{(n-1)n}{2}=n ,
\]
hence
\[
\frac{B_n-B_{n-1}}{A_n-A_{n-1}}=\frac{n a_n}{n}=a_n\longrightarrow a .
\]
By the Stolz theorem,
\[
\lim_{n\to\infty}\frac{a_1+2a_2+\cdots+n a_n}{1+2+\cdots+n}=\lim_{n\to\infty}a_n=a .
\]
(The same computation applies verbatim when $a=+\infty$ or $a=-\infty$.)
\end{enumerate}
\end{solution}

\begin{solution}{Solution 2.2.4}\label{xs:2.2.4}
Let
\[
y_n=\ln\Bigl((n!)^{\frac{1}{n^{2}}}\Bigr)=\frac{\ln n!}{n^{2}}=\frac{1}{n^{2}}\sum_{k=1}^{n}\ln k .
\]
We show that $y_n\to0$; then, since the exponential function is continuous,
\[
(n!)^{\frac{1}{n^{2}}}=e^{y_n}\longrightarrow e^{0}=1 .
\]

\emph{Upper bound.} Since $\ln k\le\ln n$ for $1\le k\le n$,
\[
0\le y_n\le\frac{n\ln n}{n^{2}}=\frac{\ln n}{n}\longrightarrow0 ,
\]
because $\frac{\ln n}{n}\to0$ as $n\to\infty$.

\emph{Lower bound.} Let $n\ge4$. Since $\ln k\ge0$ for $k\ge1$ and $\ln$ is increasing, discarding the terms of the sum with $k<\left\lceil n/2\right\rceil$ and bounding the remaining ones from below by $\ln\frac n2$ gives
\[
\sum_{k=1}^{n}\ln k\ \ge\ \Bigl(n-\left\lceil\frac n2\right\rceil+1\Bigr)\ln\frac n2\ \ge\ \frac n2\ln\frac n2 ,
\]
the last inequality because $n-\left\lceil\frac n2\right\rceil+1\ge\frac n2$ and $\ln\frac n2>0$ for $n\ge4$. Hence
\[
y_n\ge\frac{\ln\frac n2}{2n}\longrightarrow0 .
\]

By the squeeze theorem $y_n\to0$, and therefore
\[
\lim_{n\to\infty}(n!)^{\frac{1}{n^{2}}}=1 .
\]
\end{solution}

\begin{solution}{Solution 2.2.5}\label{xs:2.2.5}
Write $C_n^k=\binom{n}{k}$ and set
\[
S_n=\sum_{k=0}^{n}\ln C_n^k,\qquad T_n=\sum_{k=0}^{n}\ln k!,
\]
so that $x_n=S_n/n^{2}$. Since $C_n^k=\dfrac{n!}{k!\,(n-k)!}$,
\[
S_n=\sum_{k=0}^{n}\Bigl(\ln n!-\ln k!-\ln(n-k)!\Bigr)=(n+1)\ln n!-2T_n,
\]
the last step because $\sum_{k=0}^{n}\ln(n-k)!=\sum_{j=0}^{n}\ln j!=T_n$. Replacing $n$ by $n-1$ gives $S_{n-1}=n\ln(n-1)!-2T_{n-1}$, hence for $n\ge2$
\[
S_n-S_{n-1}=(n+1)\ln n!-n\ln(n-1)!-2\bigl(T_n-T_{n-1}\bigr)
=(n-1)\ln n!-n\ln(n-1)! ,
\]
because $T_n-T_{n-1}=\ln n!$. Using $\ln n!=\ln(n-1)!+\ln n$, we arrive at the identity
\[
S_n-S_{n-1}=(n-1)\ln n-\ln(n-1)! .
\]

We now estimate $\ln(n-1)!=\sum_{j=1}^{n-1}\ln j$ by comparing this sum with the integral of the increasing function $\ln x$. For $n\ge3$,
\[
\int_{1}^{n-1}\ln x\,dx\ \le\ \sum_{j=1}^{n-1}\ln j\ \le\ \int_{1}^{n}\ln x\,dx ,
\]
because $\ln j\ge\int_{j-1}^{j}\ln x\,dx$ for $j\ge2$, while $\ln j\le\int_{j}^{j+1}\ln x\,dx$ for $j\ge1$. Since $\int_{1}^{N}\ln x\,dx=N\ln N-N+1$, this reads
\[
(n-1)\ln(n-1)-n+2\ \le\ \ln(n-1)!\ \le\ n\ln n-n+1 .
\]
Consequently, using $\ln\dfrac{n}{n-1}\le\dfrac{1}{n-1}$,
\[
(n-1)\ln n-\ln(n-1)!\ \le\ (n-1)\ln\frac{n}{n-1}+n-2\ \le\ 1+n-2=n-1 ,
\]
while on the other side
\[
(n-1)\ln n-\ln(n-1)!\ \ge\ (n-1)\ln n-\bigl(n\ln n-n+1\bigr)=n-1-\ln n .
\]
Dividing by $n^{2}-(n-1)^{2}=2n-1$ we get, for every $n\ge3$,
\[
\frac{n-1-\ln n}{2n-1}\ \le\ \frac{S_n-S_{n-1}}{n^{2}-(n-1)^{2}}\ \le\ \frac{n-1}{2n-1},
\]
and both outer expressions tend to $\frac12$; by the squeeze theorem,
\[
\lim_{n\to\infty}\frac{S_n-S_{n-1}}{n^{2}-(n-1)^{2}}=\frac12 .
\]
The sequence $\{n^{2}\}$ is strictly increasing and tends to $+\infty$, so the Stolz theorem (Theorem 2.2.1), applied with $a_n=n^{2}$ and $b_n=S_n$, yields
\[
\lim_{n\to\infty}x_n=\lim_{n\to\infty}\frac{S_n}{n^{2}}=\frac12 .
\]
(Stirling's formula gives the same value: $\ln n!=n\ln n-n+O(\ln n)$ together with $\sum_{k=1}^{n}k\ln k=\frac12n^{2}\ln n-\frac14n^{2}+O(n\ln n)$ leads to $S_n=\frac12n^{2}+O(n\ln n)$. A numerical check confirms the limit as well: $x_{200}\approx0.4871$, $x_{1000}\approx0.4966$, $x_{5000}\approx0.4992$.)
\end{solution}

\begin{solution}{Solution 2.2.6}\label{xs:2.2.6}
Put
\[
L=\lim_{n\to\infty}\frac{A_1+A_2+\cdots+A_n}{n},\qquad S_n=A_1+A_2+\cdots+A_n ,
\]
so that $S_n/n\to L$, and put $D_j=A_j-A_{j-1}$ for $j\ge2$. The hypothesis says $jD_j=j(A_j-A_{j-1})\to0$ as $j\to\infty$.

\emph{An identity.} For every $n\ge2$,
\[
nA_n-S_n=\sum_{j=2}^{n}(j-1)\bigl(A_j-A_{j-1}\bigr)=\sum_{j=2}^{n}(j-1)D_j .
\]
Indeed, expanding the right-hand side and re-indexing the second sum,
\[
\sum_{j=2}^{n}(j-1)A_j-\sum_{j=2}^{n}(j-1)A_{j-1}
=\sum_{j=2}^{n}(j-1)A_j-\sum_{i=1}^{n-1}iA_i
=(n-1)A_n-\sum_{i=1}^{n-1}A_i ,
\]
and since $\sum_{i=1}^{n-1}A_i=S_n-A_n$, this equals
\[
(n-1)A_n-S_n+A_n=nA_n-S_n .
\]
Dividing the identity by $n$ gives
\[
A_n=\frac{S_n}{n}+\frac{1}{n}\sum_{j=2}^{n}(j-1)D_j . \tag{$\ast$}
\]

\emph{The second term tends to $0$.} By hypothesis $jD_j\to0$, and then also
\[
D_j=\frac1j\,(jD_j)\longrightarrow0 .
\]
Hence the sequence $u_j=(j-1)D_j$ satisfies $u_j=jD_j-D_j\to0$. By the Ces\`aro mean theorem (Example 2.1.4: if $u_j\to u$ then $\frac{u_1+\cdots+u_j}{j}\to u$), applied to the sequence $u_1=0$, $u_j=(j-1)D_j$ for $j\ge2$, whose limit is $0$, we obtain
\[
\frac1n\sum_{j=2}^{n}(j-1)D_j=\frac1n\sum_{j=1}^{n}u_j\longrightarrow0 .
\]

\emph{Conclusion.} Letting $n\to\infty$ in $(\ast)$ and using $S_n/n\to L$ together with the preceding result,
\[
\lim_{n\to\infty}A_n=L+0=L=\lim_{n\to\infty}\frac{A_1+A_2+\cdots+A_n}{n},
\]
which is the required equality. (If the Ces\`aro limit is $+\infty$ or $-\infty$, then $(\ast)$ shows that $A_n$ has the same infinite limit, since its second term still tends to $0$.)
\end{solution}


\srcline{Exercise set 2.3, problems 1--21}

\begin{solution}{Solution 2.3.1}\label{xs:2.3.1}
\begin{enumerate}[(1)]
  \item We first record the bound $0<x_n<2$ for all $n$, by induction: $x_1=\sqrt2<2$, and if $0<x_n<2$ then
  \[
  x_{n+1}=\sqrt{2x_n}<\sqrt{2\cdot 2}=2 .
  \]
  Moreover, for every $n$,
  \[
  \frac{x_{n+1}}{x_n}=\sqrt{\frac{2}{x_n}}>1,
  \]
  because $0<x_n<2$; hence $\{x_n\}$ is strictly increasing. Being increasing and bounded above by $2$, it converges (monotone convergence theorem). Let $L=\lim_{n\to\infty}x_n$. Since $x_n\ge x_1=\sqrt2$ for all $n$, we have $L\ge\sqrt2>0$. Squaring the recursion gives $x_{n+1}^2=2x_n$; letting $n\to\infty$ and using the limit rules for products yields $L^2=2L$, i.e.\ $L(L-2)=0$. As $L>0$, $L=2$. Thus $\lim_{n\to\infty}x_n=2$.

  \item Put $\alpha=\dfrac{1+\sqrt{1+4c}}{2}>0$ and $\beta=\dfrac{1-\sqrt{1+4c}}{2}<0$; these are the two roots of $t^2-t-c=0$, so in particular $\alpha^2=\alpha+c$. We prove by induction that $0<x_n<\alpha$: indeed $x_1=\sqrt c<\alpha$ is equivalent to $c<\alpha^2=\alpha+c$, which is true since $\alpha>0$; and if $0<x_n<\alpha$, then
  \[
  x_{n+1}=\sqrt{c+x_n}<\sqrt{c+\alpha}=\sqrt{\alpha^2}=\alpha .
  \]
  Next, for $0<x_n<\alpha$,
  \[
  x_{n+1}>x_n\iff c+x_n>x_n^2\iff x_n^2-x_n-c<0\iff (x_n-\alpha)(x_n-\beta)<0,
  \]
  and the last inequality holds because $\beta<0<x_n<\alpha$. Hence $\{x_n\}$ is strictly increasing and bounded above by $\alpha$, so it converges; let $L=\lim x_n$. Then $L=\sqrt{c+L}$, so $L^2=L+c$ with $L\ge x_1=\sqrt c>0$. Therefore $L$ is the positive root, i.e.
  \[
  \lim_{n\to\infty}x_n=\frac{1+\sqrt{1+4c}}{2}.
  \]

  \item Since $x_1>1$, we claim $x_n>1$ for every $n$: if $x_n>1$, then
  \[
  x_{n+1}-1=\frac{3x_n+1-x_n-3}{x_n+3}=\frac{2(x_n-1)}{x_n+3}>0 .
  \]
  Hence $\{x_n\}$ is bounded below by $1$. It is also decreasing, since for $x_n>1$
  \[
  x_{n+1}-x_n=\frac{3x_n+1-x_n^2-3x_n}{x_n+3}=\frac{1-x_n^2}{x_n+3}<0 .
  \]
  A decreasing sequence bounded below converges. Let $L=\lim x_n$; then $L\ge 1$ and, passing to the limit in the recursion (the function $t\mapsto(3t+1)/(t+3)$ is continuous at $L$, where the denominator is positive),
  \[
  L=\frac{3L+1}{L+3}\iff L^2+3L=3L+1\iff L^2=1 .
  \]
  Since $L\ge1$ we get $L=1$, i.e.\ $\lim_{n\to\infty}x_n=1$.
\end{enumerate}
\end{solution}

\begin{solution}{Solution 2.3.2}\label{xs:2.3.2}
We prove first that
\[
0<a_n<b_n\qquad\text{for all } n\in\mathbb N_+,
\]
by induction. This holds for $n=1$ by hypothesis. If $0<a_n<b_n$, then $a_n\ne b_n$ and the arithmetic--geometric mean inequality gives
\[
a_{n+1}=\sqrt{a_nb_n}<\frac{a_n+b_n}{2}=b_{n+1},
\]
and both numbers are positive. So the claim follows.

Now $a_{n+1}=\sqrt{a_nb_n}>a_n$ because $b_n>a_n$, so $\{a_n\}$ is strictly increasing; and $b_{n+1}=\frac{a_n+b_n}{2}<b_n$, so $\{b_n\}$ is strictly decreasing. Moreover
\[
a_1\le a_n<b_n\le b_1\qquad(n\in\mathbb N_+),
\]
so $\{a_n\}$ is increasing and bounded above by $b_1$, while $\{b_n\}$ is decreasing and bounded below by $a_1$. By the monotone convergence theorem both converge; write $a=\lim a_n$ and $b=\lim b_n$. Then
\[
b=\lim b_{n+1}=\lim\frac{a_n+b_n}{2}=\frac{a+b}{2},
\]
whence $2b=a+b$, i.e.\ $a=b$. Thus the two sequences converge to the same limit. (The common value is the arithmetic--geometric mean of $a_1$ and $b_1$; alternatively, from
\[
b_{n+1}-a_{n+1}=\frac{(\sqrt{b_n}-\sqrt{a_n})^2}{2}\le\frac{b_n-a_n}{2}
\]
one gets $0<b_n-a_n\le 2^{1-n}(b_1-a_1)\to 0$, which again shows $a=b$.)
\end{solution}

\begin{solution}{Solution 2.3.3}\label{xs:2.3.3}
For all $n$ we have $0<x_n<1$ and $(1-x_n)x_{n+1}>\frac14$. By the arithmetic--geometric mean inequality,
\[
\frac14<(1-x_n)x_{n+1}\le\left(\frac{(1-x_n)+x_{n+1}}{2}\right)^2 .
\]
Taking square roots (both sides are positive) gives
\[
\frac12<\frac{1-x_n+x_{n+1}}{2},
\]
that is,
\[
x_{n+1}>x_n\qquad(n=1,2,3,\dots).
\]
Hence $\{x_n\}$ is strictly increasing. It is bounded above by $1$ by hypothesis, so it converges; let $L=\lim_{n\to\infty}x_n$. Since $x_n>x_1>0$ and $x_n<1$, we have $0<L\le1$.

Because the limit of a product is the product of the limits,
\[
(1-x_n)x_{n+1}\longrightarrow (1-L)L .
\]
Every term of the sequence $\{(1-x_n)x_{n+1}\}$ is $>\frac14$, and a limit of numbers each $>\frac14$ is $\ge\frac14$; hence $L(1-L)\ge\frac14$. On the other hand
\[
L(1-L)=\frac14-\left(L-\frac12\right)^2\le\frac14 .
\]
Therefore $L(1-L)=\frac14$, which forces $L-\frac12=0$, i.e.\
\[
\lim_{n\to\infty}x_n=\frac12 .
\]
\end{solution}

\begin{solution}{Solution 2.3.4}\label{xs:2.3.4}
Suppose first that $\{x_n\}$ is increasing (the decreasing case is obtained by applying the increasing case to $\{-x_n\}$, which is increasing and has the convergent subsequence $\{-x_{n_k}\}$).

Let $\{x_{n_k}\}$ be a convergent subsequence with $x_{n_k}\to a$. We claim that
\[
x_n\le a\qquad\text{for every } n .
\]
Assume on the contrary that $x_m>a$ for some $m$. Put $\delta=x_m-a>0$. Since $n_k\to+\infty$, there is $k$ with $n_k>m$; as $\{x_n\}$ is increasing, $x_{n_k}\ge x_m=a+\delta$ for all such $k$, and there are infinitely many of them. But $x_{n_k}\to a$, so for $\varepsilon=\delta/2$ we must have $x_{n_k}<a+\delta/2$ for all large $k$ --- a contradiction. Hence $x_n\le a$ for all $n$.

Thus the increasing sequence $\{x_n\}$ is bounded above by $a$, and the monotone convergence theorem shows that it converges; let $L=\lim x_n$. Then $L=\sup_n x_n\le a$. On the other hand $x_{n_k}\le L$ for every $k$, so letting $k\to\infty$ gives $a\le L$. Hence $L=a$, and $\{x_n\}$ converges (to the same limit as the subsequence).
\end{solution}

\begin{solution}{Solution 2.3.5}\label{xs:2.3.5}
Write $f(t)=t(1-t)$ for $t\in(0,1)$. Since $0<t<1$ implies $0<f(t)<1$ (indeed $f(t)>0$ and $f(t)\le\frac14<1$), we get by induction
\[
0<x_n<1\qquad(n\in\mathbb N_+).
\]
Moreover
\[
x_{n+1}=x_n-x_n^2<x_n ,
\]
so $\{x_n\}$ is strictly decreasing and bounded below by $0$. By the monotone convergence theorem it converges; if $L=\lim x_n$, then $L=L(1-L)$ (continuity of $t\mapsto t(1-t)$), so $L^2=0$ and
\[
\lim_{n\to\infty}x_n=0 .
\]

Now we determine the limit of $\{nx_n\}$. For every $n$,
\[
\frac{1}{x_{n+1}}=\frac{1}{x_n(1-x_n)}=\frac{1}{x_n}+\frac{1}{1-x_n},
\]
so, summing from $k=1$ to $k=n-1$,
\[
\frac{1}{x_n}=\frac{1}{x_1}+\sum_{k=1}^{n-1}\frac{1}{1-x_k}.
\]
Since $x_k\to0$, we have $\dfrac{1}{1-x_k}\to1$. By the Ces\`aro mean theorem (if $c_k\to c$ then the arithmetic means $\frac{1}{n}\sum_{k=1}^{n}c_k\to c$), applied to $c_k=\dfrac{1}{1-x_k}$,
\[
\frac{1}{n}\sum_{k=1}^{n-1}\frac{1}{1-x_k}=\frac{n-1}{n}\cdot\frac{1}{n-1}\sum_{k=1}^{n-1}\frac{1}{1-x_k}\longrightarrow 1 .
\]
Therefore
\[
\frac{1}{nx_n}=\frac{1}{n}\cdot\frac{1}{x_n}=\frac{1}{nx_1}+\frac{1}{n}\sum_{k=1}^{n-1}\frac{1}{1-x_k}\longrightarrow 0+1=1,
\]
and consequently
\[
\lim_{n\to\infty}nx_n=1 .
\]
\end{solution}

\begin{solution}{Solution 2.3.6}\label{xs:2.3.6}
Since $l>1$, choose a number $q$ with $1<q<l$ (for instance $q=(1+l)/2$). From $\lim_{n\to\infty}\frac{a_n}{a_{n+1}}=l>q$ we get an index $N$ such that
\[
\frac{a_n}{a_{n+1}}>q\qquad\text{for all } n\ge N .
\]
As $a_{n+1}>0$, this is equivalent to
\[
0<a_{n+1}<\frac{a_n}{q}\qquad(n\ge N).
\]
Iterating this inequality, for every $n>N$,
\[
0<a_n<\frac{a_N}{q^{\,n-N}}=a_Nq^{N}q^{-n}.
\]
Since $q>1$, we have $q^{-n}\to0$ (for $q=1+r$ with $r>0$, Bernoulli's inequality gives $q^n=(1+r)^n\ge 1+nr\to\infty$, hence $q^{-n}\to0$). Therefore the right-hand side tends to $0$, and the squeeze theorem gives
\[
\lim_{n\to\infty}a_n=0 .
\]
\end{solution}

\begin{solution}{Solution 2.3.7}\label{xs:2.3.7}
For every integer $k\ge1$ and every $t\in(k,k+1)$ we have $\dfrac{1}{k+1}<\dfrac1t<\dfrac1k$. Integrating this inequality over $(k,k+1)$,
\[
\frac{1}{k+1}=\int_k^{k+1}\frac{dt}{k+1}<\int_k^{k+1}\frac{dt}{t}=\ln\frac{k+1}{k}<\int_k^{k+1}\frac{dt}{k}=\frac1k .
\]
So for all $k\ge1$,
\[
\frac{1}{k+1}<\ln\!\left(1+\frac1k\right)<\frac1k. \tag{$\ast$}
\]

\emph{Left inequality.} Summing $(\ast)$ for $k=1,\dots,n-1$ and using the telescoping sum $\sum_{k=1}^{n-1}\ln\frac{k+1}{k}=\ln n$,
\[
\frac12+\frac13+\dots+\frac1n=\sum_{k=1}^{n-1}\frac{1}{k+1}<\sum_{k=1}^{n-1}\ln\frac{k+1}{k}=\ln n<\ln(n+1).
\]
(In fact the same argument with $k=1,\dots,n$ gives the sharper $\frac12+\dots+\frac{1}{n+1}<\ln(n+1)$.)

\emph{Right inequality.} Summing $(\ast)$ for $k=1,\dots,n$ and using $\sum_{k=1}^{n}\ln\frac{k+1}{k}=\ln(n+1)$,
\[
\ln(n+1)=\sum_{k=1}^{n}\ln\frac{k+1}{k}<\sum_{k=1}^{n}\frac1k=1+\frac12+\dots+\frac1n .
\]
Combining the two estimates,
\[
\frac12+\frac13+\dots+\frac1n<\ln(n+1)<1+\frac12+\dots+\frac1n .
\]
\end{solution}

\begin{solution}{Solution 2.3.8}\label{xs:2.3.8}
Let $x_n=H_n-\ln n$, where $H_n=1+\frac12+\dots+\frac1n$.

\emph{$\{x_n\}$ is decreasing.} For every $n\ge1$,
\[
x_n-x_{n+1}=\bigl(H_n-\ln n\bigr)-\bigl(H_{n+1}-\ln(n+1)\bigr)
=\ln\frac{n+1}{n}-\frac{1}{n+1}>0,
\]
where the strict inequality is exactly $\ln\frac{n+1}{n}>\frac{1}{n+1}$, i.e.\ the left inequality $(\ast)$ of Solution 2.3.7 with $k=n$.

\emph{$\{x_n\}$ is bounded below.} By the right inequality of Exercise 2.3.7 applied with $n$ replaced by $n-1$ (for $n\ge2$),
\[
\ln n<1+\frac12+\dots+\frac{1}{n-1}<H_n,
\]
so $x_n=H_n-\ln n>0$ for $n\ge2$; and $x_1=1-\ln1=1>0$. Hence $0<x_n\le1$ for all $n$.

A decreasing sequence bounded below converges (monotone convergence theorem). Therefore $\{x_n\}$ converges; its limit is Euler's constant
\[
\gamma=\lim_{n\to\infty}\left(1+\frac12+\dots+\frac1n-\ln n\right).
\]
\end{solution}

\begin{solution}{Solution 2.3.9}\label{xs:2.3.9}
\textbf{The limit is $0$.} We use the standard series for $e$ (obtained from $e=\lim_{n\to\infty}(1+1/n)^n$ by expanding the binomial and passing to the limit):
\[
e=\sum_{j=0}^{\infty}\frac{1}{j!},\qquad 2<e<3 .
\]
Fix $n\in\mathbb N_+$. For $0\le j\le n$ the number $n!/j!$ is a positive integer, so
\[
A_n:=\sum_{j=0}^{n}\frac{n!}{j!}\in\mathbb N .
\]
Splitting the series for $n!e$ at $j=n$,
\[
n!\,e=\sum_{j=0}^{n}\frac{n!}{j!}+\sum_{j=n+1}^{\infty}\frac{n!}{j!}=A_n+r_n,
\qquad
r_n:=\sum_{j=n+1}^{\infty}\frac{n!}{j!}>0 .
\]
For $j\ge n+1$,
\[
\frac{n!}{j!}=\frac{1}{(n+1)(n+2)\cdots j}\le\frac{1}{(n+1)^{\,j-n}},
\]
hence
\[
0<r_n\le\sum_{j=n+1}^{\infty}\frac{1}{(n+1)^{\,j-n}}
=\sum_{i=1}^{\infty}\frac{1}{(n+1)^{i}}
=\frac{1}{n+1}\cdot\frac{1}{1-\frac{1}{n+1}}=\frac1n .
\]
In particular $0<r_n\le 1/n<1$ for every $n\ge 2$, while for $n=1$ we have $0<r_1=e-2<1$ because $2<e<3$. Thus in every case $0<r_n<1$, so $A_n$ is precisely the integer part of $n!\,e$; equivalently $[n!\,e]=A_n$ and
\[
n!\,e-[n!\,e]=r_n .
\]
Since $0<r_n\le \dfrac1n$ for all $n$ and $\dfrac1n\to 0$, the squeeze theorem gives
\[
\lim_{n\to\infty}\bigl(n!\,e-[n!\,e]\bigr)=\lim_{n\to\infty}r_n=0 .
\]
\end{solution}

\begin{solution}{Solution 2.3.10}\label{xs:2.3.10}
Denote
\[
S_n=\frac{1}{n+1}+\frac{1}{n+2}+\dots+\frac{1}{n+n}=\sum_{k=n+1}^{2n}\frac1k .
\]
We use the elementary estimate, valid for every $k\ge1$ (see Solution 2.3.7),
\[
\frac{1}{k+1}<\ln\frac{k+1}{k}<\frac1k .
\]

\emph{Upper bound.} Summing the left inequality over $k=n,n+1,\dots,2n-1$ and telescoping,
\[
\sum_{k=n}^{2n-1}\frac{1}{k+1}=\sum_{j=n+1}^{2n}\frac1j=S_n<\sum_{k=n}^{2n-1}\ln\frac{k+1}{k}=\ln\frac{2n}{n}=\ln 2 .
\]

\emph{Lower bound.} Summing the right inequality over the same range,
\[
\ln 2=\sum_{k=n}^{2n-1}\ln\frac{k+1}{k}<\sum_{k=n}^{2n-1}\frac1k=\frac1n+\sum_{k=n+1}^{2n-1}\frac1k
=\frac1n+S_n-\frac{1}{2n}=S_n+\frac{1}{2n},
\]
hence $S_n>\ln 2-\dfrac{1}{2n}$.

Therefore
\[
\ln 2-\frac{1}{2n}<S_n<\ln 2\qquad(n\ge1),
\]
and the squeeze theorem (since $1/(2n)\to0$) gives
\[
\lim_{n\to\infty}\left(\frac{1}{n+1}+\frac{1}{n+2}+\dots+\frac{1}{n+n}\right)=\ln 2 .
\]
\end{solution}

\begin{solution}{Solution 2.3.11}\label{xs:2.3.11}
Since $1+\dfrac{1}{2^{n+1}}>1$,
\[
x_{n+1}=x_n\left(1+\frac{1}{2^{n+1}}\right)>x_n ,
\]
so $\{x_n\}$ is strictly increasing. It remains to show that it is bounded above.

We use the following elementary fact: if $a_1,\dots,a_m\ge0$ and $s=a_1+\dots+a_m<1$, then
\[
\prod_{k=1}^{m}(1+a_k)\le\frac{1}{1-s}.
\]
(Proof by induction on $m$: for $m=1$ it is $1+a_1\le\frac{1}{1-a_1}$, i.e.\ $(1+a_1)(1-a_1)=1-a_1^2\le1$. Assume it for $m-1$ with partial sum $s'=a_1+\dots+a_{m-1}$; then, using $s=s'+a_m$,
\[
\prod_{k=1}^{m}(1+a_k)\le\frac{1+a_m}{1-s'}
\le\frac{1}{1-s},
\]
the last step being equivalent to $(1+a_m)(1-s)\le 1-s'$, i.e.\ to $-a_ms\le0$, which is true.)

Apply this with $a_k=2^{-k}$ for $k=2,\dots,n$. Since
\[
\sum_{k=2}^{n}\frac{1}{2^k}<\sum_{k=2}^{\infty}\frac{1}{2^k}=\frac12<1,
\]
we obtain
\[
x_n=\frac32\prod_{k=2}^{n}\left(1+\frac{1}{2^k}\right)
\le\frac32\cdot\frac{1}{1-\sum_{k=2}^{n}2^{-k}}
=\frac32\cdot\frac{1}{\frac12+2^{-n}}\le\frac32\cdot 2=3 .
\]
So $\{x_n\}$ is increasing and bounded above by $3$; by the monotone convergence theorem $\displaystyle\lim_{n\to\infty}x_n$ exists. (Its value is the infinite product $\prod_{k=1}^{\infty}\left(1+2^{-k}\right)\approx2.3842$.)
\end{solution}

\begin{solution}{Solution 2.3.12}\label{xs:2.3.12}
The Cauchy criterion for sequences, as established in this chapter, reads:
\[
\{x_n\}\ \text{converges}\iff
\forall\varepsilon>0\ \exists N\in\mathbb N\ \forall n>N\ \forall p\in\mathbb N_+:\ |x_{n+p}-x_n|<\varepsilon ,
\]
i.e.\ the index $N$ must be independent of $p$. We test each statement against this criterion.

\textbf{(1) Not a necessary and sufficient condition (it is necessary but not sufficient).}
The quantifier order ``$\forall p\ \exists N$'' allows $N$ to depend on $p$, so (1) says precisely
\[
\text{for every fixed }p\in\mathbb N_+:\quad \lim_{n\to\infty}(x_{n+p}-x_n)=0 .
\]
\emph{Necessary:} if $x_n\to a$, then for fixed $p$, $x_{n+p}\to a$ as $n\to\infty$, hence $x_{n+p}-x_n\to a-a=0$; so (1) holds for every convergent sequence.
\emph{Not sufficient:} take $x_n=\sqrt n$. This sequence diverges to $+\infty$, yet for every fixed $p\ge1$,
\[
0<\sqrt{n+p}-\sqrt n=\frac{p}{\sqrt{n+p}+\sqrt n}<\frac{p}{2\sqrt n}\longrightarrow0 ,
\]
so (1) is satisfied. Hence (1) does not characterize convergence. (The same conclusion holds for $x_n=1+\frac12+\dots+\frac1n$.)

\textbf{(2) Not a necessary and sufficient condition (it is necessary but not sufficient).}
Here $N$ and $p$ are both chosen existentially, i.e.\ a \emph{single} pair $(N,p)$ must make $|x_{n+p}-x_n|<\varepsilon$ for all $n>N$.
\emph{Necessary:} if $x_n\to a$, then $x_{n+1}-x_n\to0$, so for every $\varepsilon>0$ there is $N$ with $|x_{n+1}-x_n|<\varepsilon$ for all $n>N$; thus (2) holds with $p=1$.
\emph{Not sufficient:} if $0\in\mathbb N$ is allowed, then (2) holds for \emph{every} sequence (take $p=0$: $|x_n-x_n|=0<\varepsilon$), so it certainly does not force convergence; for example the divergent sequence $x_n=n$ satisfies (2). If instead $\mathbb N=\{1,2,\dots\}$, then $p\ge1$, and the divergent sequence $x_n=\sqrt n$ still satisfies (2) with $p=1$, because $|x_{n+1}-x_n|=\frac{1}{\sqrt{n+1}+\sqrt n}\to0$. So in either reading (2) is not sufficient.

\textbf{(3) Not a necessary and sufficient condition (it is necessary but not sufficient).}
Statement (3) is the disjunction of
\[
\text{(A)}\quad \forall n,p\in\mathbb N_+:\ |x_{n+p}-x_n|<\frac{p}{n},
\qquad
\text{(B)}\quad \forall p\in\mathbb N_+:\ \lim_{n\to\infty}(x_{n+p}-x_n)=0 .
\]
\emph{Necessary:} (B) is exactly condition (1), which was shown above to hold for every convergent sequence; hence the disjunction holds for every convergent sequence.
\emph{Not sufficient:} consider $x_n=\ln\ln(n+2)$ ($n\ge1$). It is strictly increasing and $x_n\to+\infty$, so it diverges. On the other hand, for all $n,p\in\mathbb N_+$,
\[
x_{n+p}-x_n=\ln\frac{\ln(n+p+2)}{\ln(n+2)}
=\ln\!\left(1+\frac{\ln\frac{n+p+2}{n+2}}{\ln(n+2)}\right)
\le\frac{\ln\left(1+\frac{p}{n+2}\right)}{\ln(n+2)}
\le\frac{p}{(n+2)\ln(n+2)}<\frac{p}{n},
\]
where we used $\ln(1+u)\le u$ twice and, in the last step, $\ln(n+2)\ge\ln3>1$, so $(n+2)\ln(n+2)>n$. Thus (A) holds for the divergent sequence $\{x_n\}$, and (3) is not sufficient. (The first alternative (A) is not necessary either: for $x_n=\sum_{2^j\le n}2\cdot2^{-j}$, a convergent sequence with jumps of size $2^{1-k}$ at $n=2^k$, we have $|x_{n+1}-x_n|=\frac{2}{n+1}>\frac1n$ at $n=2^k-1$, so (A) fails; the disjunction is nevertheless necessary because (B) holds.)

\textbf{(4) Not a necessary and sufficient condition (it is sufficient but not necessary).}
\emph{Sufficient:} assume $|x_{n+p}-x_n|<\dfrac{p}{n^2}$ for all $n,p\in\mathbb N_+$. Taking $p=1$ gives
\[
|x_{n+1}-x_n|<\frac{1}{n^2}\qquad(n\in\mathbb N_+).
\]
For $m>n\ge2$, the triangle inequality and the telescoping estimate $\frac{1}{k^2}\le\frac{1}{k(k-1)}=\frac{1}{k-1}-\frac1k$ yield
\[
|x_m-x_n|\le\sum_{k=n}^{m-1}|x_{k+1}-x_k|<\sum_{k=n}^{m-1}\frac{1}{k^2}
\le\sum_{k=n}^{\infty}\frac{1}{k(k-1)}=\frac{1}{n-1}.
\]
Given $\varepsilon>0$, choose $N>1+\frac1\varepsilon$; then for all $m>n>N$ we have $|x_m-x_n|<\frac{1}{n-1}<\varepsilon$. Thus $\{x_n\}$ is a Cauchy sequence, and by the Cauchy convergence criterion it converges.
\emph{Not necessary:} define
\[
x_n=\begin{cases}2^{-j}, & n=2^j\ \text{for some } j\in\mathbb N_+,\\[2pt] 0, & \text{otherwise.}\end{cases}
\]
Then $0\le x_n\to0$, so $\{x_n\}$ converges (to $0$). But for $n=3$ and $p=1$,
\[
|x_{4}-x_3|=\left|\frac14-0\right|=\frac14>\frac{1}{9}=\frac{p}{n^2},
\]
so (4) fails. Hence (4) is sufficient but not necessary for convergence.

\textbf{Conclusion.} None of (1)--(4) is a necessary and sufficient condition for convergence: (1), (2), (3) are necessary but not sufficient, while (4) is sufficient but not necessary.
\end{solution}

\begin{solution}{Solution 2.3.13}\label{xs:2.3.13}
We use the \emph{Cauchy convergence criterion}: a sequence $\{a_n\}$ of real numbers converges if and only if for every $\varepsilon > 0$ there is an $N \in \mathbb{N}_+$ such that
\[
|a_m - a_n| < \varepsilon \qquad \text{whenever } m > n \ge N .
\]
Throughout we use the elementary estimate, valid for $m > n \ge 1$,
\[
\frac{1}{k^2} \le \frac{1}{k(k-1)} = \frac{1}{k-1} - \frac{1}{k} \quad (k \ge 2),
\qquad\text{hence}\qquad
\sum_{k=n+1}^{m} \frac{1}{k^2} \le \sum_{k=n+1}^{m}\Bigl(\frac{1}{k-1}-\frac{1}{k}\Bigr) = \frac{1}{n} - \frac{1}{m} < \frac{1}{n}.
\]
\begin{enumerate}[(1)]
  \item For $m > n \ge 1$,
  \[
  |a_m - a_n| = \Bigl|\sum_{k=n+1}^{m} \frac{(-1)^{k-1}}{k^2}\Bigr| \le \sum_{k=n+1}^{m}\frac{1}{k^2} < \frac{1}{n}.
  \]
  Given $\varepsilon > 0$, choose $N \in \mathbb{N}_+$ with $N > 1/\varepsilon$. Then for all $m > n \ge N$,
  \[
  |a_m - a_n| < \frac{1}{n} \le \frac{1}{N} < \varepsilon .
  \]
  Thus $\{a_n\}$ satisfies the Cauchy criterion, and therefore converges.
  \item Since $|\sin kx| \le 1$ for every $k \in \mathbb{N}_+$ and every $x \in \mathbb{R}$, for $m > n \ge 1$,
  \[
  |a_m - a_n| = \Bigl|\sum_{k=n+1}^{m} \frac{\sin kx}{k^2}\Bigr| \le \sum_{k=n+1}^{m}\frac{|\sin kx|}{k^2} \le \sum_{k=n+1}^{m}\frac{1}{k^2} < \frac{1}{n}.
  \]
  Given $\varepsilon > 0$, take $N \in \mathbb{N}_+$ with $N > 1/\varepsilon$; then $|a_m - a_n| < \varepsilon$ for all $m > n \ge N$. So $\{a_n\}$ is a Cauchy sequence and hence converges.
  \item Put $u_k = \dfrac{\sin kx}{k(k+\sin kx)}$ for $k \ge 2$, so that $a_n = \sum_{k=2}^{n} u_k$ for $n \ge 2$ (this is the range of indices for which $a_n$ is defined). For $k \ge 2$,
  \[
  k + \sin kx \ge k - 1 > 0,
  \qquad\text{so}\qquad
  |u_k| = \frac{|\sin kx|}{k\,(k+\sin kx)} \le \frac{1}{k(k-1)} .
  \]
  Hence for $m > n \ge 2$,
  \[
  |a_m - a_n| = \Bigl|\sum_{k=n+1}^{m} u_k\Bigr| \le \sum_{k=n+1}^{m}|u_k| \le \sum_{k=n+1}^{m}\frac{1}{k(k-1)} = \frac{1}{n} - \frac{1}{m} < \frac{1}{n}.
  \]
  Given $\varepsilon > 0$, choose an integer $N \ge 2$ with $N > 1/\varepsilon$. Then $|a_m - a_n| < 1/n \le 1/N < \varepsilon$ for all $m > n \ge N$, so by the Cauchy criterion $\{a_n\}$ converges.
\end{enumerate}
\end{solution}

\begin{solution}{Solution 2.3.14}\label{xs:2.3.14}
Define, for $n \ge 2$,
\[
S_n = |a_2 - a_1| + |a_3 - a_2| + \cdots + |a_n - a_{n-1}| = \sum_{k=2}^{n}|a_k - a_{k-1}| .
\]
By hypothesis the sequence $\{S_n\}_{n \ge 2}$ is bounded, say $S_n \le M$ for all $n \ge 2$. Since every term $|a_k - a_{k-1}| \ge 0$, the sequence $\{S_n\}$ is nondecreasing; being nondecreasing and bounded above, it converges (its limit is $\sup_n S_n$). Write $L = \lim_{n\to\infty} S_n$.

We show that $\{a_n\}$ is a Cauchy sequence. For $m > n \ge 2$, the triangle inequality gives
\[
|a_m - a_n| = \Bigl|\sum_{k=n+1}^{m}\bigl(a_k - a_{k-1}\bigr)\Bigr|
\le \sum_{k=n+1}^{m}|a_k - a_{k-1}| = S_m - S_n .
\]
Let $\varepsilon > 0$. Since $S_n \to L$, there is $N \ge 2$ with $L - S_N < \varepsilon$. Then for all $m > n \ge N$, using $S_n \ge S_N$ and $S_m \le L$,
\[
|a_m - a_n| \le S_m - S_n \le L - S_n \le L - S_N < \varepsilon .
\]
Hence $\{a_n\}$ satisfies the Cauchy convergence criterion and therefore converges.
\end{solution}

\begin{solution}{Solution 2.3.15}\label{xs:2.3.15}
First observe that $b_n > 0$ for every $n \in \mathbb{N}_+$. Indeed, a strictly decreasing sequence is bounded below by its limit, and here that limit is $0$; hence $b_n \ge 0$ for all $n$, and since $b_n \to 0$ and $\{b_n\}$ is strictly decreasing, in fact $b_n > 0$ for all $n$. In particular $b_n - b_{n-1} < 0$ and $b_n - b_{n-1} \neq 0$ for $n \ge 2$.

\medskip
\noindent\textbf{Case 1: $A$ is a finite number.}
Put $u_n = a_n - A b_n$. Then $u_n \to 0 - A\cdot 0 = 0$, and for $n \ge 2$,
\[
\frac{u_n - u_{n-1}}{b_n - b_{n-1}}
= \frac{a_n - a_{n-1}}{b_n - b_{n-1}} - A \longrightarrow 0 .
\]
So it suffices to prove the following claim.

\emph{Claim.} If $c_n \to 0$, $b_1 > b_2 > b_3 > \cdots$ with $b_n \to 0$, and
$\dfrac{c_n - c_{n-1}}{b_n - b_{n-1}} \to 0$, then $\dfrac{c_n}{b_n} \to 0$.

\emph{Proof of the claim.} Let $\varepsilon > 0$. Choose $N$ such that
\[
\Bigl|\frac{c_n - c_{n-1}}{b_n - b_{n-1}}\Bigr| < \frac{\varepsilon}{2}
\qquad \text{for all } n > N,
\]
that is, since $b_{n-1} - b_n > 0$,
\[
|c_n - c_{n-1}| \le \frac{\varepsilon}{2}\,(b_{n-1} - b_n) \qquad (n > N).
\]
For $m > n > N$ the triangle inequality and a telescoping sum give
\[
|c_m - c_n| \le \sum_{k=n+1}^{m}|c_k - c_{k-1}|
\le \frac{\varepsilon}{2}\sum_{k=n+1}^{m}(b_{k-1} - b_k)
= \frac{\varepsilon}{2}\,(b_n - b_m) < \frac{\varepsilon}{2}\,b_n .
\]
Letting $m \to \infty$ and using $c_m \to 0$ yields $|c_n| \le \frac{\varepsilon}{2} b_n$, hence
\[
\Bigl|\frac{c_n}{b_n}\Bigr| \le \frac{\varepsilon}{2} < \varepsilon
\qquad (n > N),
\]
which proves the claim. Applying the claim to $c_n = u_n = a_n - A b_n$ gives
\[
\frac{a_n}{b_n} - A = \frac{a_n - A b_n}{b_n} = \frac{u_n}{b_n} \longrightarrow 0,
\qquad\text{so}\qquad \frac{a_n}{b_n} \longrightarrow A .
\]

\medskip
\noindent\textbf{Case 2: $A = +\infty$.}
Let $M > 0$ be given. Since $\dfrac{a_n - a_{n-1}}{b_n - b_{n-1}} \to +\infty$, there is $N$ such that
\[
\frac{a_n - a_{n-1}}{b_n - b_{n-1}} > M \qquad (n > N).
\]
As $b_n - b_{n-1} < 0$, multiplying by it reverses the inequality:
\[
a_n - a_{n-1} < M\,(b_n - b_{n-1})
\qquad\Longleftrightarrow\qquad
\bigl(a_n - M b_n\bigr) - \bigl(a_{n-1} - M b_{n-1}\bigr) < 0 .
\]
Thus $v_n := a_n - M b_n$ is strictly decreasing for $n \ge N$ (i.e.\ $v_{N} > v_{N+1} > \cdots$), and $v_n \to 0 - M \cdot 0 = 0$. A decreasing sequence is bounded below by its limit, so $v_n \ge 0$ for all $n \ge N$. Hence
\[
a_n \ge M b_n, \qquad\text{that is,}\qquad \frac{a_n}{b_n} \ge M \qquad (n \ge N).
\]
Since $M > 0$ is arbitrary, $\dfrac{a_n}{b_n} \to +\infty = A$.

\medskip
\noindent\textbf{Case 3: $A = -\infty$.}
Let $M > 0$ be given. Choose $N$ such that
\[
\frac{a_n - a_{n-1}}{b_n - b_{n-1}} < -M \qquad (n > N).
\]
Multiplying by $b_n - b_{n-1} < 0$ reverses the inequality:
\[
a_n - a_{n-1} > -M\,(b_n - b_{n-1}) = M\,(b_{n-1} - b_n),
\]
that is,
\[
\bigl(a_n + M b_n\bigr) - \bigl(a_{n-1} + M b_{n-1}\bigr) > 0 \qquad (n > N).
\]
So $w_n := a_n + M b_n$ is strictly increasing for $n \ge N$, and $w_n \to 0$. An increasing sequence is bounded above by its limit, so $w_n \le 0$ for all $n \ge N$. Hence
\[
a_n \le -M b_n, \qquad\text{that is,}\qquad \frac{a_n}{b_n} \le -M \qquad (n \ge N).
\]
Since $M > 0$ is arbitrary, $\dfrac{a_n}{b_n} \to -\infty = A$. This completes the proof.
\end{solution}

\begin{solution}{Solution 2.3.16}\label{xs:2.3.16}
Since $a \le a_n \le b$ for all $n \in \mathbb{N}_+$, the sequence $\{a_n\}$ is bounded. By the Bolzano--Weierstrass theorem, every bounded sequence of real numbers has a convergent subsequence; hence there is a subsequence $a_{n_k}$ and a number $\xi \in \mathbb{R}$ with
\[
a_{n_k} \longrightarrow \xi .
\]
Because $a \le a_{n_k} \le b$ for all $k$, passing to the limit gives $\xi \in [a,b]$.

Now suppose, for contradiction, that \emph{every} convergent subsequence of $\{a_n\}$ had the limit $\xi$. Then, in particular, the whole sequence could not converge to any number different from $\xi$; and since $\{a_n\}$ diverges by hypothesis, $\{a_n\}$ does not converge to $\xi$ either. We now exhibit a second subsequence with a limit different from $\xi$.

The statement ``$a_n \not\to \xi$'' means that there is an $\varepsilon > 0$ such that for every $N \in \mathbb{N}_+$ there exists an index $n > N$ with
\[
|a_n - \xi| \ge \varepsilon .
\]
Choose such indices recursively: pick $m_1 \in \mathbb{N}_+$ with $|a_{m_1} - \xi| \ge \varepsilon$; having chosen $m_k$, pick $m_{k+1} > m_k$ with $|a_{m_{k+1}} - \xi| \ge \varepsilon$. Then $(m_k)$ is strictly increasing, so $\{a_{m_k}\}$ is a subsequence of $\{a_n\}$, and
\[
|a_{m_k} - \xi| \ge \varepsilon \qquad \text{for all } k.
\]
This subsequence is again bounded (all its terms lie in $[a,b]$), so by the Bolzano--Weierstrass theorem it has a convergent subsequence $\{a_{m_{k_j}}\}_{j}$, say $a_{m_{k_j}} \to \eta$. The subsequence $\{a_{m_{k_j}}\}_j$ of $\{a_n\}$ is also a subsequence of $\{a_{m_k}\}$ and satisfies $|a_{m_{k_j}} - \xi| \ge \varepsilon$ for every $j$; letting $j \to \infty$ and using the continuity of the absolute value,
\[
|\eta - \xi| \ge \varepsilon > 0, \qquad\text{so}\qquad \eta \neq \xi .
\]
Thus $\{a_n\}$ has the subsequence $\{a_{n_k}\}$ converging to $\xi$ and the subsequence $\{a_{m_{k_j}}\}_j$ converging to $\eta \neq \xi$: two subsequences with different limits, as required.
\end{solution}

\begin{solution}{Solution 2.3.17}\label{xs:2.3.17}
We recall the Cauchy convergence criterion: $\{a_n\}$ converges if and only if for every $\varepsilon > 0$ there is $N \in \mathbb{N}_+$ such that $|a_m - a_n| < \varepsilon$ whenever $m > n \ge N$.

Let $\{a_n\}$ be monotone and bounded. We treat the nondecreasing case $a_1 \le a_2 \le a_3 \le \cdots$; the nonincreasing case follows by applying the argument to the nondecreasing bounded sequence $\{-a_n\}$.

So suppose $a_1 \le a_2 \le \cdots$ and the set $\{a_n : n \in \mathbb{N}_+\}$ is bounded above. By the least upper bound property of the real numbers, this nonempty bounded-above set has a supremum
\[
M = \sup\{a_n : n \in \mathbb{N}_+\} \in \mathbb{R},
\qquad\text{and}\qquad a_n \le M \ \ \text{for all } n .
\]
Let $\varepsilon > 0$. Since $M - \varepsilon$ is not an upper bound of the set, there is $N \in \mathbb{N}_+$ with
\[
a_N > M - \varepsilon .
\]
For all $m > n \ge N$, monotonicity gives $a_m \ge a_n$, so
\[
0 \le a_m - a_n \le M - a_n \le M - a_N < \varepsilon ,
\]
where we used $a_n \ge a_N$ (as $n \ge N$) and $a_m \le M$. Hence $|a_m - a_n| < \varepsilon$ for all $m > n \ge N$, i.e.\ $\{a_n\}$ is a Cauchy sequence. By the Cauchy convergence criterion, $\{a_n\}$ converges. (In fact the argument shows $\lim_{n\to\infty} a_n = M$.)
\end{solution}

\begin{solution}{Solution 2.3.18}\label{xs:2.3.18}
We use the \emph{nested interval theorem}; for reference we also recall the statement of the \emph{Bolzano--Weierstrass theorem}, which is what has to be proved.

\emph{Nested interval theorem.} If $I_k = [c_k, d_k]$ $(k \in \mathbb{N}_+)$ are closed intervals with
\[
I_1 \supseteq I_2 \supseteq I_3 \supseteq \cdots
\qquad\text{and}\qquad d_k - c_k \longrightarrow 0,
\]
then $\bigcap_{k=1}^{\infty} I_k$ consists of exactly one point $\xi$; moreover $c_k \to \xi$ and $d_k \to \xi$.

\emph{Bolzano--Weierstrass theorem.} Every bounded sequence of real numbers has a convergent subsequence.

\medskip
Let $\{a_n\}$ be an arbitrary bounded sequence: there is $M > 0$ with $|a_n| \le M$ for all $n \in \mathbb{N}_+$, so every term lies in $I_1 := [-M, M]$. We construct closed intervals
\[
I_1 \supseteq I_2 \supseteq I_3 \supseteq \cdots
\]
with the property
\[
(P_k) \qquad a_n \in I_k \ \text{for infinitely many } n .
\]
$I_1$ has property $(P_1)$ by the choice of $M$. Suppose $I_k = [c_k, d_k]$ has property $(P_k)$, and let $p_k = \tfrac{1}{2}(c_k + d_k)$ be its midpoint. The two closed halves $[c_k, p_k]$ and $[p_k, d_k]$ cover $I_k$, so
\[
\{n : a_n \in I_k\} = \{n : a_n \in [c_k,p_k]\} \cup \{n : a_n \in [p_k,d_k]\},
\]
and the set on the left is infinite; hence at least one of the two sets on the right is infinite. Choose that half as $I_{k+1}$; then $I_{k+1} \subseteq I_k$ and $(P_{k+1})$ holds. This completes the recursion.

Each $I_k$ is a nonempty closed interval and
\[
d_k - c_k = \frac{d_1 - c_1}{2^{\,k-1}} = \frac{2M}{2^{\,k-1}} \longrightarrow 0 .
\]
By the nested interval theorem there is exactly one point $\xi$ with $\xi \in \bigcap_k I_k$, and $c_k \to \xi$, $d_k \to \xi$.

Finally, because the set $\{n : a_n \in I_k\}$ is infinite for every $k$, we may choose indices recursively by
\[
n_1 < n_2 < n_3 < \cdots, \qquad a_{n_k} \in I_k \ \ (k \in \mathbb{N}_+):
\]
indeed, having chosen $n_k$, the set $\{n : a_n \in I_{k+1}\}$ is infinite, so it contains some $n > n_k$, which we take as $n_{k+1}$. Then $(n_k)$ is strictly increasing, so $\{a_{n_k}\}$ is a subsequence of $\{a_n\}$, and for every $k$,
\[
|a_{n_k} - \xi| \le d_k - c_k = \frac{2M}{2^{\,k-1}} \longrightarrow 0 .
\]
Hence $a_{n_k} \to \xi$: the bounded sequence $\{a_n\}$ has a convergent subsequence. Since $\{a_n\}$ was an arbitrary bounded sequence, the Bolzano--Weierstrass theorem is proved.
\end{solution}

\begin{solution}{Solution 2.3.19}\label{xs:2.3.19}
\emph{Cauchy convergence criterion.} A sequence $\{a_n\}$ of real numbers converges if and only if it is a Cauchy sequence, i.e.\ for every $\varepsilon > 0$ there is $N \in \mathbb{N}_+$ such that $|a_m - a_n| < \varepsilon$ whenever $m, n \ge N$.

We use the nested interval theorem: if $I_k = [c_k,d_k]$ are closed intervals with $I_1 \supseteq I_2 \supseteq \cdots$ and $d_k - c_k \to 0$, then $\bigcap_k I_k$ consists of exactly one point $\xi$, and $c_k \to \xi$, $d_k \to \xi$.

\medskip
\noindent\emph{Necessity.} Suppose $a_n \to \xi$. Given $\varepsilon > 0$, choose $N$ with $|a_n - \xi| < \varepsilon/2$ for all $n \ge N$. Then for $m, n \ge N$,
\[
|a_m - a_n| \le |a_m - \xi| + |\xi - a_n| < \frac{\varepsilon}{2} + \frac{\varepsilon}{2} = \varepsilon .
\]
So every convergent sequence is a Cauchy sequence.

\medskip
\noindent\emph{Sufficiency.} Let $\{a_n\}$ be a Cauchy sequence.

\emph{Step 1: $\{a_n\}$ is bounded.} Take $N_0$ such that $|a_m - a_n| < 1$ for all $m,n \ge N_0$. Then for every $m \ge N_0$, $|a_m| \le |a_{N_0}| + 1$, so with
\[
M = \max\{|a_1|, |a_2|, \dots, |a_{N_0-1}|, |a_{N_0}| + 1\}
\]
we have $|a_n| \le M$ for all $n \in \mathbb{N}_+$. In particular every term lies in $I_1 := [-M, M]$.

\emph{Step 2: construction of nested intervals.} We construct closed intervals $I_1 \supseteq I_2 \supseteq \cdots$ such that
\[
(Q_k) \qquad a_n \in I_k \ \text{for infinitely many } n .
\]
$I_1$ satisfies $(Q_1)$. Suppose $I_k = [c_k,d_k]$ satisfies $(Q_k)$ and let $p_k = \tfrac12(c_k+d_k)$. Since $[c_k,p_k] \cup [p_k,d_k] = I_k$, at least one of the two sets $\{n : a_n \in [c_k,p_k]\}$ and $\{n : a_n \in [p_k,d_k]\}$ is infinite; choose the corresponding half as $I_{k+1}$. Then $I_{k+1} \subseteq I_k$ and $(Q_{k+1})$ holds. Moreover
\[
d_k - c_k = \frac{2M}{2^{\,k-1}} \longrightarrow 0 .
\]
By the nested interval theorem, $\bigcap_k I_k = \{\xi\}$ for a single point $\xi$, and $c_k \to \xi$, $d_k \to \xi$.

\emph{Step 3: $a_n \to \xi$.} Let $\varepsilon > 0$. Choose $k$ with $d_k - c_k < \varepsilon/2$. Since $\{a_n\}$ is a Cauchy sequence, there is $N \in \mathbb{N}_+$ with
\[
|a_m - a_n| < \frac{\varepsilon}{2} \qquad (m, n \ge N).
\]
By $(Q_k)$ the set $\{n : a_n \in I_k\}$ is infinite, so we may choose an index $m \ge N$ with $a_m \in I_k$. Then $|a_m - \xi| \le d_k - c_k < \varepsilon/2$. Consequently, for every $n \ge N$,
\[
|a_n - \xi| \le |a_n - a_m| + |a_m - \xi| < \frac{\varepsilon}{2} + \frac{\varepsilon}{2} = \varepsilon .
\]
Hence $a_n \to \xi$: every Cauchy sequence of real numbers converges. Together with the necessity part, this is the Cauchy convergence criterion.
\end{solution}

\begin{solution}{Solution 2.3.20}\label{xs:2.3.20}
Set
\[
S_n = \sum_{k=1}^{n} x_k^2 \qquad (n \in \mathbb{N}_+), \qquad S_0 = 0 .
\]
Then $S_n - S_{n-1} = x_n^2 \ge 0$ for every $n \ge 1$, so $\{S_n\}$ is nondecreasing, and the hypothesis reads
\[
x_n S_n \longrightarrow 1 .
\]

\emph{Step 1: $S_n \to +\infty$.} If $\{S_n\}$ were bounded, then, being nondecreasing and bounded, it would converge to some $S \in \mathbb{R}$; then $x_n^2 = S_n - S_{n-1} \to 0$, hence $x_n \to 0$ and $x_n S_n \to 0 \cdot S = 0$, contradicting $x_n S_n \to 1$. Therefore $S_n \to +\infty$.

\emph{Step 2: $x_n \to 0$ and $S_{n-1}/S_n \to 1$.} Since $x_n S_n \to 1$, the sequence $\{x_n S_n\}$ is bounded, say $|x_n S_n| \le C$ for all $n$. As $S_n \to +\infty$,
\[
|x_n| = \frac{|x_n S_n|}{S_n} \le \frac{C}{S_n} \longrightarrow 0 ,
\]
so $x_n \to 0$, and then $x_n^2 = S_n - S_{n-1} \to 0$ as well. Also
\[
\Bigl|\frac{x_n}{S_n}\Bigr| = \frac{|x_n S_n|}{S_n^{2}} \le \frac{C}{S_n^{2}} \longrightarrow 0 ,
\qquad\text{so}\qquad
\Bigl|\frac{x_n^2}{S_n}\Bigr| = |x_n| \cdot \Bigl|\frac{x_n}{S_n}\Bigr| \longrightarrow 0 .
\]
Consequently
\[
\frac{S_{n-1}}{S_n} = 1 - \frac{x_n^2}{S_n} \longrightarrow 1 .
\]

\emph{Step 3: $S_n^3 - S_{n-1}^3 \to 3$.} For $n \ge 1$,
\[
S_n^3 - S_{n-1}^3 = \bigl(S_n - S_{n-1}\bigr)\bigl(S_n^2 + S_n S_{n-1} + S_{n-1}^2\bigr)
= (x_n S_n)^2\Bigl(1 + \frac{S_{n-1}}{S_n} + \Bigl(\frac{S_{n-1}}{S_n}\Bigr)^2\Bigr).
\]
By Steps 1--2, $(x_n S_n)^2 \to 1$ and $S_{n-1}/S_n \to 1$; hence
\[
S_n^3 - S_{n-1}^3 \longrightarrow 1\cdot(1 + 1 + 1) = 3 .
\]

\emph{Step 4: $S_n^3/n \to 3$.} Write $d_n = S_n^3 - S_{n-1}^3$, so that $d_n \to 3$. Given $\varepsilon > 0$, choose $N$ with $|d_n - 3| < \varepsilon$ for all $n > N$. For $n > N$, summation over $k = N+1, \dots, n$ gives
\[
S_n^3 - S_N^3 - 3(n-N) = \sum_{k=N+1}^{n}(d_k - 3),
\qquad\text{hence}\qquad
\bigl|S_n^3 - S_N^3 - 3(n-N)\bigr| \le \varepsilon\,(n-N).
\]
Dividing by $n > 0$,
\[
\Bigl|\frac{S_n^3}{n} - 3\Bigr| \le \varepsilon\,\frac{n-N}{n} + \frac{|S_N^3 - 3N|}{n} \le \varepsilon + \frac{|S_N^3 - 3N|}{n},
\]
and letting $n \to \infty$ yields $\limsup_{n\to\infty}\bigl|S_n^3/n - 3\bigr| \le \varepsilon$. Since $\varepsilon > 0$ is arbitrary, $S_n^3/n \to 3$.

\emph{Step 5: conclusion.} Since $S_n \to +\infty$, we have $S_n > 0$ for all sufficiently large $n$, and therefore
\[
\frac{S_n}{\sqrt[3]{3n}} = \Bigl(\frac{S_n^3}{3n}\Bigr)^{1/3} \longrightarrow \Bigl(\frac{3}{3}\Bigr)^{1/3} = 1
\]
by Step 4 and the continuity of the cube root. Finally, for such $n$,
\[
\sqrt[3]{3n}\; x_n = (x_n S_n)\cdot \frac{\sqrt[3]{3n}}{S_n}
= (x_n S_n)\cdot \Bigl(\frac{S_n}{\sqrt[3]{3n}}\Bigr)^{-1} \longrightarrow 1 \cdot 1^{-1} = 1 ,
\]
which is the required limit.
\end{solution}

\begin{solution}{Solution 2.3.21}\label{xs:2.3.21}
\begin{enumerate}[(1)]
  \item We prove by induction that $a_n \ge 3$ for all $n \ge 0$. Indeed, $a_0 = 3$, and if $a_{n-1} \ge 3$ then
  \[
  a_n = a_{n-1}^2 - 2 \ge 3^2 - 2 = 7 \ge 3 .
  \]
  Consequently, for every $n \ge 1$,
  \[
  a_n - a_{n-1} = a_{n-1}^2 - 2 - a_{n-1} = (a_{n-1} - 2)(a_{n-1} + 1) \ge (3-2)(3+1) = 4 .
  \]
  Summing these inequalities from $n = 1$ to $N$ gives, for all $N \in \mathbb{N}_+$,
  \[
  a_N \ge a_0 + 4N = 3 + 4N .
  \]
  Hence for every $G > 0$, every $N > (G-3)/4$ satisfies $a_N > G$; that is, $\lim_{n\to\infty} a_n = +\infty$.
  \item Put
  \[
  t = \frac{3+\sqrt{5}}{2} > 1 .
  \]
  Then $t^{-1} = \dfrac{2}{3+\sqrt5} = \dfrac{2(3-\sqrt5)}{(3+\sqrt5)(3-\sqrt5)} = \dfrac{3-\sqrt5}{2}$, so that
  \[
  t + t^{-1} = 3, \qquad t - t^{-1} = \sqrt{5} .
  \]
  \emph{Closed form of $a_n$.} We claim that
  \[
  a_n = t^{\,2^n} + t^{-2^n} \qquad (n \ge 0).
  \]
  For $n = 0$ this is $t + t^{-1} = 3 = a_0$. If $a_{n-1} = t^{\,2^{n-1}} + t^{-2^{n-1}}$, then, using $2\cdot 2^{n-1} = 2^n$,
  \[
  a_n = a_{n-1}^2 - 2 = \bigl(t^{\,2^{n-1}} + t^{-2^{n-1}}\bigr)^2 - 2 = t^{\,2^n} + 2 + t^{-2^n} - 2 = t^{\,2^n} + t^{-2^n},
  \]
  which proves the claim by induction. In particular $a_n \ge 3 > 0$ for all $n$, so the quotient in (2) is well defined.

  \emph{A telescoping product.} For every $k \ge 0$, applying $(u - u^{-1})(u + u^{-1}) = u^2 - u^{-2}$ with $u = t^{\,2^k} \neq 1$ (note $t > 1$) gives
  \[
  t^{\,2^k} + t^{-2^k} = \frac{t^{\,2^{k+1}} - t^{-2^{k+1}}}{t^{\,2^k} - t^{-2^k}} .
  \]
  Multiplying these identities for $k = 0, 1, \dots, n-1$ (the denominators cancel successively) yields the telescoping formula
  \[
  a_0 a_1 \cdots a_{n-1} = \prod_{k=0}^{n-1}\bigl(t^{\,2^k} + t^{-2^k}\bigr)
  = \frac{t^{\,2^n} - t^{-2^n}}{t - t^{-1}}
  = \frac{t^{\,2^n} - t^{-2^n}}{\sqrt5} .
  \]
  Therefore, using the closed form of $a_n$,
  \[
  \frac{a_n}{a_0 a_1 \cdots a_{n-1}}
  = \sqrt5\;\frac{t^{\,2^n} + t^{-2^n}}{t^{\,2^n} - t^{-2^n}}
  = \sqrt5\;\frac{1 + t^{-2^{n+1}}}{1 - t^{-2^{n+1}}} ,
  \]
  where the last expression is obtained by dividing numerator and denominator by $t^{\,2^n}$ and using $t^{-2^n}/t^{\,2^n} = t^{-2^{n+1}}$.

  \emph{Limit.} Since $0 < t^{-1} < 1$, we have $t^{-2^{n+1}} = (t^{-1})^{2^{n+1}} \to 0$ as $n \to \infty$. Hence
  \[
  \frac{1 + t^{-2^{n+1}}}{1 - t^{-2^{n+1}}} \longrightarrow \frac{1+0}{1-0} = 1 ,
  \]
  and consequently
  \[
  \lim_{n\to\infty} \frac{a_n}{a_0 a_1 \cdots a_{n-1}} = \sqrt5 .
  \]
\end{enumerate}
\end{solution}


\srcline{Exercise set 2.4, problems 1--10}

\begin{solution}{Solution 2.4.1}\label{xs:2.4.1}
Recall the definitions. Let $A$ be a nonempty set of real numbers. A number $M$ is called the \emph{supremum} (least upper bound) of $A$, written $M = \sup A$, if
\begin{enumerate}[(i)]
  \item $x \le M$ for every $x \in A$ (that is, $M$ is an upper bound of $A$), and
  \item $M \le M'$ for every upper bound $M'$ of $A$ (that is, $M$ is the least such bound).
\end{enumerate}
Dually, $m$ is the \emph{infimum} (greatest lower bound) of $A$ if $m \le x$ for every $x \in A$ and $m' \le m$ for every lower bound $m'$ of $A$.

\medskip
\noindent\textbf{Uniqueness of the supremum.} Suppose $M_1$ and $M_2$ both satisfy (i) and (ii). Since $M_1$ is an upper bound of $A$ and $M_2$ is a least upper bound, (ii) applied to $M_2$ with the upper bound $M' = M_1$ gives $M_2 \le M_1$. Interchanging the roles of $M_1$ and $M_2$, the same argument gives $M_1 \le M_2$. Hence $M_1 \le M_2$ and $M_2 \le M_1$, and by the trichotomy of the order of $\mathbb{R}$ we get $M_1 = M_2$. Therefore the supremum, if it exists, is unique.

\medskip
\noindent\textbf{Uniqueness of the infimum.} Suppose $m_1$ and $m_2$ are both infima of $A$. Since $m_2$ is a lower bound of $A$ and $m_1$ is a greatest lower bound, the defining property of $m_1$ gives $m_2 \le m_1$. Symmetrically $m_1 \le m_2$. Hence $m_1 = m_2$, so the infimum, if it exists, is also unique.

\medskip
For a set that is unbounded above the book adopts the convention $\sup A = +\infty$, and for a set unbounded below the convention $\inf A = -\infty$; these extended values are unique as well, by the same comparison argument. This proves the uniqueness of the supremum and of the infimum in all cases.
\end{solution}

\begin{solution}{Solution 2.4.2}\label{xs:2.4.2}
Throughout, $A \ne \varnothing$; we also use that an element $M \in A$ with $x \le M$ for all $x \in A$ is called the maximum of $A$ and is written $M = \max A$, and dually for the minimum.

\medskip
\noindent\textbf{(1)} 

\emph{Sufficiency.} Suppose $A$ has a maximum $\max A$. Then $\max A \in A$, and by the definition of a maximum, $x \le \max A$ for every $x \in A$; thus $\max A$ is an upper bound of $A$. Let $M'$ be any upper bound of $A$. Since $\max A \in A$, the defining property of an upper bound gives $\max A \le M'$. Hence $\max A$ is the least upper bound of $A$, that is, $\sup A = \max A$, and in particular $\sup A \in A$.

\emph{Necessity.} Suppose $\sup A \in A$. By definition of the supremum, $x \le \sup A$ for every $x \in A$, and $\sup A$ itself belongs to $A$. Therefore $\sup A$ is an element of $A$ that dominates every element of $A$, i.e.\ $A$ has a maximum and $\max A = \sup A$. In particular $\sup A = \max A$, as claimed.

\medskip
\noindent\textbf{(2)} 

\emph{Sufficiency.} Suppose $A$ has a minimum $\min A$. Then $\min A \in A$ and $\min A \le x$ for every $x \in A$, so $\min A$ is a lower bound of $A$. If $m'$ is any lower bound of $A$, then, since $\min A \in A$, we get $m' \le \min A$. Hence $\min A$ is the greatest lower bound of $A$, that is, $\inf A = \min A$, and $\inf A \in A$.

\emph{Necessity.} Suppose $\inf A \in A$. By definition of the infimum, $\inf A \le x$ for every $x \in A$, and $\inf A \in A$. Hence $\inf A$ is an element of $A$ that is dominated by no element of $A$, i.e.\ $A$ has a minimum and $\min A = \inf A$; in particular $\inf A = \min A$.
\end{solution}

\begin{solution}{Solution 2.4.3}\label{xs:2.4.3}
The correct conclusion is
\[
\sup S \le a .
\]
Indeed, assume that $S$ is nonempty and bounded above, so that $\sup S$ is a real number. Since $x < a$ for every $x \in S$, we have $x \le a$ for every $x \in S$, and therefore $a$ is an upper bound of $S$. By the defining property of the supremum as the \emph{least} upper bound, $\sup S \le a$.

The strict inequality $\sup S < a$ is \emph{not} a valid conclusion. For example, take $S = (0,1)$ and $a = 1$. Then $x < 1$ for every $x \in S$, while $\sup S = 1 = a$, so $\sup S < a$ fails. (Here $S$ is a nonempty set bounded above, and its supremum exists by the supremum principle.) Thus, from the hypothesis $x < a$ for all $x \in S$ one may conclude only $\sup S \le a$; the inequality can degenerate into an equality precisely when $a$ is the least upper bound of $S$ but does not belong to $S$.
\end{solution}

\begin{solution}{Solution 2.4.4}\label{xs:2.4.4}
We use the fact that the supremum is the least upper bound and the infimum the greatest lower bound, together with the density of $\mathbb{Q}$ in $\mathbb{R}$ (between any two distinct real numbers there lies a rational number) and the convention $\sup A = +\infty$ for a set unbounded above, $\inf A = -\infty$ for a set unbounded below.

\medskip
\noindent\textbf{(1)} $S = \{x \in \mathbb{Q} \mid x > \sqrt{2}\}$.

\emph{Infimum.} Since $\sqrt{2} < x$ for every $x \in S$, the number $\sqrt{2}$ is a lower bound of $S$. Let $b$ be any lower bound of $S$ and suppose, for contradiction, that $b > \sqrt{2}$. By the density of $\mathbb{Q}$ in $\mathbb{R}$ there is a rational number $q$ with $\sqrt{2} < q < b$. Then $q \in \mathbb{Q}$ and $q > \sqrt{2}$, so $q \in S$; but $q < b$, contradicting the fact that $b$ is a lower bound of $S$. Hence $b \le \sqrt{2}$ for every lower bound $b$, so $\sqrt{2}$ is the greatest lower bound:
\[
\inf S = \sqrt{2}.
\]
Note that $\sqrt{2} \notin \mathbb{Q}$, hence $\inf S \notin S$, so $S$ has no minimum.

\emph{Supremum.} $S$ is unbounded above: given any real $M$, choose (by the Archimedean property) a natural number $n$ with $n > \max\{M, 2\}$; then $n \in \mathbb{Q}$ and $n > \sqrt{2}$, that is, $n \in S$ and $n > M$. Hence $S$ has no upper bound and
\[
\sup S = +\infty .
\]

\medskip
\noindent\textbf{(2)} $S = \{y \in \mathbb{R} \mid y = x^2,\ x \in (-\sqrt{3}, 10)\}$.

Since $x^2 \ge 0$ for all real $x$, we have $y \ge 0$ for all $y \in S$, and $0 \in S$ because $x = 0 \in (-\sqrt{3}, 10)$ gives $y = 0$. Hence $0$ is a lower bound of $S$ that belongs to $S$, so
\[
\inf S = 0 = \min S .
\]
For the upper end, if $x \in (-\sqrt{3}, 10)$ then $|x| < 10$, hence $y = x^2 < 100$; so $100$ is an upper bound of $S$. Conversely, let $c < 100$. Choose the real number $x = \sqrt{t}$ where $t$ is any number with $\max\{0, c\} < t < 100$ (for instance $t = \frac{\max\{0,c\} + 100}{2}$ if $c < 100$). Then $0 \le t < 100$ gives $x = \sqrt{t} \in [0, 10) \subset (-\sqrt{3}, 10)$, so $y = t \in S$; and, if $c \ge 0$, then $t > c$; if $c < 0$, then $t \ge 0 > c$. In both cases some element of $S$ exceeds $c$, so $c$ is not an upper bound of $S$. Hence $100$ is the least upper bound:
\[
\sup S = 100 .
\]
Since $x^2 = 100$ would force $x = \pm 10$, and neither of these lies in $(-\sqrt{3}, 10)$, we have $100 \notin S$ and $S$ has no maximum. Thus $S = [0, 100)$.

\medskip
\noindent\textbf{(3)} $S = \{y \in \mathbb{R} \mid y = \arctan x,\ x \in \mathbb{Q}\}$.

The function $\arctan$ is strictly increasing on $\mathbb{R}$ with $-\frac{\pi}{2} < \arctan x < \frac{\pi}{2}$ for every $x \in \mathbb{R}$. Hence $-\frac{\pi}{2}$ is a lower bound and $\frac{\pi}{2}$ is an upper bound of $S$.

Let $b < \frac{\pi}{2}$. Choose a real number $t$ with $\max\left\{b, -\frac{\pi}{2}\right\} < t < \frac{\pi}{2}$; then $\tan t$ is a real number, and by the density of $\mathbb{Q}$ in $\mathbb{R}$ there is a rational number $q$ with $q > \tan t$. Since $\arctan$ is strictly increasing and $\arctan(\tan t) = t$ for $t \in \left(-\frac{\pi}{2}, \frac{\pi}{2}\right)$,
\[
\arctan q > \arctan(\tan t) = t > b .
\]
Thus some element $\arctan q$ of $S$ exceeds $b$, so no $b < \frac{\pi}{2}$ is an upper bound of $S$, and therefore
\[
\sup S = \frac{\pi}{2} .
\]
Similarly, let $c > -\frac{\pi}{2}$ and choose a real number $t$ with $-\frac{\pi}{2} < t < \min\left\{c, \frac{\pi}{2}\right\}$; by density of $\mathbb{Q}$ choose $q \in \mathbb{Q}$ with $q < \tan t$. Then $q \in \mathbb{Q}$ and
\[
\arctan q < \arctan(\tan t) = t < c ,
\]
so some element of $S$ is smaller than $c$ and $c$ is not a lower bound of $S$. Hence
\[
\inf S = -\frac{\pi}{2} .
\]
Since $-\frac{\pi}{2} < \arctan q < \frac{\pi}{2}$ for every $q \in \mathbb{Q}$, neither endpoint belongs to $S$ (all elements of $S$ equal $\arctan q$ with $q \in \mathbb{Q}$); in particular $S$ has neither a maximum nor a minimum. So $S$ is a dense subset of the open interval $\left(-\frac{\pi}{2}, \frac{\pi}{2}\right)$ whose supremum and infimum are its endpoints.

\medskip
\noindent\textbf{(4)} $S = \left\{1, \dfrac{1}{2}, \dfrac{2}{3}, \dfrac{3}{4}, \cdots \right\} = \{1\} \cup \left\{\dfrac{n}{n+1} : n \in \mathbb{N}_+\right\}$.

For every $n \in \mathbb{N}_+$ we have $\dfrac{n}{n+1} < 1$, and $1 \in S$; hence $1$ is an element of $S$ that is an upper bound of $S$, so
\[
\sup S = 1 = \max S .
\]
For the lower end, $\dfrac{n}{n+1} \ge \dfrac{1}{2}$ for every $n \in \mathbb{N}_+$ (equivalently $2n \ge n+1$, which holds for every $n \ge 1$), and $\dfrac{1}{2} \in S$; hence $\dfrac{1}{2}$ is a lower bound of $S$ belonging to $S$, so
\[
\inf S = \frac{1}{2} = \min S .
\]
Indeed the sequence $\dfrac{n}{n+1}$ is strictly increasing (since $\dfrac{n+1}{n+2} - \dfrac{n}{n+1} = \dfrac{1}{(n+1)(n+2)} > 0$), so its smallest term is $\dfrac{1}{2}$ and its terms approach $1$ from below, which is consistent with the values found above.
\end{solution}

\begin{solution}{Solution 2.4.5}\label{xs:2.4.5}
Let $A \ne \varnothing$ and $-A = \{x \mid -x \in A\}$. The map $x \mapsto -x$ is a bijection of $A$ onto $-A$ that reverses the order: for real numbers $u, v$ we have $u \le v \iff -u \ge -v$. We prove both assertions from the definitions.

\medskip
\noindent\textbf{(1) $\sup(-A) = -\inf A$.}

Assume first that $A$ is bounded below, so that $m = \inf A \in \mathbb{R}$. For every $x \in -A$ we have $-x \in A$, hence $-x \ge m$, that is, $x \le -m$. Thus $-m$ is an upper bound of $-A$. Now let $c$ be any upper bound of $-A$; for every $a \in A$ we have $-a \in -A$, hence $-a \le c$, that is, $a \ge -c$. Therefore $-c$ is a lower bound of $A$, and, since $m = \inf A$ is the greatest lower bound, $m \ge -c$, i.e.\ $c \ge -m$. So $-m$ is the least upper bound of $-A$, giving $\sup(-A) = -m = -\inf A$.

If $A$ is unbounded below, then $\inf A = -\infty$ by convention, and for every real $M$ there exists $a \in A$ with $a < -M$; then $-a \in -A$ and $-a > M$, so $-A$ is unbounded above and $\sup(-A) = +\infty = -(-\infty) = -\inf A$. The identity therefore also holds in the extended sense.

\medskip
\noindent\textbf{(2) $\inf(-A) = -\sup A$.}

Assume first that $A$ is bounded above, so that $M_0 = \sup A \in \mathbb{R}$. For every $x \in -A$ we have $-x \in A$, hence $-x \le M_0$, that is, $x \ge -M_0$; so $-M_0$ is a lower bound of $-A$. Now let $c$ be any lower bound of $-A$; for every $a \in A$, $-a \in -A$, hence $-a \ge c$, that is, $a \le -c$. Therefore $-c$ is an upper bound of $A$, and, since $M_0 = \sup A$ is the least upper bound, $M_0 \le -c$, i.e.\ $c \le -M_0$. Thus $-M_0$ is the greatest lower bound of $-A$, and $\inf(-A) = -M_0 = -\sup A$.

If $A$ is unbounded above, then $\sup A = +\infty$ and for every real $M$ there is $a \in A$ with $a > -M$, whence $-a \in -A$ and $-a < M$; so $-A$ is unbounded below and $\inf(-A) = -\infty = -(+\infty) = -\sup A$. Again the identity holds in the extended sense.
\end{solution}

\begin{solution}{Solution 2.4.6}\label{xs:2.4.6}
Let $A, B$ be nonempty sets of real numbers bounded above, and put
\[
A + B = \{x + y \mid x \in A,\ y \in B\}.
\]
\emph{Step 1: $\sup A + \sup B$ is an upper bound of $A + B$.} For any $x \in A$ and $y \in B$ we have $x \le \sup A$ and $y \le \sup B$ by the definition of the supremum as an upper bound; adding these inequalities gives
\[
x + y \le \sup A + \sup B .
\]
Hence every element of $A + B$ is at most $\sup A + \sup B$, so $\sup A + \sup B$ is an upper bound of $A + B$ and therefore
\[
\sup(A + B) \le \sup A + \sup B .
\]

\emph{Step 2: the reverse inequality is also true (so in fact equality holds).} Let $\varepsilon > 0$. By the defining property of the supremum as the \emph{least} upper bound, $\sup A - \frac{\varepsilon}{2}$ is not an upper bound of $A$, so there is $x_0 \in A$ with $x_0 > \sup A - \frac{\varepsilon}{2}$; likewise there is $y_0 \in B$ with $y_0 > \sup B - \frac{\varepsilon}{2}$. Then $x_0 + y_0 \in A + B$ and
\[
x_0 + y_0 > \sup A + \sup B - \varepsilon .
\]
Since every element of $A + B$ is at most $\sup(A+B)$, this shows that $\sup(A + B) > \sup A + \sup B - \varepsilon$ for every $\varepsilon > 0$, whence $\sup(A+B) \ge \sup A + \sup B$.

Combining the two steps, $\sup(A + B) = \sup A + \sup B$.

\emph{Conclusion for $S$.} Since $S \subset \{x + y \mid x \in A,\ y \in B\} = A + B$, every element of $S$ is an element of $A + B$ and hence is at most $\sup A + \sup B$; thus $\sup A + \sup B$ is an upper bound of $S$, and consequently
\[
\sup S \le \sup A + \sup B .
\]

\emph{The strict inequality need not hold.} Take $A = B = [0,1]$, so that $\sup A = \sup B = 1$ and $\sup A + \sup B = 2$, and take
\[
S = \{x + y \mid x \in A,\ y = 1\} = [1, 2] \subset [0, 2] = \{x + y \mid x \in A,\ y \in B\}.
\]
Then $\sup S = 2 = \sup A + \sup B$, so the inequality $\sup S \le \sup A + \sup B$ is an equality here; the strict inequality $\sup S < \sup A + \sup B$ is therefore not a valid conclusion.
\end{solution}

\begin{solution}{Solution 2.4.7}\label{xs:2.4.7}
Let $A, B$ be nonempty sets of real numbers bounded above, and suppose that the set $S$ satisfies
\[
S \supset \{x + y \mid x \in A,\ y \in B\} = A + B .
\]
(We also assume $S$ is bounded above, so that $\sup S$ is a real number; otherwise $\sup S = +\infty$ by convention and the inequality $\sup S \ge \sup A + \sup B$ is trivially true.)

\emph{Step 1: every $x + y$ with $x \in A$, $y \in B$ is at most $\sup S$.} Indeed, $x + y \in A + B \subset S$, and $\sup S$ is an upper bound of $S$. Hence
\[
x + y \le \sup S \qquad \text{for all } x \in A,\ y \in B. \tag{1}
\]

\emph{Step 2: $\sup S - \sup B$ is an upper bound of $A$.} Fix $x \in A$. By (1), $x + y \le \sup S$ for every $y \in B$, that is, $y \le \sup S - x$ for every $y \in B$. Hence $\sup S - x$ is an upper bound of $B$, and, since $\sup B$ is the least upper bound of $B$,
\[
\sup B \le \sup S - x, \qquad \text{i.e.} \qquad x \le \sup S - \sup B .
\]
As $x \in A$ was arbitrary, $\sup S - \sup B$ is an upper bound of $A$.

\emph{Step 3: conclusion.} Since $\sup A$ is the least upper bound of $A$, Step 2 gives $\sup A \le \sup S - \sup B$, that is,
\[
\sup S \ge \sup A + \sup B .
\]

\emph{The strict inequality need not hold.} Take $A = B = [0,1]$ and $S = [0,2]$. Then $\sup A = \sup B = 1$, so $\sup A + \sup B = 2$; and
\[
\{x + y \mid x \in A,\ y \in B\} = [0,2] \subset S,
\]
while $\sup S = 2 = \sup A + \sup B$. Thus the inequality $\sup S \ge \sup A + \sup B$ becomes an equality in this example, and the strict inequality $\sup S > \sup A + \sup B$ is not a valid conclusion in general. (For comparison, with the same $A$ and $B$ but $S = [0,3]$ one does get the strict inequality $3 > 2$.)
\end{solution}

\begin{solution}{Solution 2.4.8}\label{xs:2.4.8}
Let $A, B \ne \varnothing$ be sets of real numbers and let $\alpha \ge 0$. For a set $C$ of real numbers and a real number $\lambda$ write $\lambda C = \{\lambda c \mid c \in C\}$, so that
\[
A + \alpha B = \{a + \alpha b \mid a \in A,\ b \in B\} = A + (\alpha B).
\]
We first record two lemmas.

\medskip
\noindent\textbf{Lemma 1 (supremum and infimum of a sum).} If $A, B \ne \varnothing$, then $\sup(A+B) = \sup A + \sup B$ and $\inf(A+B) = \inf A + \inf B$ (in the extended sense).

\emph{Proof of the sup part.} For $a \in A$, $b \in B$ we have $a \le \sup A$ and $b \le \sup B$, hence $a + b \le \sup A + \sup B$; so $\sup A + \sup B$ is an upper bound of $A + B$ and $\sup(A+B) \le \sup A + \sup B$. If, say, $\sup A = +\infty$, then $A + B$ is unbounded above: fix $b_0 \in B$; for every real $R$ there is $a \in A$ with $a > R - b_0$, whence $a + b_0 \in A + B$ and $a + b_0 > R$. Hence $\sup(A+B) = +\infty = \sup A + \sup B$ in that case, and the case $\sup B = +\infty$ is identical. It therefore remains to treat the case in which both suprema are finite. Let $\varepsilon > 0$. Since $\sup A - \frac{\varepsilon}{2}$ is not an upper bound of $A$, there is $a_0 \in A$ with $a_0 > \sup A - \frac{\varepsilon}{2}$, and similarly there is $b_0 \in B$ with $b_0 > \sup B - \frac{\varepsilon}{2}$. Then $a_0 + b_0 \in A + B$ and $a_0 + b_0 > \sup A + \sup B - \varepsilon$, while $a_0 + b_0 \le \sup(A+B)$; hence $\sup(A+B) > \sup A + \sup B - \varepsilon$ for every $\varepsilon > 0$, which forces $\sup(A+B) \ge \sup A + \sup B$.

\emph{Proof of the inf part.} For $a \in A$, $b \in B$ we have $a \ge \inf A$ and $b \ge \inf B$, hence $a + b \ge \inf A + \inf B$; so $\inf A + \inf B$ is a lower bound of $A+B$ and $\inf(A+B) \ge \inf A + \inf B$. If, say, $\inf A = -\infty$, then $A+B$ is unbounded below: fix $b_0 \in B$; for every real $R$ there is $a \in A$ with $a < R - b_0$, whence $a + b_0 \in A + B$ and $a + b_0 < R$. Hence $\inf(A+B) = -\infty = \inf A + \inf B$ in that case, and the case $\inf B = -\infty$ is identical. It therefore remains to treat the case in which both infima are finite. Given $\varepsilon > 0$, the number $\inf A + \frac{\varepsilon}{2}$ is not a lower bound of $A$, so there is $a_0 \in A$ with $a_0 < \inf A + \frac{\varepsilon}{2}$, and likewise there is $b_0 \in B$ with $b_0 < \inf B + \frac{\varepsilon}{2}$. Then $a_0 + b_0 \in A + B$ and $a_0 + b_0 < \inf A + \inf B + \varepsilon$, while $a_0 + b_0 \ge \inf(A+B)$; hence $\inf(A+B) < \inf A + \inf B + \varepsilon$ for every $\varepsilon > 0$, giving $\inf(A+B) \le \inf A + \inf B$. $\square$

\medskip
\noindent\textbf{Lemma 2 (multiplication by $\alpha \ge 0$).} If $B \ne \varnothing$ and $\alpha \ge 0$, then $\sup(\alpha B) = \alpha \sup B$ and $\inf(\alpha B) = \alpha \inf B$.

\emph{Proof.} If $\alpha > 0$: for $b \in B$ we have $b \le \sup B$, hence $\alpha b \le \alpha \sup B$, so $\alpha \sup B$ is an upper bound of $\alpha B$. (If $\sup B = +\infty$ then $\alpha B$ is unbounded above, so $\sup(\alpha B) = +\infty = \alpha \sup B$; if $\inf B = -\infty$ then $\alpha B$ is unbounded below, so $\inf(\alpha B) = -\infty = \alpha \inf B$. Hence we may suppose both are finite.) Given $\varepsilon > 0$, the number $\sup B - \frac{\varepsilon}{\alpha}$ is not an upper bound of $B$, so there is $b_0 \in B$ with $b_0 > \sup B - \frac{\varepsilon}{\alpha}$, whence $\alpha b_0 \in \alpha B$ and $\alpha b_0 > \alpha \sup B - \varepsilon$. Thus $\sup(\alpha B) = \alpha \sup B$. The analogous argument with lower bounds (for $b \in B$, $b \ge \inf B$ gives $\alpha b \ge \alpha \inf B$; and some $b_0 \in B$ satisfies $b_0 < \inf B + \frac{\varepsilon}{\alpha}$, whence $\alpha b_0 < \alpha \inf B + \varepsilon$) gives $\inf(\alpha B) = \alpha \inf B$. If $\alpha = 0$, then $\alpha B = \{0\}$, so $\sup(\alpha B) = 0 = 0 \cdot \sup B$ and $\inf(\alpha B) = 0 = 0 \cdot \inf B$ (using the usual convention $0 \cdot (\pm\infty) = 0$ in the unbounded case). $\square$

\medskip
\noindent\textbf{(1)} Since $A + \alpha B = A + (\alpha B)$, Lemma 1 (sup part) and Lemma 2 give
\[
\sup(A + \alpha B) = \sup\bigl(A + (\alpha B)\bigr) = \sup A + \sup(\alpha B) = \sup A + \alpha \sup B .
\]

\medskip
\noindent\textbf{(2)} Likewise, by Lemma 1 (inf part) and Lemma 2,
\[
\inf(A + \alpha B) = \inf\bigl(A + (\alpha B)\bigr) = \inf A + \inf(\alpha B) = \inf A + \alpha \inf B .
\]
This completes the proof of both identities. (The hypotheses $\alpha \ge 0$ and $A, B \ne \varnothing$ are exactly what the two lemmas require: without $\alpha \ge 0$ the set $\alpha B$ would be reflected and the formulas for the supremum and infimum would be interchanged.)
\end{solution}

\begin{solution}{Solution 2.4.9}\label{xs:2.4.9}
Let $A, B$ be sets of real numbers with $A \subset B$, and assume both are nonempty. (If $A = \varnothing$, then $\sup A = -\infty$ and $\inf A = +\infty$ by the usual convention, and both inequalities are trivially true.)

\medskip
\noindent\textbf{$\sup A \le \sup B$.} Assume $B$ is bounded above, so that $\sup B \in \mathbb{R}$; otherwise $\sup B = +\infty$ and the inequality is trivial. Let $x \in A$ be arbitrary. Since $A \subset B$, we have $x \in B$, and $\sup B$ is an upper bound of $B$, so $x \le \sup B$. As $x \in A$ was arbitrary, $\sup B$ is an upper bound of $A$. By the defining property of the supremum of $A$ as the \emph{least} upper bound of $A$,
\[
\sup A \le \sup B .
\]

\medskip
\noindent\textbf{$\inf A \ge \inf B$.} Assume $B$ is bounded below, so that $\inf B \in \mathbb{R}$; otherwise $\inf B = -\infty$ and the inequality is trivial. Let $x \in A$ be arbitrary. Since $A \subset B$, we have $x \in B$, and $\inf B$ is a lower bound of $B$, so $x \ge \inf B$. As $x \in A$ was arbitrary, $\inf B$ is a lower bound of $A$. By the defining property of the infimum of $A$ as the \emph{greatest} lower bound of $A$,
\[
\inf A \ge \inf B .
\]
Both statements are therefore proved, and they express the obvious fact that enlarging a set can only raise its supremum and lower its infimum.
\end{solution}

\begin{solution}{Solution 2.4.10}\label{xs:2.4.10}
Let $A, B$ be sets of real numbers such that
\[
x \le y \qquad \text{for every } x \in A \text{ and every } y \in B. \tag{1}
\]
Assume $A \ne \varnothing$ and $B \ne \varnothing$ (if $A = \varnothing$, then $\sup A = -\infty$ and the inequality $\sup A \le \inf B$ is trivial; if $B = \varnothing$, then $\inf B = +\infty$ and it is again trivial).

\emph{Step 1: $B$ is bounded below and $A$ is bounded above.} Fix $x_0 \in A$. By (1), $x_0 \le y$ for every $y \in B$; thus $x_0$ is a lower bound of $B$, so $B$ is nonempty and bounded below and, by the supremum principle (applied to the set of lower bounds of $B$), its greatest lower bound $\inf B \in \mathbb{R}$ exists. Similarly, fix $y_0 \in B$; by (1), $x \le y_0$ for every $x \in A$, so $y_0$ is an upper bound of $A$ and $\sup A \in \mathbb{R}$ exists.

\emph{Step 2: $\inf B$ is an upper bound of $A$.} Let $x \in A$. By (1), $x \le y$ for every $y \in B$, so $x$ is a lower bound of $B$. Since $\inf B$ is the greatest lower bound of $B$, we get
\[
x \le \inf B .
\]
As $x \in A$ was arbitrary, $\inf B$ is an upper bound of the set $A$.

\emph{Step 3: conclusion.} Since $\sup A$ is the least upper bound of $A$ and $\inf B$ is an upper bound of $A$ (Step 2),
\[
\sup A \le \inf B ,
\]
which is the desired inequality.
\end{solution}

\end{document}
